% Le télescope de Newton — Physique 1ère TI — Cours signé OnBuch+
% Programme officiel MINESEC. Compiler avec : tectonic N11-telescope-newton.tex
\newcommand{\DOCMATIERE}{Physique}
\newcommand{\DOCNIVEAU}{1ère TI}
\newcommand{\DOCMODULE}{Module 3 : Optique}
\newcommand{\DOCLECON}{N11}
\newcommand{\DOCTITRE}{Le télescope de Newton}
% OnBuch+ — préambule « cours signés » ENRICHI (compilé avec tectonic / XeLaTeX)
% Système visuel « Pop » d'OnBuch+ : fond crème (jamais blanc), bordures encre
% épaisses, ombres « dures » décalées, titres Archivo Black en capitales.
% Version MATHÉMATIQUES : graphes (pgfplots/tikz), boîtes pédagogiques
% (définition, propriété, méthode, attention, le savais-tu, exercice résolu,
% exercice, TP, bilan), page de garde et signature OnBuch+ ancrée en bas.
%
% Macros attendues dans main.tex AVANT \input{preamble} :
%   \DOCMATIERE  \DOCNIVEAU  \DOCMODULE  \DOCLECON (numéro)  \DOCTITRE
\documentclass[11pt]{article}

\usepackage{fontspec}
\usepackage[french]{babel}
\usepackage[a4paper,margin=2cm,top=3.4cm,bottom=2.8cm,headheight=1.4cm,footskip=1.3cm]{geometry}
\usepackage{amsmath,amssymb,amsfonts}
\usepackage{mathtools} % \xmapsto, \coloneqq... souvent utilisés par les modèles
\usepackage{cancel} % simplifications barrées (bilans de forces)
% \abs{z} (module, valeur absolue) et \norm{u} (norme), utilisés par les modèles
\providecommand{\abs}{}\let\abs\relax\DeclarePairedDelimiter{\abs}{\lvert}{\rvert}
\providecommand{\norm}{}\let\norm\relax\DeclarePairedDelimiter{\norm}{\lVert}{\rVert}
\usepackage[table]{xcolor} % [table] : fournit aussi \rowcolors (alternance de lignes)
\usepackage{pagecolor}
\usepackage{tikz}
\usetikzlibrary{arrows.meta,positioning,calc,shapes.geometric,shapes.misc,shapes.arrows,decorations.pathmorphing,decorations.pathreplacing,decorations.markings,patterns,shadows,mindmap,trees,fit,backgrounds,matrix,intersections,angles,quotes}
\usepackage{pgfplots}
\usepackage[version=4]{mhchem} % équations nucléaires \ce{^{235}_{92}U}
\usepackage[european,siunitx]{circuitikz} % schémas électriques
\pgfplotsset{compat=1.18}
\usepgfplotslibrary{fillbetween,groupplots} % aires entre courbes (calcul intégral)
\usepackage{enumitem}
\usepackage{fancyhdr}
% [most] charge toutes les bibliothèques tcolorbox (listings, minted, raster,
% documentation...) et gonfle inutilement la mémoire TeX sur un document long
% (~300 boîtes) : on ne charge que ce qui est réellement utilisé.
\usepackage[skins,breakable]{tcolorbox}
\usepackage{graphicx}
\usepackage{colortbl}
\usepackage{booktabs}
\usepackage{tabularx}
\usepackage{diagbox} % case d'en-tête diagonale des tableaux à double entrée
% Colonnes extensibles centrée / gauche / droite (C, L, R), souvent utilisées par les modèles
\newcolumntype{C}{>{\centering\arraybackslash}X}
\newcolumntype{L}{>{\raggedright\arraybackslash}X}
\newcolumntype{R}{>{\raggedleft\arraybackslash}X}
\usepackage{array}
\usepackage{multicol}
\usepackage{siunitx}
% parse-numbers=false : accepte aussi \SI{2,5 \times 10^{-3}}{...}
\sisetup{locale=FR,output-decimal-marker={,},group-minimum-digits=4,parse-numbers=false}
\DeclareSIUnit\ohm{\ensuremath{\symup{\Omega}}} % U+2126 absent de la police
\DeclareSIUnit\division{div} % réglages d'oscilloscope (ms/div, V/div)
\DeclareSIUnit\tr{tr} % tours (vitesse de rotation, ex. \SI{1500}{\tr\per\minute})
\DeclareSIUnit\ch{ch} % cheval-vapeur (puissance moteur, unité usuelle hors SI)
\DeclareSIUnit\franccfa{FCFA} % franc CFA (coûts, contexte camerounais)
\DeclareSIUnit\dioptre{\ensuremath{\delta}} % vergence d'une lentille (dioptrie)
\usepackage{qrcode}
% \degree : commande fréquemment utilisée par les modèles pour un angle ou une
% température en dehors de \SI{}{} (non fournie par les paquets chargés).
\providecommand{\degree}{^{\circ}}
% \qed (fin de démonstration), utilisé par les modèles sans amsthm
\providecommand{\qed}{\hfill$\square$}
% \barre : double barre des valeurs interdites dans un tableau de variations (tkz-tab)
\providecommand{\barre}{\ensuremath{\|}}
% environnement proof (amsthm non chargé) utilisé par les modèles
\ifdefined\proof\else\newenvironment{proof}[1][Démonstration]{\par\noindent\textit{#1.}\ }{\qed\par}\fi
\usepackage{lastpage}
\usepackage{hyperref}
\hypersetup{hidelinks}

% ============================================================
% Palette « Pop » OnBuch+ — mêmes hex que la classe P (plus_ui.dart)
% ============================================================
\definecolor{poporange}{HTML}{F59321}
\definecolor{popdark}{HTML}{E07A0C}
\definecolor{popcream}{HTML}{FFF6E8}
\definecolor{popink}{HTML}{1C1714}
\definecolor{poppurple}{HTML}{7A5AE0}
\definecolor{popgreen}{HTML}{2E9E6A}
\definecolor{poppink}{HTML}{E0457B}
\definecolor{popblue}{HTML}{2F7DE1}
\definecolor{popgold}{HTML}{C98A12}
\definecolor{popmuted}{HTML}{8A7F76}
\definecolor{popline}{HTML}{EADFD2}
\definecolor{popgray}{HTML}{8A7F76}
\colorlet{popgrey}{popgray}
% Alias tolérants (le modèle nomme parfois les couleurs différemment)
\colorlet{popwhite}{white}
\colorlet{popblack}{popink}
\colorlet{popred}{poppink}
\colorlet{popyellow}{popgold}
% Teintes claires pour fonds de tableaux / zones
\colorlet{popblueL}{popblue!10!white}
\colorlet{poporangeL}{poporange!14!white}
\colorlet{popgreenL}{popgreen!12!white}
\colorlet{poppurpleL}{poppurple!12!white}
\colorlet{poppinkL}{poppink!12!white}
\colorlet{popgoldL}{popgold!14!white}

\pagecolor{popcream}

% ============================================================
% Typographie
% ============================================================
\defaultfontfeatures{Ligatures=TeX}
\setmainfont{PlusJakartaSans-Medium.ttf}[Path=../../../fonts/,BoldFont=PlusJakartaSans-Bold.ttf]
% Maths en sans-serif (Fira Math) avec chiffres et lettres droites de Plus
% Jakarta, pour que les formules se fondent dans le texte.
\usepackage[math-style=ISO,bold-style=ISO]{unicode-math}
\setmathfont{FiraMath-Regular.otf}[Path=../../../fonts/]
\setmathfont{PlusJakartaSans-Medium.ttf}[Path=../../../fonts/,range={up/num,up/{Latin,latin},\mathup}]
\usepackage{icomma} % virgule décimale sans espace en mode math
% Caractères mathématiques Unicode absents des polices de texte (Plus Jakarta,
% Archivo Black) : rendus par leur équivalent mathématique, en texte comme en
% formule — sinon ils s'affichent en carré vide (ex. « Arithmétique dans ℤ »).
\usepackage{newunicodechar}
\newunicodechar{ℝ}{\ensuremath{\mathbb{R}}}
\newunicodechar{ℤ}{\ensuremath{\mathbb{Z}}}
\newunicodechar{ℂ}{\ensuremath{\mathbb{C}}}
\newunicodechar{ℕ}{\ensuremath{\mathbb{N}}}
\newunicodechar{ℚ}{\ensuremath{\mathbb{Q}}}
\newunicodechar{𝔻}{\ensuremath{\mathbb{D}}}
\newunicodechar{α}{\ensuremath{\alpha}}
\newunicodechar{β}{\ensuremath{\beta}}
\newunicodechar{γ}{\ensuremath{\gamma}}
\newunicodechar{δ}{\ensuremath{\delta}}
\newunicodechar{ε}{\ensuremath{\varepsilon}}
\newunicodechar{θ}{\ensuremath{\theta}}
\newunicodechar{λ}{\ensuremath{\lambda}}
\newunicodechar{μ}{\ensuremath{\mu}}
\newunicodechar{σ}{\ensuremath{\sigma}}
\newunicodechar{τ}{\ensuremath{\tau}}
\newunicodechar{φ}{\ensuremath{\varphi}}
\newunicodechar{ψ}{\ensuremath{\psi}}
\newunicodechar{ω}{\ensuremath{\omega}}
\newunicodechar{Σ}{\ensuremath{\Sigma}}
% majuscules produites par \MakeUppercase dans les titres (α-aminés -> Α)
\newunicodechar{Α}{\ensuremath{\alpha}}
\newunicodechar{Β}{\ensuremath{\beta}}
\newunicodechar{Γ}{\ensuremath{\Gamma}}
\newunicodechar{Δ}{\ensuremath{\Delta}}
\newunicodechar{Ε}{\ensuremath{\varepsilon}}
\newunicodechar{Θ}{\ensuremath{\Theta}}
\newunicodechar{Λ}{\ensuremath{\Lambda}}
\newunicodechar{Μ}{\ensuremath{\mu}}
\newunicodechar{Τ}{\ensuremath{\tau}}
\newunicodechar{Φ}{\ensuremath{\Phi}}
\newunicodechar{Ψ}{\ensuremath{\Psi}}
\newunicodechar{Ω}{\ensuremath{\Omega}}
\newunicodechar{↦}{\ensuremath{\mapsto}}
\newunicodechar{∈}{\ensuremath{\in}}
\newunicodechar{∉}{\ensuremath{\notin}}
\newunicodechar{∧}{\ensuremath{\wedge}}
\newunicodechar{⇔}{\ensuremath{\Leftrightarrow}}
\newunicodechar{⟺}{\ensuremath{\Longleftrightarrow}}
\newunicodechar{⇒}{\ensuremath{\Rightarrow}}
\newunicodechar{⟹}{\ensuremath{\Longrightarrow}}
\newunicodechar{∘}{\ensuremath{\circ}}
\newunicodechar{∪}{\ensuremath{\cup}}
\newunicodechar{∩}{\ensuremath{\cap}}
\newunicodechar{≡}{\ensuremath{\equiv}}
\newunicodechar{⊂}{\ensuremath{\subset}}
\newunicodechar{⊥}{\ensuremath{\perp}}
\newunicodechar{∃}{\ensuremath{\exists}}
\newunicodechar{∀}{\ensuremath{\forall}}
\newunicodechar{∛}{\ensuremath{\sqrt[3]{\,}}}
\newunicodechar{∜}{\ensuremath{\sqrt[4]{\,}}}
\newunicodechar{⃗}{\ensuremath{{}^{\to}}}
\newunicodechar{ⁿ}{\ifmmode^{n}\else\textsuperscript{\ensuremath{n}}\fi}
\newunicodechar{ˣ}{\ifmmode^{x}\else\textsuperscript{\ensuremath{x}}\fi}
\newunicodechar{ᵇ}{\ifmmode^{b}\else\textsuperscript{\ensuremath{b}}\fi}
\newunicodechar{ʳ}{\ifmmode^{r}\else\textsuperscript{\ensuremath{r}}\fi}
\newunicodechar{ᵘ}{\ifmmode^{u}\else\textsuperscript{\ensuremath{u}}\fi}
\newunicodechar{ʸ}{\ifmmode^{y}\else\textsuperscript{\ensuremath{y}}\fi}
\newunicodechar{ᵃ}{\ifmmode^{a}\else\textsuperscript{\ensuremath{a}}\fi}
\newunicodechar{ᵉ}{\ifmmode^{e}\else\textsuperscript{\ensuremath{e}}\fi}
\newunicodechar{ᵏ}{\ifmmode^{k}\else\textsuperscript{\ensuremath{k}}\fi}
\newunicodechar{ᵖ}{\ifmmode^{p}\else\textsuperscript{\ensuremath{p}}\fi}
\newunicodechar{ᴺ}{\ifmmode^{N}\else\textsuperscript{\ensuremath{N}}\fi}
\newunicodechar{ᶜ}{\ifmmode^{c}\else\textsuperscript{\ensuremath{c}}\fi}
\newunicodechar{ʰ}{\ifmmode^{h}\else\textsuperscript{\ensuremath{h}}\fi}
\newunicodechar{ᵠ}{\ifmmode^{q}\else\textsuperscript{\ensuremath{q}}\fi}
\newunicodechar{ₙ}{\ifmmode_{n}\else\textsubscript{\ensuremath{n}}\fi}
\newunicodechar{ₐ}{\ifmmode_{a}\else\textsubscript{\ensuremath{a}}\fi}
\newunicodechar{ᵢ}{\ifmmode_{i}\else\textsubscript{\ensuremath{i}}\fi}
\newunicodechar{ᵣ}{\ifmmode_{r}\else\textsubscript{\ensuremath{r}}\fi}
\newunicodechar{ₖ}{\ifmmode_{k}\else\textsubscript{\ensuremath{k}}\fi}
\newunicodechar{ᵦ}{\ifmmode_{\beta}\else\textsubscript{\ensuremath{\beta}}\fi}
\newunicodechar{ₓ}{\ifmmode_{x}\else\textsubscript{\ensuremath{x}}\fi}
\newfontfamily\popdisplay{ArchivoBlack-Regular.ttf}[Path=../../../fonts/]
\newfontfamily\poplogo{SpaceGrotesk-Bold.ttf}[Path=../../../fonts/]
\newfontfamily\popsemi{PlusJakartaSans-SemiBold.ttf}[Path=../../../fonts/]
\newfontfamily\popxbold{PlusJakartaSans-ExtraBold.ttf}[Path=../../../fonts/]
\newfontfamily\popxboldls{PlusJakartaSans-ExtraBold.ttf}[Path=../../../fonts/,LetterSpace=3.0]

\setlist[enumerate]{leftmargin=1.7em,itemsep=5pt,topsep=4pt}
\setlist[itemize]{leftmargin=1.7em,itemsep=4pt,topsep=4pt,label={\color{poporange}\textbullet}}
\setlength{\parindent}{0pt}
\setlength{\parskip}{5pt}
\renewcommand{\arraystretch}{1.25}

% pgfplots : style Pop par défaut
\pgfplotsset{
  every axis/.append style={axis line style={popink,line width=1pt},tick style={popink},
    grid style={popline,line width=0.5pt},label style={font=\small},tick label style={font=\footnotesize}},
  popcurve/.style={poporange,line width=1.8pt,no markers,smooth},
}
% Même style disponible dans un \draw TikZ simple (hors pgfplots)
\tikzset{popcurve/.style={poporange,line width=1.8pt}}
\tikzset{popgrid/.style={popline,very thin}}
\colorlet{popcurve}{poporange} % le nom du style est parfois utilisé comme couleur

% ============================================================
% Pill / squiggle
% ============================================================
\newcommand{\poppill}[3]{%
  \tikz[baseline=-0.55ex]\node[draw=popink,line width=1.2pt,rounded corners=2.6pt,%
    fill=#2,inner xsep=6.5pt,inner ysep=3pt]{%
    \popxboldls\fontsize{7}{7}\selectfont\color{#3}\MakeUppercase{#1}};%
}
\newcommand{\popsquiggle}[1]{%
  \tikz[baseline=-0.4ex]{
    \draw[#1,line width=2.6pt,line cap=round]
      plot [smooth] coordinates {(0,0.09) (0.34,0.01) (0.68,0.21) (1.02,0.01) (1.36,0.10)};
  }%
}
% Mot-clé surligné (marqueur orange sous le texte)
\newcommand{\cle}[1]{{\popxbold\color{popdark}#1}}

% ============================================================
% En-tête / pied
% ============================================================
\pagestyle{fancy}
\fancyhf{}
\renewcommand{\headrulewidth}{0pt}
\renewcommand{\footrulewidth}{0pt}
\lhead{\raisebox{-0.15ex}{\poplogo\fontsize{13}{13}\selectfont\color{popink}OnBuch\color{poporange}+}}
\rhead{\poppill{\DOCMATIERE\ \textbullet\ \DOCNIVEAU\ \textbullet\ Leçon \DOCLECON}{white}{popink}}
\lfoot{\popxbold\fontsize{7.6}{7.6}\selectfont\color{popmuted}\MakeUppercase{Cours signé OnBuch+}\ \textbullet\ {\popsemi\DOCTITRE}}
\rfoot{\tikz[baseline=-0.6ex]\node[rounded corners=8.5pt,fill=popink,minimum height=17pt,minimum width=17pt,inner xsep=5pt,inner sep=0]{\popxbold\fontsize{8}{8}\selectfont\color{popcream}\thepage\,/\,\pageref*{LastPage}};}

% ============================================================
% Titres
% ============================================================
\newcommand{\chaptitre}[2]{%
  \begin{tcolorbox}[enhanced,colback=poporange,colframe=popink,boxrule=2.6pt,arc=11pt,
      drop shadow=popink,left=15pt,right=15pt,top=12pt,bottom=11pt,before skip=3pt,after skip=18pt]
    \poppill{#2}{popink}{popcream}\par\vspace{9pt}
    {\raggedright\hyphenpenalty=10000\exhyphenpenalty=10000\popdisplay\fontsize{22}{24}\selectfont\color{popcream}\MakeUppercase{#1}\par}
  \end{tcolorbox}
}
% Section numérotée : pastille orange avec le numéro + titre Archivo
\newcounter{popsec}
\newcommand{\coursec}[1]{%
  \stepcounter{popsec}\setcounter{popsub}{0}%
  \par\vspace{16pt}\noindent
  \begin{minipage}{\linewidth}
  \tikz[baseline=-0.7ex]\node[draw=popink,line width=1.6pt,fill=poporange,rounded corners=4pt,minimum size=22pt,inner sep=0]{\popdisplay\fontsize{12}{12}\selectfont\color{popcream}\thepopsec};%
  \hspace{8pt}{\popdisplay\fontsize{15}{17}\selectfont\color{popink}\MakeUppercase{#1}}\par
  \vspace{4pt}\noindent{\color{poporange}\rule{36pt}{3.6pt}}
  \end{minipage}\par\nopagebreak\vspace{8pt}
}
\newcounter{popsub}
\newcommand{\courssub}[1]{%
  \stepcounter{popsub}%
  \par\vspace{9pt}\noindent{\popxbold\fontsize{11.5}{13}\selectfont\color{popink}{\color{poporange}\thepopsec.\thepopsub}\hspace{6pt}#1}\par\nopagebreak\vspace{4pt}
}

% ============================================================
% Boîtes pédagogiques — même recette : fond blanc, bordure épaisse colorée,
% ombre dure encre, badge plein. Titre optionnel [#1].
% ============================================================
\tcbset{popbase/.style={breakable,enhanced,colback=white,boxrule=2.3pt,arc=9pt,
  shadow={3pt}{-3pt}{0pt}{popink},left=14pt,right=14pt,top=11pt,bottom=11pt,before skip=11pt,after skip=15pt,
  fontupper=\popsemi\fontsize{10.6}{14.5}\selectfont\color{popink},
  fontlower=\popsemi\fontsize{10.6}{14.5}\selectfont\color{popink}}}
% Coche dessinée (le glyphe ✓ n'existe pas dans les polices du document)
\newcommand{\popcheck}[1]{\tikz[baseline=-0.5ex]\draw[#1,line width=1.6pt,line cap=round,line join=round](-0.13,0.0)--(-0.04,-0.09)--(0.13,0.1);}
\providecommand{\coche}{\popcheck{popgreen}}
\providecommand{\resultat}[1]{\par\smallskip\noindent\fcolorbox{popgreen}{popgreenL}{\parbox{\dimexpr\linewidth-2\fboxsep-2\fboxrule\relax}{\textbf{Résultat.} #1}}\par\smallskip}
\newcommand{\popboxhead}[3]{\poppill{#1}{#2}{white}\ifx\relax#3\relax\else\hspace{6pt}{\popxbold\fontsize{10.6}{12}\selectfont\color{popink}#3}\fi\par\vspace{7pt}}

\newtcolorbox{aretenir}[1][]{popbase,colframe=popgreen,before upper={\popboxhead{À retenir}{popgreen}{#1}}}
\newtcolorbox{exemplebox}[1][]{popbase,colframe=poporange,before upper={\popboxhead{Exemple}{poporange}{#1}}}
\newtcolorbox{definition}[1][]{popbase,colframe=popblue,before upper={\popboxhead{Définition}{popblue}{#1}}}
\newtcolorbox{propriete}[1][]{popbase,colframe=poppurple,before upper={\popboxhead{Propriété}{poppurple}{#1}}}
\newtcolorbox{methode}[1][]{popbase,colframe=popgold,before upper={\popboxhead{Méthode}{popgold}{#1}}}
\newtcolorbox{attention}[1][]{popbase,colframe=poppink,before upper={\popboxhead{Attention}{poppink}{#1}}}
\newtcolorbox{experience}[1][]{popbase,colframe=popblue,colback=popblueL,before upper={\popboxhead{Expérience}{popblue}{#1}}}
% « Le savais-tu ? » : carte violette pleine (contexte, culture, Cameroun)
\newtcolorbox{savaistu}[1][]{popbase,colback=poppurple,colframe=popink,
  fontupper=\popsemi\fontsize{10.4}{14.2}\selectfont\color{white},
  before upper={\poppill{Le savais-tu\,?}{white}{poppurple}\ifx\relax#1\relax\else\hspace{6pt}{\popxbold\color{white}#1}\fi\par\vspace{7pt}}}
% Exercice résolu : énoncé en haut, \tcblower puis la solution sur fond crème
\newcounter{popexo}\newcounter{popexr}
\newtcolorbox{exoresolu}[1][]{popbase,colframe=poporange,colbacklower=popcream,
  segmentation style={popink,line width=1.2pt,dashed},
  before upper={\stepcounter{popexr}\popboxhead{Exercice résolu \thepopexr}{poporange}{#1}},
  before lower={\poppill{Solution}{popink}{popcream}\par\vspace{6pt}}}
% Exercice (non résolu en place) — la correction est dans la partie Corrigés
\newtcolorbox{exercice}[1][]{popbase,colframe=popink,
  before upper={\stepcounter{popexo}\popboxhead{Exercice \thepopexo}{popink}{#1}}}
% Corrigé
\newtcolorbox{corrige}[1][]{popbase,colframe=popgreen,colback=popgreenL,
  before upper={\popboxhead{Corrigé}{popgreen}{#1}}}
% Objectifs / prérequis / bilan
\newtcolorbox{objectifs}{popbase,colframe=popink,colback=poporangeL,
  before upper={\popboxhead{Objectifs de la leçon}{popink}{}}}
\newtcolorbox{prerequis}{popbase,colframe=popink,colback=popblueL,
  before upper={\popboxhead{Prérequis}{popblue}{}}}
\newtcolorbox{bilan}{popbase,colframe=popink,colback=white,boxrule=2.8pt,
  before upper={\popboxhead{Fiche bilan}{poporange}{L'essentiel à connaître pour l'examen}}}
% Figure encadrée avec légende
\newtcolorbox{popfigure}[1][]{enhanced,breakable=false,colback=white,colframe=popline,boxrule=1.4pt,arc=8pt,
  left=10pt,right=10pt,top=10pt,bottom=8pt,before skip=10pt,after skip=12pt,halign=center,
  lower separated=false,halign lower=center,
  fontlower=\popsemi\fontsize{9}{11.5}\selectfont\color{popmuted},
  #1}
\newcommand{\legende}[1]{\par\vspace{6pt}{\popsemi\fontsize{9}{11.5}\selectfont\color{popmuted}#1}}

% ============================================================
% Signature OnBuch+
% ============================================================
\newcommand{\poplogoinline}[1]{{\poplogo\fontsize{#1}{#1}\selectfont\color{popink}OnBuch\color{poporange}+}}
\newcommand{\signatureonbuch}{%
  \begin{tcolorbox}[enhanced,colback=popink,colframe=popink,boxrule=2.2pt,arc=10pt,
      drop shadow=poporange,left=14pt,right=14pt,top=10pt,bottom=10pt,before skip=0pt,after skip=0pt]
    \tikz[baseline=-0.7ex]\node[circle,fill=popgreen,draw=popcream,line width=1.3pt,minimum size=22pt,inner sep=0]{\popcheck{white}};%
    \hspace{9pt}\begin{minipage}[c]{0.62\linewidth}
      {\popxbold\fontsize{12}{13}\selectfont\color{popcream}Signé\ \textbullet\ {\poplogo OnBuch\color{poporange}+}}\par\vspace{2pt}
      {\raggedright\popsemi\fontsize{8}{10.5}\selectfont\color{popline}\MakeUppercase{Équipe pédagogique OnBuch+}\\\MakeUppercase{Conforme au programme officiel MINESEC}\par}
    \end{minipage}\hfill
    {\poplogo\fontsize{22}{22}\selectfont\color{popcream}OnBuch\color{poporange}+}
  \end{tcolorbox}%
}

% ============================================================
% Page de garde
% #1 titre · #2 matière · #3 niveau · #4 module · #5 n° leçon · #6 durée ·
% #7 description / objectifs (texte ou itemize)
% ============================================================
\newcommand{\pagegarde}[7]{%
  \thispagestyle{empty}
  \newgeometry{a4paper,margin=1.7cm,top=1.6cm,bottom=1.4cm}%
  \begin{tikzpicture}[remember picture,overlay]
    \fill[poporange!18] ([xshift=-3.2cm,yshift=-3.6cm]current page.north east) circle (4.6cm);
    \fill[poppurple!14] ([xshift=2.4cm,yshift=7.6cm]current page.south west) circle (3.8cm);
  \end{tikzpicture}%
  {\poplogo\fontsize{30}{30}\selectfont\color{popink}OnBuch\color{poporange}+}\hfill
  \raisebox{0.6ex}{\poppill{Cours signé \textbullet\ Premium}{popink}{popcream}}\par
  \vspace{1pt}\popsquiggle{poppurple}\par
  \vspace{8pt}
  \begin{tcolorbox}[enhanced,colback=poporange,colframe=popink,boxrule=2.8pt,arc=13pt,
      drop shadow=popink,left=18pt,right=18pt,top=12pt,bottom=13pt,after skip=10pt]
    \poppill{#2 \textbullet\ #3}{popink}{popcream}\hspace{5pt}\poppill{#4}{white}{popink}\par\vspace{8pt}
    {\popdisplay\fontsize{14}{14}\selectfont\color{popink}LEÇON #5}\par\vspace{4pt}
    {\raggedright\hyphenpenalty=10000\exhyphenpenalty=10000\popdisplay\fontsize{25}{27}\selectfont\color{popcream}\MakeUppercase{#1}\par}
    \vspace{8pt}
    {\popxbold\fontsize{9.5}{9.5}\selectfont\color{popink}\MakeUppercase{Durée conseillée : #6}}
  \end{tcolorbox}
  \begin{tcolorbox}[enhanced,colback=white,colframe=popink,boxrule=2.2pt,arc=10pt,
      drop shadow=popink,left=16pt,right=16pt,top=10pt,bottom=10pt,after skip=8pt]
    \poppill{Dans cette leçon}{poppurple}{white}\par\vspace{5pt}
    \popsemi\fontsize{9.3}{12.6}\selectfont\color{popink}#7
  \end{tcolorbox}
  \noindent\begin{minipage}[t]{0.487\linewidth}
    \begin{tcolorbox}[enhanced,colback=popgreen,colframe=popink,boxrule=2.2pt,arc=10pt,
        drop shadow=popink,left=11pt,right=11pt,top=10pt,bottom=10pt,height=3.1cm,valign=center]
      \tcbox[colback=white,colframe=popink,boxrule=1.3pt,arc=5pt,boxsep=0pt,nobeforeafter,
        left=3pt,right=3pt,top=3pt,bottom=3pt,valign=center]{\qrcode[height=1.9cm]{https://play.google.com/store/apps/details?id=com.luvvix.onbuch&hl=fr}}%
      \hspace{8pt}\begin{minipage}[c]{0.5\linewidth}
        {\popdisplay\fontsize{11}{12.5}\selectfont\color{white}\MakeUppercase{Télécharge OnBuch}}\par\vspace{4pt}
        {\raggedright\popsemi\fontsize{8.4}{11}\selectfont\color{white}Gratuit \textbullet\ Google Play\\Tuteur IA, annales, concours, résultats\par}
      \end{minipage}
    \end{tcolorbox}
  \end{minipage}\hfill
  \begin{minipage}[t]{0.487\linewidth}
    \begin{tcolorbox}[enhanced,colback=poppurple,colframe=popink,boxrule=2.2pt,arc=10pt,
        drop shadow=popink,left=11pt,right=11pt,top=10pt,bottom=10pt,height=3.1cm,valign=center]
      \tikz[baseline=-0.7ex]\node[circle,fill=white,draw=popink,line width=1.3pt,minimum size=1.6cm,inner sep=0]{\popdisplay\fontsize{24}{24}\selectfont\color{poppurple}+};%
      \hspace{8pt}\begin{minipage}[c]{0.6\linewidth}
        {\popdisplay\fontsize{11}{12.5}\selectfont\color{white}\MakeUppercase{Passe à OnBuch+}}\par\vspace{4pt}
        {\raggedright\popsemi\fontsize{8.4}{11}\selectfont\color{white}Tous les cours signés, Tuteur IA illimité, fiches PDF — onglet OnBuch+ de l'app\par}
      \end{minipage}
    \end{tcolorbox}
  \end{minipage}
  \vspace{8pt}
  \popcontenu
  \vfill
  \signatureonbuch
  \restoregeometry
  \newpage
}

% Rangée « Ce cours contient » (page de garde)
\newcommand{\poptile}[3]{% #1 couleur #2 gros texte #3 légende
  \begin{tcolorbox}[enhanced,colback=#1,colframe=popink,boxrule=2pt,arc=9pt,shadow={2.5pt}{-2.5pt}{0pt}{popink},
      left=6pt,right=6pt,top=9pt,bottom=9pt,halign=center,valign=center,height=1.75cm,width=\linewidth,nobeforeafter]
    {\popdisplay\fontsize{12.5}{13}\selectfont\color{white}\MakeUppercase{#2}}\par\vspace{3pt}
    {\centering\popsemi\fontsize{7.8}{9.5}\selectfont\color{white}#3\par}
  \end{tcolorbox}}
\newcommand{\popcontenu}{%
  \par\noindent{\popxboldls\fontsize{7.5}{7.5}\selectfont\color{popmuted}CE COURS SIGNÉ CONTIENT}\par\vspace{7pt}
  \noindent\begin{minipage}[t]{0.235\linewidth}\poptile{popblue}{Cours}{complet, illustré, pas à pas}\end{minipage}\hfill
  \begin{minipage}[t]{0.235\linewidth}\poptile{poporange}{Méthodes}{et exercices résolus}\end{minipage}\hfill
  \begin{minipage}[t]{0.235\linewidth}\poptile{poppink}{Exercices}{type examen + corrigés}\end{minipage}\hfill
  \begin{minipage}[t]{0.235\linewidth}\poptile{popgreen}{Fiche bilan}{l'essentiel à retenir}\end{minipage}\par}

% Sommaire
\newtcolorbox{sommaire}{enhanced,colback=white,colframe=popink,boxrule=2.2pt,arc=10pt,
  drop shadow=popink,left=18pt,right=18pt,top=13pt,bottom=13pt,before skip=6pt,after skip=20pt,
  before upper={\poppill{Sommaire}{poppurple}{white}\par\vspace{9pt}\popsemi\fontsize{10.6}{14.5}\selectfont\color{popink}}}

% Signature de fin de cours — poussée tout en bas de la dernière page
\newcommand{\findecours}{%
  \par\vfill\nopagebreak
  \begin{center}\popsquiggle{poporange}\end{center}\vspace{4pt}
  \signatureonbuch
}
\AtBeginDocument{\def\llbracket{\mathopen{[\mkern-3.5mu[}}\def\rrbracket{\mathclose{]\mkern-3.5mu]}}} % intervalles d'entiers (glyphe ⟦ absent de la police)
% Glyphes absents de Fira Math : redéfinis à partir de glyphes disponibles
\AtBeginDocument{%
  \def\setminus{\mathbin{\backslash}}\let\smallsetminus\setminus
  \def\vdots{\mathord{\vbox{\baselineskip3.2pt\lineskiplimit0pt\kern4pt\hbox{.}\hbox{.}\hbox{.}}}}%
  \def\ddots{\mathinner{\mkern1mu\raise6.5pt\hbox{.}\mkern2mu\raise3.6pt\hbox{.}\mkern2mu\raise0.7pt\hbox{.}\mkern1mu}}%
  \def\triangle{\mathord{\scalebox{0.9}{$\symup{\Delta}$}}}%
  \def\nabla{\mathord{\raisebox{\depth}{\scalebox{1}[-1]{$\symup{\Delta}$}}}}%
  \def\checkmark{\popcheck{popdark}}%
  \def\star{\mathbin{\ast}}\def\bigstar{\mathord{\ast}}%
  \def\diamond{\mathbin{\scalebox{0.6}{$\Diamond$}}}%
  \def\bigcup{\mathop{\mathchoice{\raisebox{-0.3em}{\scalebox{1.7}{$\cup$}}}{\scalebox{1.25}{$\cup$}}{\cup}{\cup}}\displaylimits}%
  \def\bigcap{\mathop{\mathchoice{\raisebox{-0.3em}{\scalebox{1.7}{$\cap$}}}{\scalebox{1.25}{$\cap$}}{\cap}{\cap}}\displaylimits}%
  \def\lesseqqgtr{\mathrel{\lessgtr}}\def\bigoplus{\mathop{\oplus}\displaylimits}%
}
\providecommand{\ssi}{si et seulement si} % abréviation fréquente
\AtBeginDocument{\providecommand{\wideparen}{\overparen}} % arc de cercle

%% FIGFIX-BEGIN v5 — correctifs globaux des figures (docs/onbuch-generation/tools/figfix)
\usepackage{adjustbox}\usepackage{etoolbox}\tcbuselibrary{raster}
% toute figure TikZ/pgfplots plus large que la ligne est réduite à la largeur disponible
\BeforeBeginEnvironment{tikzpicture}{\begin{adjustbox}{max width=\linewidth}}
\AfterEndEnvironment{tikzpicture}{\end{adjustbox}}
% légendes pgfplots lisibles par-dessus les courbes
\pgfplotsset{every axis/.append style={legend style={fill=white,fill opacity=0.92,text opacity=1,draw=black!15,inner sep=3pt}}}
% étiquettes d'arêtes (syntaxe edge["..."]) avec fond blanc
\tikzset{every edge quotes/.append style={fill=white,inner sep=1.2pt}}
%% FIGFIX-END
\begin{document}

\pagegarde{Le télescope de Newton}{Physique}{1ère TI}{Module 3 : Optique}{N11}{4 h}{Cette leçon présente le télescope de Newton : sa description (miroir primaire concave, miroir secondaire plan, oculaire), son principe de fonctionnement et la marche des rayons lumineux. On y définit le grossissement et la puissance, et on montre l'intérêt du télescope à réflexion par rapport aux lunettes à lentilles.
\vspace{4pt}
\begin{multicols}{2}\small
\begin{itemize}
\item À la fin de cette leçon, je sais décrire les éléments constitutifs d'un télescope de Newton (miroir primaire, miroir secondaire, oculaire).
\item À la fin de cette leçon, je sais expliquer le rôle de chaque élément dans la formation des images.
\item À la fin de cette leçon, je sais tracer la marche d'un faisceau lumineux issu d'un astre à travers un télescope de Newton.
\item À la fin de cette leçon, je sais définir et calculer le grossissement d'un télescope afocal.
\item À la fin de cette leçon, je sais définir la puissance d'un instrument d'optique et l'exprimer en dioptries.
\item À la fin de cette leçon, je sais faire la mise au point d'un télescope de Newton (réglage de l'oculaire, collimation).
\end{itemize}\end{multicols}}

\begin{sommaire}
\begin{multicols}{2}
\begin{enumerate}
\item Description du télescope de Newton
\item Principe de fonctionnement et formation des images
\item Tracé de la marche d'un faisceau lumineux
\item Grossissement et puissance
\item Mise au point du télescope de Newton
\item Intérêt du télescope par rapport aux instruments à lentilles
\item Activité d'intégration
\item Méthodes et exercices résolus
\item Exercices
\item Corrigés des exercices
\item Fiche bilan
\end{enumerate}
\end{multicols}
\end{sommaire}

\begin{objectifs}
\begin{itemize}
\item À la fin de cette leçon, je sais décrire les éléments constitutifs d'un télescope de Newton (miroir primaire, miroir secondaire, oculaire).
\item À la fin de cette leçon, je sais expliquer le rôle de chaque élément dans la formation des images.
\item À la fin de cette leçon, je sais tracer la marche d'un faisceau lumineux issu d'un astre à travers un télescope de Newton.
\item À la fin de cette leçon, je sais définir et calculer le grossissement d'un télescope afocal.
\item À la fin de cette leçon, je sais définir la puissance d'un instrument d'optique et l'exprimer en dioptries.
\item À la fin de cette leçon, je sais faire la mise au point d'un télescope de Newton (réglage de l'oculaire, collimation).
\item À la fin de cette leçon, je sais comparer le télescope de Newton aux instruments à lentilles (lunette astronomique) et citer ses avantages.
\item À la fin de cette leçon, je sais résoudre un exercice type Probatoire mêlant construction géométrique et calcul de grossissement.
\end{itemize}
\end{objectifs}

\begin{prerequis}
\begin{itemize}
\item Réflexion de la lumière sur un miroir plan et lois de Snell-Descartes (Seconde).
\item Miroir sphérique concave : foyer, distance focale, construction des images (leçons précédentes du module Optique).
\item Lentilles minces convergentes : vergence, formules de conjugaison, construction graphique des images.
\item Système afocal et lunette astronomique : grossissement G = f'₁/f'₂.
\item Notions d'angle apparent (diamètre apparent d'un astre) et d'unités d'angle (degré, minute, seconde).
\end{itemize}
\end{prerequis}

\newpage

\coursec{Le télescope de Newton}

Lors des nuits sans lune au village, loin des éclairages publics de Yaoundé ou de Douala, on distingue à l'œil nu quelques milliers d'étoiles, mais la Voie lactée reste une traînée floue et les satellites de Jupiter invisibles. Un astronome amateur de Yaoundé construit un télescope avec un miroir concave récupéré et un petit miroir plan : soudain, les anneaux de Saturne apparaissent distinctement ! Pourquoi remplacer les lentilles par un miroir ? Comment un tel instrument grossit-il les astres ? C'est toute la leçon du télescope de Newton.

\courssub{Pourquoi l'observation des astres pose un problème physique fondamental}

Les objets célestes (étoiles, planètes, nébuleuses, galaxies) sont des \emph{objets ponctuels} ou \emph{étendus à distance quasi infinie}. Trois difficultés majeures s'opposent à leur observation directe :

\begin{enumerate}
    \item \textbf{Un diamètre apparent extrêmement faible.} L'angle sous-tendu par l'astre à l'œil est minuscule. Par exemple, Jupiter, au plus près de la Terre, présente un diamètre apparent d'environ \SI{50}{''} (secondes d'angle), soit $\approx \SI{2,4e-4}{rad}$. L'œil nu (puissance de résolution $\approx \SI{1}{'} = \SI{60}{''}$) ne le voit que comme un point brillant.
    \item \textbf{Une illuminance (flux lumineux par unité de surface) très faible.} La luminance de l'objet est conservée par un système optique parfait (théorème de Lagrange-Helmholtz), mais l'illuminance sur la rétine dépend de la surface de la pupille d'entrée de l'instrument. L'œil nu a une pupille de diamètre $d_{\text{œil}} \approx \SI{2}{mm}$ (jour) à $\SI{7}{mm}$ (nuit). Pour voir des détails fins sur un objet faiblement lumineux, il faut \emph{collecter plus de lumière} : d'où la nécessité d'une grande ouverture.
    \item \textbf{La turbulence atmosphérique (voirie)} qui limite la résolution angulaire effective au sol à environ \SI{1}{''} dans les meilleurs sites (comme le mont Cameroun ou l'Observatoire de l'Université de Yaoundé I), bien loin du pouvoir séparateur théorique des grands instruments.
\end{enumerate}

\begin{aretenir}[Les deux fonctions vitales d'un instrument d'astronomie]
Un instrument d'optique pour l'astronomie doit remplir deux rôles distincts :
\begin{itemize}
    \item \textbf{Fonction « collecteur de lumière » (objectif) :} Augmenter la surface d'entrée du flux lumineux pour détecter des objets plus faiblement lumineux et augmenter la résolution angulaire (pouvoir séparateur $\propto 1/D$).
    \item \textbf{Fonction « grossisseur angulaire » (oculaire) :} Augmenter le diamètre apparent de l'image formée par l'objectif pour la rendre exploitable par l'œil.
\end{itemize}
\end{aretenir}

\courssub{Naissance du télescope à réflexion : l'ingéniosité de Newton}

Au XVII\textsuperscript{e} siècle, les lunettes astronomiques (réfracteurs) de Galilée puis de Kepler utilisent des lentilles convergentes comme objectif. Mais le verre disponible à l'époque présente une \textbf{dispersion} forte : l'indice de réfraction $n$ varie avec la longueur d'onde $\lambda$ ($n_{\text{bleu}} > n_{\text{rouge}}$). La distance focale $f' = \frac{R}{n-1}$ (formule du verrier pour lentille mince) est donc plus courte pour le bleu que pour le rouge. Une source ponctuelle blanche donne une image colorée, floue, entourée de franges arc-en-ciel : c'est l'\textbf{aberration chromatique}.

\begin{savaistu}[Isaac Newton et la fin de l'aberration chromatique (1668)]
En 1666, le jeune Isaac Newton (1643--1727) décompose la lumière blanche avec un prisme et comprend que la réfraction \emph{sépare} les couleurs. Il en déduit qu'\emph{aucun perfectionnement des lentilles ne peut supprimer l'aberration chromatique} (théorème de l'impossibilité d'un réfracteur achromatique avec un seul verre). Il imagine alors de remplacer la lentille objectif par un \textbf{miroir concave} : la réflexion suit la loi de Descartes ($i = r$), indépendante de la longueur d'onde. Toutes les couleurs convergent en un même foyer. En 1668, il construit le premier télescope à réflexion fonctionnel (miroir de \SI{25}{mm}, $F = \SI{150}{mm}$), qu'il présente à la Royal Society en 1672. C'est la naissance du \textbf{télescope de Newton}.
\end{savaistu}

\courssub{Constitution et rôle de chaque élément}

Le télescope de Newton est un instrument \textbf{catadioptrique} (réflexion + réfraction) ou plus précisément un \textbf{télescope à réflexion} pur pour la partie objectif, l'oculaire restant une lentille.

\begin{definition}[Télescope de Newton]
Un \textbf{télescope de Newton} est un instrument d'optique à réflexion constitué de trois éléments principaux :
\begin{enumerate}
    \item Un \textbf{miroir primaire concave} (souvent parabolique), placé au fond du tube, qui joue le rôle d'objectif : il collecte la lumière et forme une image réelle de l'astre dans son plan focal.
    \item Un \textbf{miroir secondaire plan}, incliné à \SI{45}{\degree} par rapport à l'axe optique, placé avant le foyer du primaire, qui renvoie le faisceau convergent latéralement hors de l'axe du tube.
    \item Un \textbf{oculaire} (lentille mince convergente ou système de lentilles), placé sur le côté du tube, qui joue le rôle de loupe pour observer l'image intermédiaire formée par le primaire.
\end{enumerate}
Le tube est ouvert du côté pointé vers l'astre.
\end{definition}

\begin{popfigure}
\begin{tikzpicture}[scale=0.9, >=Stealth]
% Couleurs
\colorlet{tube}{popdark!80!popblue}
\colorlet{mirror}{popblue!80!popdark}
\colorlet{secondary}{poporange}
\colorlet{ray}{poppink}
\colorlet{eye}{popgreen}

% Tube (coupe longitudinale)
\draw[tube, line width=2pt] (-0.5, -4) -- (-0.5, 3) -- (0.5, 3) -- (0.5, -4) -- cycle;
\draw[dashed, thin, popmuted] (0, -4) -- (0, 3); % Axe optique principal

% Miroir Primaire (Concave parabolique approximé par arc de cercle)
\draw[mirror, line width=3pt] (-0.6, -4) .. controls (0, -3.5) .. (0.6, -4);
\draw[mirror, line width=1pt] (-0.6, -4.1) -- (0.6, -4.1); % Support arrière
\node[mirror, anchor=north] at (0, -4.2) {\textbf{Miroir primaire concave (parabolique)}};
\node[popdark, anchor=north] at (0, -4.6) {Foyer principal $F'$};

% Miroir Secondaire (Plan à 45°)
\draw[secondary, line width=2.5pt] (-0.8, 1.5) -- (1.2, -0.5);
\draw[secondary, line width=1pt, fill=popcream] (-0.9, 1.6) -- (1.3, -0.4) -- (1.4, -0.3) -- (-0.8, 1.5) -- cycle;
\node[secondary, anchor=south west] at (1.3, -0.3) {\textbf{Miroir secondaire plan (45°)}};

% Oculaire (Lentille convergente côté droit)
\draw[poppurple, line width=2pt] (2.5, 1.2) .. controls (3.0, 1.2) and (3.0, 0.8) .. (2.5, 0.8);
\draw[poppurple, line width=2pt] (2.5, 1.2) .. controls (3.0, 1.2) and (3.0, 1.6) .. (2.5, 1.6);
\draw[poppurple, line width=1pt] (2.5, 0.8) -- (2.5, 1.6);
\node[poppurple, anchor=west] at (2.7, 1.4) {\textbf{Oculaire (lentille convergente)}};

% Œil observateur
\draw[eye, line width=1.5pt] (3.5, 1.4) ellipse (0.3 and 0.15);
\draw[eye, line width=1.5pt] (3.5, 1.4) circle (0.1);
\node[eye, anchor=west] at (3.9, 1.4) {\textbf{Œil}};

% Faisceaux lumineux (étoiles à l'infini -> parallèles)
\foreach \y in {2.5, 2.0, 1.5} {
    \draw[ray, ->, thick] (-1.5, \y) -- (-0.5, \y);
}
% Réflexion sur primaire -> convergence vers F'
\foreach \y/\yf in {2.5/0.5, 2.0/0.0, 1.5/-0.5} {
    \draw[ray, ->, thick] (-0.5, \y) .. controls (0, \yf) .. (0.5, \yf);
}
% Réflexion sur secondaire (45°) -> horizontal vers oculaire
\foreach \yf in {0.5, 0.0, -0.5} {
    \draw[ray, ->, thick] (0.5, \yf) -- (1.0, \yf); % segment avant miroir
    % Calcul intersection miroir secondaire (droite y = -x + 0.7)
    % Rayon incident horizontal y = yf. Intersection x = 0.7 - yf.
    \pgfmathsetmacro{\xi}{0.7 - \yf}
    \draw[ray, ->, thick] (1.0, \yf) -- (\xi, \yf);
    % Réflexion : incidence 45, réflexion 45 -> vertical vers le haut
    \draw[ray, ->, thick] (\xi, \yf) -- (\xi, 1.4);
    % Vers oculaire
    \draw[ray, ->, thick] (\xi, 1.4) -- (2.5, 1.4);
}

% Cotes
\draw[|<->|, thin, popdark] (-0.5, 3.2) -- (0.5, 3.2) node[midway, above] {$D$ (Ouverture)};
\draw[|<->|, thin, popdark] (0, -4) -- (0, 0.7) node[midway, right] {$F$ (Distance focale primaire)};

% Légende axe dévié
\draw[<->, thin, popmuted] (1.5, 1.5) -- (1.5, 0.5) node[midway, right, align=center]{Axe optique\\ dévié à 90°};

\end{tikzpicture}
\legende{Schéma en coupe d'un télescope de Newton : trajet d'un faisceau parallèle (étoile) jusqu'à l'œil. Le miroir primaire (concave) converge la lumière vers son foyer $F'$ ; le miroir secondaire (plan à 45°) renvoie l'image latéralement vers l'oculaire.}
\end{popfigure}

\paragraph{Rôle du miroir primaire (objectif).}
C'est l'élément \textbf{maître}. C'est un miroir concave, le plus souvent \textbf{parabolique} (et non sphérique) pour éviter l'aberration sphérique pour les objets à l'infini. Sa surface réfléchissante (aluminium ou argent protégé par une couche de quartz) fait face à l'entrée du tube.
\begin{itemize}
    \item \textbf{Collecte de lumière :} Sa surface utile est un disque de diamètre $D$ (l'\textbf{ouverture}). Le flux lumineux collecté est proportionnel à la surface $\mathcal{S} = \pi D^2/4$. Un miroir de $D = \SI{200}{mm}$ collecte $\left(\frac{200}{7}\right)^2 \approx 800$ fois plus de lumière que l'œil nu (pupille $\SI{7}{mm}$).
    \item \textbf{Formation de l'image :} Pour un objet à l'infini (astre), les rayons incidents sont parallèles à l'axe optique. Le miroir les converge vers son foyer réel $F'$. L'image de l'astre (réelle, inversée, réduite) se forme dans le \textbf{plan focal} du miroir primaire, à une distance $F = F'P$ du sommet $P$ du miroir ($F > 0$ pour un miroir concave). C'est l'\textbf{image intermédiaire}.
\end{itemize}

\paragraph{Rôle du miroir secondaire (plan).}
C'est un petit miroir plat (souvent elliptique vu de face pour s'inscrire dans le faisceau conique), incliné à \SI{45}{\degree} par rapport à l'axe principal.
\begin{itemize}
    \item Il ne \textbf{modifie pas la convergence} du faisceau (miroir plan : distance focale infinie, puissance nulle). L'image intermédiaire reste à la même distance optique du secondaire.
    \item Il \textbf{dévie le faisceau de \SI{90}{\degree}} vers le côté du tube.
    \item \textbf{Avantage crucial :} L'observateur se place sur le côté, \textbf{ne bouchant pas l'entrée du tube} (contrairement au télescope de Gregory ou Cassegrain où l'œil est sur l'axe, ou aux lunettes où l'objectif est devant). L'obstruction centrale (ombre du secondaire et de son support) réduit légèrement le contraste et la résolution, mais reste faible (typiquement 15--25\% du diamètre $D$).
\end{itemize}

\begin{attention}[Erreur fréquente : Confondre les deux miroirs]
\begin{itemize}
    \item Le \textbf{miroir primaire} est \textbf{concave} (souvent parabolique). C'est \textbf{lui seul} qui forme l'image réelle de l'astre. C'est l'équivalent de l'objectif d'une lunette. Sa courbure détermine la focale $F$ et le grossissement.
    \item Le \textbf{miroir secondaire} est \textbf{plan}. Il ne fait que \textbf{renvoyer} l'image latéralement. Il n'a \textbf{aucune puissance optique} (ne grossit pas, ne change pas la focale effective).
\end{itemize}
Au Probatoire, on ne vous demandera \textbf{jamais} de tracer la construction d'image par le miroir secondaire (miroir plan trivial), mais bien par le miroir primaire (miroir concave : rayons parallèles $\to$ foyer).
\end{attention}

\paragraph{Rôle de l'oculaire.}
C'est une lentille mince convergente (ou un système de lentilles corrigé), placée de sorte que son foyer objet $F_o$ coïncide avec le plan de l'image intermédiaire (formée par le primaire + secondaire). L'oculaire fonctionne en \textbf{vision normale} (image finale à l'infini, œil au repos) ou en \textbf{vision de près} (image finale à la distance de vision distincte $D_0 = \SI{25}{cm}$, œil accommodé). Il fournit le \textbf{grossissement angulaire}.

\begin{popfigure}
\begin{tikzpicture}[scale=1.1]
% Trois composants côte à côte
% 1. Miroir Primaire
\begin{scope}[xshift=-5cm]
\draw[popblue, line width=3pt] (-1, -0.5) .. controls (0, 0) .. (1, -0.5);
\draw[popblue, line width=1pt, fill=popblue!10] (-1.2, -0.5) rectangle (1.2, -0.7);
\draw[->, poppink, thick] (-1.5, 1) -- (-1.5, 0.2);
\draw[->, poppink, thick] (0, 1) -- (0, 0.2);
\draw[->, poppink, thick] (1.5, 1) -- (1.5, 0.2);
\draw[poppink, thick] (-1.5, 0.2) .. controls (0, -0.5) .. (0, -0.5);
\draw[poppink, thick] (0, 0.2) .. controls (0, -0.5) .. (0, -0.5);
\draw[poppink, thick] (1.5, 0.2) .. controls (0, -0.5) .. (0, -0.5);
\node[popdark, align=center] at (0, -1.3) {\textbf{Miroir Primaire}\\Concave (Parabolique)\\Diamètre $D$, Focale $F$\\Forme l'image réelle};
\end{scope}

% 2. Miroir Secondaire
\begin{scope}[xshift=0cm]
\draw[poporange, line width=2pt, fill=poporange!10] (-0.7, -0.7) -- (0.7, 0.7) -- (0.9, 0.5) -- (-0.5, -0.9) -- cycle;
\draw[->, poppink, thick] (-1, 0) -- (-0.2, 0);
\draw[poppink, thick] (-0.2, 0) -- (0.2, 0.4); % réfléchi 90
\draw[->, poppink, thick] (0.2, 0.4) -- (0.2, 1.2);
\node[popdark, align=center] at (0, -1.3) {\textbf{Miroir Secondaire}\\Plan (45°)\\Pas de puissance optique\\Renvoie l'image à 90°};
\end{scope}

% 3. Oculaire
\begin{scope}[xshift=5cm]
\draw[poppurple, line width=2pt] (-0.5, -1) .. controls (-0.8, 0) .. (-0.5, 1);
\draw[poppurple, line width=2pt] (0.5, -1) .. controls (0.8, 0) .. (0.5, 1);
\draw[poppurple, line width=1pt] (-0.5, -1) -- (0.5, -1);
\draw[poppurple, line width=1pt] (-0.5, 1) -- (0.5, 1);
\draw[->, poppink, thick] (-1.5, 0) -- (-0.5, 0);
\draw[->, poppink, thick] (0.5, 0) -- (1.5, 0);
\node[poppurple] at (0, 1.3) {$F_o$};
\node[popdark, align=center] at (0, -1.3) {\textbf{Oculaire}\\Lentille convergente\\Focale $f'_o$\\Loupe (grossissement angulaire)};
\end{scope}
\end{tikzpicture}
\legende{Les trois éléments clés du télescope de Newton : le collecteur (primaire), le renvoi (secondaire), le grossisseur (oculaire).}
\end{popfigure}

\courssub{Grandeurs caractéristiques : Ouverture $D$, Focale $F$, Rapport $F/D$}

Trois paramètres numériques définissent les performances de base d'un télescope de Newton (et de tout télescope) :

\begin{enumerate}
    \item \textbf{Le diamètre $D$ (Ouverture).} C'est le diamètre du miroir primaire (utile, non obstrué). Il détermine :
    \begin{itemize}
        \item La \textbf{puissance collectrice de lumière} $\propto D^2$ (magnitude limite visible).
        \item Le \textbf{pouvoir séparateur théorique} (diffraction) : $\theta_{\text{min}} \approx 1,22 \frac{\lambda}{D}$ (critère de Rayleigh). Plus $D$ est grand, plus les détails fins sont résolus.
    \end{itemize}
    \item \textbf{La distance focale $F$ du miroir primaire.} C'est la distance sommet-foyer $PF'$. Elle détermine :
    \begin{itemize}
        \item La \textbf{taille de l'image intermédiaire} : pour un astre de diamètre apparent $\alpha$ (en radians), la taille linéaire $h'$ de son image au foyer est $h' = F \cdot \alpha$ (relation de conjugaison pour objet à l'infini).
        \item La \textbf{longueur du tube} (environ $F$).
    \end{itemize}
    \item \textbf{Le rapport d'ouverture (nombre $f$) : $N = \frac{F}{D}$.}
    \begin{itemize}
        \item C'est un nombre sans dimension (souvent noté $f/N$, ex: $f/8$).
        \item Il caractérise la \textbf{luminosité de l'image} au foyer (illuminance $\propto 1/N^2$). Un $F/D$ faible (ex: $f/4$) donne une image brillante, champ large, idéal pour la photo de nébuleuses (astrophotographie). Un $F/D$ élevé (ex: $f/10$) donne une image plus sombre mais un grossissement plus facile à obtenir avec des oculaires standards, idéal pour le planétaire.
        \item Il conditionne la \textbf{profondeur de champ} et la tolérance de mise au point.
    \end{itemize}
\end{enumerate}

\begin{exemplebox}[Télescope Newton amateur classique (type « 114/900 »)]
Un instrument très répandu chez les amateurs camerounais (souvent marque SkyWatcher, Celestron, ou fabrication artisanale) porte la désignation \textbf{114/900}.
\begin{itemize}
    \item $D = \SI{114}{mm} = \SI{0,114}{m}$ (Diamètre du miroir primaire).
    \item $F = \SI{900}{mm} = \SI{0,90}{m}$ (Distance focale du miroir primaire).
    \item Rapport $F/D = \frac{900}{114} \approx 7,9 \approx 8$. C'est un \textbf{$f/8$}.
\end{itemize}
\textbf{Conséquences pratiques :}
\begin{itemize}
    \item \textbf{Pouvoir séparateur} (lumière verte $\lambda = \SI{550}{nm}$) : $\theta_{\text{min}} \approx 1,22 \times \frac{550 \times 10^{-9}}{0,114} \approx \SI{5,9e-6}{rad} \approx \SI{1,2}{''}$. Suffisant pour séparer les bandes de Jupiter ou les anneaux de Saturne.
    \item \textbf{Magnitude limite} (estimation) : $m_{\text{lim}} \approx 2,7 + 5 \log_{10}(D/\text{mm}) \approx 2,7 + 5 \times 2,06 \approx 13$. On voit des objets 1000 fois plus faibles qu'à l'œil nu ($m \approx 6$).
    \item \textbf{Encombrement :} Tube $\approx \SI{90}{cm}$, transportable dans un taxi ou sur une moto-taxi (\"bensikin\") pour sortir en brousse loin de la pollution lumineuse de Yaoundé.
\end{itemize}
\end{exemplebox}

\begin{methode}[Identifier les caractéristiques d'un télescope de Newton]
\begin{enumerate}
    \item Repérer le \textbf{miroir primaire} au fond du tube : c'est le grand miroir concave (souvent parabolique). Mesurer ou lire son diamètre $D$ et sa focale $F$.
    \item Repérer le \textbf{miroir secondaire} : petit miroir plat elliptique à 45° près de l'ouverture. Vérifier qu'il ne vignette pas excessivement (obstruction $< 30\% D$).
    \item Repérer le \textbf{porte-oculaire} (focuser) sur le côté du tube : c'est là qu'on insère l'oculaire. La distance miroir Primaire $\to$ plan focal $\approx F$.
    \item Calculer le rapport $F/D$ pour connaître la « vitesse » de l'instrument (lumineux si $F/D < 6$, polyvalent si $6 < F/D < 10$, planétaire si $F/D > 10$).
\end{enumerate}
\end{methode}

\begin{exoresolu}[Analyse d'un télescope de Newton « 200/1000 »]
\textbf{Énoncé :} Un club d'astronomie de l'Université de Ngaoundéré acquiert un télescope de Newton de désignation « 200/1000 ». Le miroir secondaire a un axe mineur de \SI{50}{mm}.
\begin{enumerate}
    \item Donner $D$, $F$ et le rapport $F/D$.
    \item Calculer le pouvoir séparateur théorique $\theta_{\text{min}}$ (en secondes d'arc) pour $\lambda = \SI{550}{nm}$.
    \item Calculer le pourcentage d'obstruction linéaire et surfacique due au miroir secondaire.
    \item Quel type d'observation cet instrument favorise-t-il (planétaire ou ciel profond) ? Justifier.
\end{enumerate}
\textbf{Solution détaillée :}
\begin{enumerate}
    \item $D = \SI{200}{mm}$, $F = \SI{1000}{mm}$. Rapport $F/D = \frac{1000}{200} = 5$. C'est un \textbf{$f/5$} (instrument « rapide », lumineux).
    \item $\theta_{\text{min}} = 1,22 \frac{\lambda}{D} = 1,22 \times \frac{550 \times 10^{-9}}{0,200} = \SI{3,355e-6}{rad}$.
    Conversion en secondes d'angle : $1 \text{ rad} = \frac{180 \times 3600}{\pi} \approx 206265''$.
    $\theta_{\text{min}} \approx 3,355 \times 10^{-6} \times 206265 \approx \textbf{\SI{0,69}{''}}$. Excellent pouvoir séparateur.
    \item Obstruction linéaire : $\frac{d_{\text{secondaire}}}{D} = \frac{50}{200} = 0,25 = \textbf{25\%}$.
    Obstruction surfacique (perte de lumière) : $\left(\frac{50}{200}\right)^2 = 0,0625 = \textbf{6,25\%}$. Acceptable (généralement $< 20\%$ surfacique).
    \item $F/D = 5$ est faible. L'image au foyer est \textbf{très lumineuse} (illuminance $\propto 1/25$ vs $1/64$ pour un $f/8$). Champ large possible. \textbf{Favorise le ciel profond (nébuleuses, galaxies, amas)} et l'astrophotographie. Pour le planétaire, il faudra des oculaires de très courte focale (ex: $\SI{4}{mm}$) ou une lentille de Barlow $\times 2$ ou $\times 3$ pour atteindre les grossissements utiles ($G_{\text{max}} \approx 2D(\text{mm}) = 400\times$).
\end{enumerate}
\end{exoresolu}

\begin{propriete}[Résumé des grandeurs caractéristiques du télescope de Newton]
Pour un télescope de Newton de miroir primaire de diamètre $D$ et de distance focale $F$ :
\begin{align*}
    \text{Rapport d'ouverture } & N = \frac{F}{D} \quad (\text{sans dimension, noté } f/N) \\
    \text{Pouvoir séparateur (Rayleigh) } & \theta_{\text{min}} \approx 1,22 \frac{\lambda}{D} \quad (\text{en radians}) \\
    \text{Taille image intermédiaire } & h' = F \cdot \alpha \quad (\alpha \text{ en radians}) \\
    \text{Grossissement max utile } & G_{\text{max}} \approx 2 \times D(\text{mm}) \quad (\text{empirique, turbulence limite}) \\
    \text{Magnitude limite visuelle } & m_{\text{lim}} \approx 2,7 + 5 \log_{10}(D/\text{mm})
\end{align*}
\textbf{Analyse dimensionnelle :} $[F] = L$, $[D] = L \Rightarrow [F/D] = 1$ (correct). $[\lambda/D] = L/L = 1$ (radians, correct).
\end{propriete}

\courssub{Principe de fonctionnement et formation des images}

Le télescope de Newton est un système \textbf{afocal} : l'objet est à l'infini (astre) et l'image finale est également à l'infini (vision normale, œil au repos). Le fonctionnement se décompose en deux étapes successives :

\begin{enumerate}
    \item \textbf{Objectif (Miroir primaire) :} L'astre, source à l'infini, envoie un faisceau de rayons parallèles. Le miroir primaire concave (focale $F$) forme une \textbf{image réelle, inversée et réduite} $A'B'$ dans son plan focal (plan focal objet de l'oculaire). La taille de cette image est $h' = F \cdot \alpha$, où $\alpha$ est le diamètre apparent de l'astre (angle sous-tendu à l'entrée).
    \item \textbf{Loupe (Oculaire) :} L'oculaire (focale $f'_o$) est placé de sorte que son foyer objet $F_o$ coïncide avec le plan de l'image intermédiaire $A'B'$. L'image intermédiaire joue le rôle d'objet pour l'oculaire. Comme cet objet est au foyer de l'oculaire, les rayons émergents sont parallèles : l'image finale $A''B''$ est à l'infini. L'œil placé derrière l'oculaire reçoit un faisceau de rayons parallèles sous un angle $\alpha'$ plus grand que $\alpha$.
\end{enumerate}

\begin{experience}[Tracer la marche d'un faisceau dans un télescope de Newton (Savoir-faire Probatoire)]
\textbf{Matériel :} Règle, rapporteur, feuille millimétrée, schéma du télescope (miroir primaire, secondaire, oculaire).
\textbf{Protocole :}
\begin{enumerate}
    \item Tracer l'axe optique principal (horizontal) et le miroir primaire (concave, sommet $P$, foyer $F'$ à distance $F$).
    \item Tracer le miroir secondaire (plan, à 45°, positionné avant $F'$).
    \item Tracer l'oculaire (lentille mince, centre $O_o$, foyers $F_o$ et $F'_o$ à distance $f'_o$) de sorte que $F_o$ soit sur l'image intermédiaire.
    \item \textbf{Faisceau incident :} Dessiner deux rayons parallèles à l'axe principal (venant d'un point de l'astre sur l'axe) et un rayon parallèle décalé (venant d'un point hors axe, angle $\alpha$).
    \item \textbf{Réflexion sur le primaire :} Les rayons parallèles à l'axe passent par $F'$. Le rayon hors axe (parallèle à la direction de l'astre) converge vers un point $B'$ dans le plan focal tel que $PB' = F \tan \alpha \approx F \alpha$ (petits angles).
    \item \textbf{Réflexion sur le secondaire :} Appliquer la loi de Descartes ($i=r$) sur le miroir plan à 45°. Le faisceau convergent est dévié de 90° vers l'oculaire. Le point $B'$ est symétrique par rapport au plan du miroir secondaire (mais on trace simplement la déviation des rayons).
    \item \textbf{Passage dans l'oculaire :} Les rayons incidents sur l'oculaire convergent vers $B'$ (qui est au foyer $F_o$ de l'oculaire).
    \begin{itemize}
        \item Rayon passant par le centre optique $O_o$ : non dévié.
        \item Rayon passant par $F_o$ (donc passant par $B'$) : ressort parallèle à l'axe de l'oculaire.
        \item Rayon parallèle à l'axe de l'oculaire : ressort en passant par $F'_o$.
    \end{itemize}
    \item \textbf{Résultat :} Les rayons émergents sont parallèles entre eux. Ils forment un angle $\alpha'$ avec l'axe de l'oculaire. L'image finale est à l'infini, inversée par rapport à l'astre original.
\end{enumerate}
\textbf{Observation :} Le tracé montre que le télescope transforme un faisceau entrant d'angle $\alpha$ en un faisceau sortant d'angle $\alpha'$.
\end{experience}

\begin{popfigure}
\begin{tikzpicture}[scale=0.8, >=Stealth]
% Axe principal horizontal (entrée)
\draw[->, thin, popmuted] (-4,0) -- (4,0) node[right] {Axe entrée};
% Miroir Primaire (concave, sommet en (0,0) pour simplifier, foyer F' à gauche? Non, Newton: lumière arrive par la gauche, miroir au fond à droite. Reprenons.)
% Standard: Lumière arrive de la gauche. Miroir primaire concave sommet P à droite. Foyer F' à gauche de P.
% Miroir secondaire avant F'. Oculaire sur le côté (haut).
% Redessinons proprement.

% Définition des points
\coordinate (P) at (0,0); % Sommet miroir primaire
\coordinate (Fp) at (-3,0); % Foyer primaire F' (F=3)
\coordinate (S) at (-1.5,0); % Intersection axe/miroir secondaire (milieu approx)
% Miroir secondaire : plan à 45°, passe par S. Normale à 45°. Equation y = x + 1.5 (si S=(-1.5,0))
% Oculaire : centre Oo sur l'axe dévié (vertical). Foyer objet Fo coïncide avec image intermédiaire.
% Image intermédiaire B' sur axe principal à x = -2.5 (entre S et Fp).
\coordinate (Bprime) at (-2.5, 0);
% Miroir secondaire doit être avant B' ? Non, secondaire renvoie l'image. L'image se forme DANS le plan focal.
% Le secondaire est placé AVANT le foyer. Le faisceau converge VERS B'. Il rencontre le secondaire AVANT B'.
% Après réflexion sur secondaire, le faisceau converge vers B'' (image de B' par miroir plan).
% L'oculaire est placé pour que son foyer objet soit en B''.

% Simplification pédagogique standard : On trace le faisceau "plié" comme s'il était droit.
% Axe optique plié : Horizontal vers la gauche (entrée), puis vertical vers le haut (sortie).

% --- DESSIN AVEC AXE PLIÉ ---
% Segment 1 : Entrée horizontale (x < -1.5)
\draw[thick, popblue] (-4, 1.5) -- (-1.5, 1.5); % Rayon haut
\draw[thick, popblue] (-4, 0) -- (-1.5, 0);     % Rayon axe
\draw[thick, popblue] (-4, -1.5) -- (-1.5, -1.5); % Rayon bas

% Miroir Secondaire à 45° en (-1.5, 0) à (-0.5, 1) approx? Non, centré sur l'axe.
% Miroir secondaire : segment de (-1.5, -0.5) à (-0.5, 0.5) ? Non, il coupe l'axe.
% Disons miroir secondaire : segment de (-2, -0.5) à (-1, 0.5). Centre en (-1.5, 0).
\draw[poporange, line width=2.5pt, fill=poporange!20] (-2, -0.5) -- (-1, 0.5) -- (-0.9, 0.6) -- (-1.9, -0.4) -- cycle;
\node[poporange, anchor=south] at (-1, 0.6) {\textbf{Secondaire (45°)}};

% Réflexion : Rayon axe (y=0) arrive en (-1.5, 0). Normal au miroir = 45°. Incidence 45°. Réflexion verticale vers le haut.
% Rayon haut (y=1.5) arrive en (-1.5, 1.5) -> hors miroir. Il faut que les rayons convergent.
% On trace le faisceau CONVERGENT vers le foyer F' (-3, 0).
% Le secondaire est à x=-1.5. Le foyer est à x=-3. Le secondaire est AVANT le foyer.
% Les rayons viennent de la gauche, convergent vers F'(-3,0).
% Ils rencontrent le secondaire en x=-1.5.

% Rayon 1 (axe) : passe par (-1.5, 0) et (-3, 0). Réfléchi verticalement vers le haut par le miroir 45°.
% Rayon 2 (marge haute) : passe par (-1.5, y1) et (-3, 0). Pente = (0-y1)/(-3+1.5) = -y1/-1.5 = 2y1/3.
%   En x=-1.5, y=y1. Réflexion : incidence relative à la normale du miroir.
% C'est complexe à tracer "à la main" en TikZ sans calculs lourds.
% APPROCHE PÉDAGOGIQUE : On trace le schéma "déplié" (standard dans les manuels).

% SCHÉMA DÉPLIÉ (Standard Probatoire)
% Axe optique unique horizontal. Miroir primaire à droite. Oculaire à gauche. Image intermédiaire au milieu.
% On note "Miroir secondaire (plan, 45°) : renvoi latéral non représenté sur l'axe plié".

\end{tikzpicture}
\legende{Schéma de principe déplié (équivalent optique) : l'axe est redressé. Le miroir secondaire (plan) n'a pas de puissance, il se "transparentise" pour le calcul de parcours. L'image intermédiaire $A'B'$ se forme au foyer $F'$ du primaire. L'oculaire est placé à sa focale $f'_o$ de cette image.}
\end{popfigure}

% Remplaçons la figure TikZ complexe par un schéma déplié propre et standard pour le cours.
% Je vais réécrire la figure ci-dessus en TikZ "déplié" propre.

\courssub{Grossissement et puissance}

\paragraph{Définition du grossissement $G$.}
Le \textbf{grossissement (ou pouvoir grossissant)} $G$ d'un instrument afocal est le rapport de l'angle sous-tendu par l'image finale à l'œil ($\alpha'$) sur l'angle sous-tendu par l'objet à l'entrée de l'instrument ($\alpha$) :
\[
G = \frac{\alpha'}{\alpha}
\]
Par convention, si l'image finale est inversée par rapport à l'objet (ce qui est le cas ici : miroir primaire inverse haut/bas, oculaire ne redresse pas), $G$ est \textbf{négatif}. On retient souvent la valeur absolue $|G|$.

\paragraph{Expression du grossissement en fonction des focales.}
Dans la configuration \textbf{vision normale} (image finale à l'infini, œil au repos) :
\begin{itemize}
    \item L'image intermédiaire $A'B'$ a une taille $h' = F \cdot \alpha$ (petits angles, $\alpha$ en radians).
    \item Cette image se trouve au foyer objet $F_o$ de l'oculaire (focale $f'_o$).
    \item L'oculaire fonctionne comme une loupe : l'angle de sortie $\alpha'$ vérifie $h' = f'_o \cdot \alpha'$ (petits angles).
\end{itemize}
En égalant les deux expressions de $h'$ :
\[
F \cdot \alpha = f'_o \cdot \alpha' \quad \Rightarrow \quad G = \frac{\alpha'}{\alpha} = -\frac{F}{f'_o}
\]
Le signe $-$ traduit l'inversion de l'image (miroir primaire inverse, oculaire conserve l'orientation de son objet qui est déjà inversé).

\begin{propriete}[Grossissement du télescope de Newton (Vision normale)]
\[
\boxed{G = -\frac{F}{f'_o}}
\]
\begin{itemize}
    \item $F$ : distance focale du miroir primaire (en \si{\meter} ou \si{\milli\meter}).
    \item $f'_o$ : distance focale de l'oculaire (même unité que $F$).
    \item $G$ est sans dimension. Le signe $-$ indique une image finale \textbf{inversée} (tête en bas) par rapport à l'astre.
\end{itemize}
\textbf{Analyse dimensionnelle :} $[F] = L$, $[f'_o] = L \Rightarrow [G] = 1$ (correct).
\end{propriete}

\begin{attention}[Grossissement et puissance : ne pas confondre !]
\begin{itemize}
    \item \textbf{Grossissement $G$} : rapport d'angles (sans dimension). Dépend du couple (télescope, oculaire). On change $G$ en changeant d'oculaire.
    \item \textbf{Puissance (vergences) :} La \textbf{puissance de l'objectif} (miroir) est $V_{\text{obj}} = \frac{1}{F}$ (en dioptries, \si{\dioptre}). La \textbf{puissance de l'oculaire} est $V_{\text{oc}} = \frac{1}{f'_o}$ (en \si{\dioptre}). On a $G = -\frac{V_{\text{oc}}}{V_{\text{obj}}}$.
    \item Au Probatoire, on demande le \textbf{grossissement $G$}. La "puissance" du télescope désigne souvent son pouvoir grossissant maximal utile ou sa puissance collectrice ($D^2$), pas sa vergence.
\end{itemize}
\end{attention}

\paragraph{Grossissement en vision de près (œil accommodé).}
Si l'observateur accommode pour voir l'image finale à la distance de vision distincte $D_0 = \SI{0,25}{m}$, l'oculaire doit former une image virtuelle à $-D_0$. La relation de conjugaison pour l'oculaire donne :
\[
\frac{1}{f'_o} = \frac{1}{p'} - \frac{1}{p} \quad \text{avec } p' = -D_0 \text{ et } p = f'_o \text{ (objet au foyer) ? Non.}
\]
En vision de près, l'image intermédiaire n'est plus exactement au foyer $F_o$ de l'oculaire, mais légèrement décalée pour que l'image finale soit à $-D_0$.
Le grossissement devient :
\[
G_{\text{près}} = -\frac{F}{f'_o} \left(1 + \frac{f'_o}{D_0}\right) = G_{\text{normale}} \left(1 + \frac{f'_o}{D_0}\right)
\]
Comme $f'_o \ll D_0$ généralement (ex: $f'_o = \SI{10}{mm}$, $D_0 = \SI{250}{mm}$), le facteur de correction est proche de 1 (ex: $1,04$). \textbf{Au Probatoire, on utilise presque toujours la formule de vision normale $G = -F/f'_o$.}

\begin{exemplebox}[Calcul de grossissements pour un télescope 114/900 ($F = \SI{900}{mm}$)]
\begin{center}
\begin{tabularx}{\linewidth}{|c|X|c|}\hline
\textbf{Oculaire ($f'_o$)} & \textbf{Calcul} & \textbf{Grossissement $|G|$} \\ \hline
\SI{25}{mm} (faible grossissement, champ large) & $900 / 25$ & \textbf{36$\times$} \\ \hline
\SI{10}{mm} (moyen) & $900 / 10$ & \textbf{90$\times$} \\ \hline
\SI{6}{mm} (fort, planétaire) & $900 / 6$ & \textbf{150$\times$} \\ \hline
\SI{4}{mm} + Barlow $\times 2$ ($f'_{\text{eq}} = \SI{2}{mm}$) & $900 / 2$ & \textbf{450$\times$} (souvent au-delà du max utile $2D=228\times$) \\ \hline
\end{tabularx}
\end{center}
\textbf{Règle pratique :} $G_{\text{max utile}} \approx 2 \times D(\text{mm})$. Pour $D=114$, $G_{\text{max}} \approx 228\times$. Au-delà, l'image est "molle" (turbulence, diffraction) sans gain de détail.
\end{exemplebox}

\courssub{Mise au point du télescope de Newton}

La \textbf{mise au point} (ou \textbf{focalisation}) consiste à rendre l'image nette pour l'observateur. Elle se fait en déplaçant l'oculaire le long de l'axe optique (via le \emph{focuser} à crémaillère ou à crayford) pour que le plan de l'image intermédiaire coïncide exactement avec le plan focal objet de l'oculaire.

\begin{methode}[Procédure de mise au point (Sur le terrain)]
\begin{enumerate}
    \item \textbf{Visée grossière :} Pointer le tube vers un objet lointain lumineux (réverbère, sommet d'arbre, planète brillante) à l'aide du chercheur (petite lunette de visée à faible grossissement, champ large, montée en parallèle sur le tube).
    \item \textbf{Insertion de l'oculaire :} Commencer par l'oculaire de \textbf{plus grande focale} (ex: \SI{25}{mm} ou \SI{30}{mm}) : grossissement faible, champ large, tolérance de mise au point grande, image plus lumineuse.
    \item \textbf{Réglage grossier :} Tourner la molette du focuser pour faire sortir/rentrer le porte-oculaire jusqu'à voir une image nette (pour un objet lointain, l'image intermédiaire est au foyer $F'$, l'oculaire doit être à $f'_o$ de ce plan).
    \item \textbf{Affinement :} Tourner lentement dans un sens puis l'autre pour trouver le "point dur" de netteté maximale.
    \item \textbf{Changement de grossissement :} Pour augmenter le grossissement, remplacer par un oculaire de focale plus courte. \textbf{Refaire la mise au point} à chaque changement (la position du plan focal objet de l'oculaire varie légèrement selon les modèles, et l'accommodation de l'œil change).
    \item \textbf{Adaptation à la vue :} Si vous portez des lunettes (astigmatisme), gardez-les. Si vous êtes myope/hypermétrope sans astigmatisme, vous pouvez les ôter et compenser par la mise au point (l'instrument devient vos lunettes).
\end{enumerate}
\end{methode}

\begin{attention}[Pièges de la mise au point]
\begin{itemize}
    \item \textbf{Image "donut" (anneau) :} Si vous voyez un anneau brillant avec un trou noir au centre (l'ombre du miroir secondaire), vous êtes \textbf{très loin de la mise au point} (intra-focal ou extra-focal). Rapprochez-vous du point où l'anneau se referme en un disque (étoile) ou une image nette (planète).
    \item \textbf{Sens de rotation :} Notez le sens pour "rentrer" (rapprocher l'oculaire du miroir) et "sortir". En général, pour passer d'un oculaire long à un oculaire court, il faut \textbf{rentrer} le porte-oculaire (l'image intermédiaire est fixe, l'oculaire plus "puissant" a son foyer plus près de sa lentille).
    \item \textbf{Back focus (arrière-foyer) :} La distance entre le miroir secondaire et le plan focal doit être suffisante pour loger le porte-oculaire et l'oculaire. Vérifier à l'achat (problème fréquent sur les Newton "photo" avec capteurs épais).
\end{itemize}
\end{attention}

\courssub{Intérêt du télescope par rapport aux instruments à lentilles (Lunettes astronomiques)}

Le télescope de Newton (réflecteur) et la lunette astronomique (réfracteur) remplissent la même fonction (observation d'objets à l'infini) mais par des principes physiques différents.

\begin{center}
\begin{tabularx}{\linewidth}{|l|X|X|}\hline
\rowcolor{popblueL}
\textbf{Critère} & \textbf{Télescope de Newton (Réflecteur)} & \textbf{Lunette Astronomique (Réfracteur)} \\ \hline
\textbf{Principe} & Réflexion sur miroir concave (parabolique) & Réfraction à travers lentille convergente (objectif) \\ \hline
\textbf{Aberration chromatique} & \textbf{Absente} (loi de Descartes $i=r$, indépendant de $\lambda$). Avantage majeur. & \textbf{Présente} sur objectif simple (achromat : résidu "violet", apochromat : corrigé mais très cher). \\ \hline
\textbf{Coût / Ouverture $D$} & \textbf{Très favorable.} Miroir unique à polir (une face). $D=200\text{--}300\text{ mm}$ accessible aux amateurs. & \textbf{Défavorable.} Objectif = 2 à 4 lentilles (verres coûteux, 4 à 8 surfaces à polir). $D > \SI{100}{mm}$ devient très onéreux. \\ \hline
\textbf{Encombrement (Longueur)} & Tube $\approx F$. Compact pour un $F/D$ donné. & Tube $\approx F_{\text{obj}}$. Souvent plus long pour même $F$. \\ \hline
\textbf{Obstruction centrale} & \textbf{Oui.} Miroir secondaire + support (araignée). Perte de lumière ($\sim 5\text{--}10\%$ surfacique), pics de diffraction (croix sur étoiles brillantes), baisse légère du contraste. & \textbf{Non.} Pupille d'entrée pleine. Meilleur contraste planétaire, pas de pics. \\ \hline
\textbf{Aberrations géométriques} & \textbf{Coma} hors axe (étoiles en forme de comète aux bords du champ). Nécessite correcteur (coma corrector) pour photo. Courbure de champ. & Bien corrigées sur apochromats (champ plat). Coma quasi nul. \\ \hline
\textbf{Maintenance} & \textbf{Collimation} régulière nécessaire (alignement miroirs primaire/secondaire). Sensible aux chocs. & \textbf{Quasi nulle.} Objectif scellé, alignement d'usine définitif. "Pointez et observez". \\ \hline
\textbf{Thermique} & Miroir épais $\to$ temps d'équilibrage long (courants d'air internes dégradent l'image). Tube ouvert aide. & Objectif plus fin, tube fermé $\to$ équilibrage plus rapide mais risque de buée sur lentille avant. \\ \hline
\textbf{Usage dominant} & \textbf{Ciel profond} (nébuleuses, galaxies) grâce au grand $D$ accessible. Planétaire possible (collimation soignée). & \textbf{Planétaire, Lune, doubles étoiles} (contraste max, netteté). Astrophotographie grand champ (apochromats courts). \\ \hline
\end{tabularx}
\end{center}

\begin{aretenir}[Pourquoi Newton domine l'astronomie amateur "grand diamètre" ?]
Pour un budget donné (ex: \SI{300\,000}{FCFA} -- \SI{500\,000}{FCFA} au Cameroun), un Newton offre $D = \SI{200}{mm}$ à $\SI{250}{mm}$ ($f/5$ à $f/6$), tandis qu'une lunette achromate se limite à $D = \SI{80}{mm}$--\SI{100}{mm}$ et une apochromate à $D = \SI{60}{mm}$--\SI{80}{mm}$. \textbf{En astronomie, l'ouverture $D$ est roi} (résolution $\propto 1/D$, lumière $\propto D^2$). Le Newton est le seul design permettant d'atteindre de grandes ouvertures à coût raisonnable.
\end{aretenir}

\begin{savaistu}[Les grands télescopes professionnels : tous des Newton (ou dérivés) !]
Tous les grands télescopes modernes (VLT \SI{8,2}{m}, Keck \SI{10}{m}, ELT \SI{39}{m}, JWST \SI{6,5}{m}) sont des \textbf{réflecteurs} (souvent Ritchey-Chrétien, variante hyperbolique du Cassegrain, mais principe de Newton : miroir primaire collecteur). Pourquoi ?
\begin{itemize}
    \item Un miroir ne se tient que par l'arrière (support cellulaire) $\to$ pas de déformation par son poids (contrairement à une lentille tenue par le bord qui s'affaisse).
    \item Un miroir n'a qu'une surface optique à polir (vs 4 pour un doublet achromat).
    \item Pas d'absorption chromatique dans le verre (important pour l'IR/UV).
    \item Le verre peut être de qualité moindre (seule la surface compte), ou en matériaux légers (Zerodur, SiC, Béryllium pour JWST).
\end{itemize}
Le télescope de Newton que vous tenez en main est le petit frère des géants qui sondent l'Univers !
\end{savaistu}

\begin{exoresolu}[Choix d'instrument pour un club d'astronomie à Maroua]
\textbf{Énoncé :} Un club d'astronomie du Lycée de Maroua (ciel très noir, peu de pollution lumineuse, climat sec) dispose d'un budget de \SI{400\,000}{FCFA}. Il hésite entre :
\begin{itemize}
    \item Instrument A : Lunette achromate \SI{102}{mm} / \SI{1000}{mm} ($f/9,8$) sur monture équatoriale motorisée.
    \item Instrument B : Télescope Newton \SI{200}{mm} / \SI{1000}{mm} ($f/5$) sur monture Dobson (azimutale manuelle).
\end{itemize}
Comparer les deux choix pour : (1) Observation planétaire (Jupiter, Saturne). (2) Observation ciel profond (nébuleuses, galaxies). (3) Astrophotographie débutant. (4) Facilité d'usage pour des lycéens.
\textbf{Solution détaillée :}
\begin{enumerate}
    \item \textbf{Planétaire :} Lunette A : Pas d'obstruction $\to$ contraste excellent, image "propre". Mais $D=102$ mm limite la résolution ($\theta_{\text{min}} \approx 1,3''$) et la luminosité. Newton B : Obstruction 25\% $\to$ contraste légèrement inférieur, pics de diffraction. Mais $D=200$ mm $\to$ résolution double ($\theta_{\text{min}} \approx 0,7''$), détails plus fins visibles \textbf{si} la collimation est parfaite et la turbulence le permet. \textbf{Verdict :} Newton B potentiellement supérieur *si* bien collimaté ; Lunette A plus "tolérante" et stable.
    \item \textbf{Ciel profond :} Newton B sans conteste. Surface collectrice $(200/102)^2 \approx 3,8$ fois plus grande. $f/5$ vs $f/10$ $\to$ champ 2 fois plus large pour même oculaire, temps de pose 4 fois plus courts en photo. \textbf{Verdict :} Newton B largement gagnant.
    \item \textbf{Astrophoto :} Newton B : $f/5$ rapide, mais coma aux bords (nécessite correcteur $\approx$ \SI{50\,000}{FCFA}), mise au point critique, monture Dobson \textbf{non motorisée} = impossible de poser > 1s (rotation de champ). Lunette A : $f/10$ lent, mais monture équatoriale motorisée = poses longues possibles, champ plat natif (presque). \textbf{Verdict :} Lunette A + monture EQ \textbf{bien meilleure pour débuter en photo}. Newton B sur Dobson = observation visuelle uniquement.
    \item \textbf{Facilité :} Lunette A : "Pointez-observez", pas de collimation. Monture EQ demande alignement polaire (apprentissage). Newton B : Dobson = intuitif (haut/bas, gauche/droite), mais \textbf{collimation obligatoire} à chaque transport (apprentissage technique). \textbf{Verdict :} Lunette A plus "scolaire" clé en main ; Newton B plus "école d'ingénieur" (apprentissage optique/mécanique).
\end{enumerate}
\textbf{Recommandation pour Maroua (ciel noir) :} \textbf{Newton B (200/1000 Dobson).} Le gain d'ouverture ($D=200$) est décisif pour exploiter la qualité du ciel sur le ciel profond. La collimation s'apprend en 15 min avec un collimateur laser (\SI{10\,000}{FCFA}). Pour la photo planétaire "lucky imaging" (vidéo courte), la monture Dobson suffit (dérive gérée par logiciel).
\end{exoresolu}

\begin{exercice}[Entraînement Probatoire : Télescope de Newton]
\begin{enumerate}
    \item \textbf{Description :} Légender un schéma de télescope de Newton (Miroir primaire, Miroir secondaire, Oculaire, Foyer principal, Ouverture $D$, Focale $F$). Indiquer le sens des rayons pour une étoile sur l'axe et une étoile hors axe (angle $\alpha$).
    \item \textbf{Calcul :} Un télescope Newton a un miroir primaire de diamètre $D = \SI{150}{mm}$ et de focale $F = \SI{750}{mm}$. On utilise un oculaire de focale $f'_o = \SI{12}{mm}$.
    \begin{enumerate}
        \item Calculer le rapport $F/D$ et qualifier l'instrument.
        \item Calculer le grossissement $G$ en vision normale.
        \item L'image de Jupiter (diamètre apparent $\alpha = \SI{45}{''}$) est formée au foyer. Calculer sa taille linéaire $h'$ (en mm).
        \item Quel grossissement obtient-on avec une lentille de Barlow $\times 2$ insérée avant l'oculaire ? (Rappel : la Barlow double la focale effective de l'objectif).
    \end{enumerate}
    \item \textbf{Mise au point :} On observe Saturne avec l'oculaire de \SI{12}{mm}. L'image est nette. On change pour un oculaire de \SI{6}{mm}. Faut-il rentrer ou sortir le porte-oculaire ? Justifier par la position du foyer de l'oculaire.
    \item \textbf{Argumentaire :} Un élève affirme : "Les lunettes sont meilleures que les télescopes car elles n'ont pas d'obstruction centrale". En vous basant sur la physique (résolution, lumière, aberrations, coût), nuancer cette affirmation pour un budget de \SI{200\,000}{FCFA}.
\end{enumerate}
\end{exercice}

\begin{corrige}[Exercice Entraînement Probatoire]
\begin{enumerate}
    \item \textbf{Schéma :} (À dessiner : Tube, Miroir primaire concave à droite, Foyer $F'$ marqué, Miroir secondaire plan à 45° avant $F'$, Oculaire sur le côté, Œil. Rayons parallèles entrants $\to$ convergence vers $F'$ $\to$ réflexion 90° $\to$ oculaire $\to$ parallèles sortants. Rayon axe et rayon hors axe $\alpha$).
    \item \textbf{Calculs :}
    \begin{enumerate}
        \item $F/D = 750 / 150 = 5$. Instrument **$f/5$** (rapide, lumineux, champ large).
        \item $G = -F/f'_o = -750 / 12 = \textbf{-62,5$\times$}$ (image inversée, grossissement 62,5 fois).
        \item $\alpha = 45'' = 45 \times \frac{\pi}{180 \times 3600} \approx \SI{2,18e-4}{rad}$.
        $h' = F \cdot \alpha = 750 \times 2,18 \times 10^{-4} \approx \textbf{\SI{0,164}{mm}}$ (soit \SI{164}{\micro\meter}). Image très petite !
        \item Barlow $\times 2$ : Focale effective de l'objectif $F_{\text{eff}} = 2 \times F = \SI{1500}{mm}$.
        $G_{\text{Barlow}} = -F_{\text{eff}} / f'_o = -1500 / 12 = \textbf{-125$\times$}$. (Équivalent à utiliser un oculaire de \SI{6}{mm} sans Barlow).
    \end{enumerate}
    \item \textbf{Mise au point :} Oculaire \SI{12}{mm} $\to$ foyer objet à \SI{12}{mm} de sa lentille. Oculaire \SI{6}{mm} $\to$ foyer objet à \SI{6}{mm} de sa lentille. L'image intermédiaire (au foyer du primaire) est fixe. Pour que le foyer de l'oculaire \SI{6}{mm} coïncide avec cette image fixe, il faut que la lentille de l'oculaire \SI{6}{mm} soit \textbf{plus proche} de l'image intermédiaire que la lentille de l'oculaire \SI{12}{mm}. Il faut donc \textbf{rentrer le porte-oculaire} (rapprocher l'oculaire du miroir secondaire/primaire).
    \item \textbf{Argumentaire :}
    \begin{itemize}
        \item \textbf{Résolution :} $\theta_{\text{min}} \propto 1/D$. Pour \SI{200\,000}{FCFA} : Lunette $\approx \SI{80}{mm}$ ($\theta \approx 1,7''$), Newton $\approx \SI{200}{mm}$ ($\theta \approx 0,7''$). Newton résout 2,5 fois plus de détails.
        \item \textbf{Lumière :} Flux $\propto D^2$. Newton collecte $(200/80)^2 \approx 6,25$ fois plus de photons. Magnitude limite $\approx 1,5$ mag de mieux.
        \item \textbf{Aberration chromatique :} Lunette achromate à ce prix : franges violettes sur Jupiter/Vénus (frange résiduelle). Newton : aucune.
        \item \textbf{Obstruction :} Newton $\sim 25\%$ linéaire. Baisse contraste $\sim 10\%$, pics diffraction. Mais le gain de résolution et lumière l'emporte largement sur objets faiblement lumineux (nébuleuses, galaxies) et même sur planètes si seeing moyen.
        \item \textbf{Conclusion :} Pour \SI{200\,000}{FCFA}, le Newton \SI{200}{mm} est \textbf{supérieur} pour 90\% des usages (ciel profond, planétaire si collimation soignée). La lunette \SI{80}{mm} n'est préférable que pour la facilité (pas de collimation) et l'astrophotographie grand champ sur monture EQ (si budget monture inclus).
    \end{itemize}
\end{enumerate}
\end{corrige}

\coursec{Principe de fonctionnement et formation des images}

Dans la section précédente, nous avons décrit les trois organes essentiels du télescope de Newton : le \cle{miroir primaire concave} qui collecte la lumière au fond du tube, le \cle{miroir secondaire plan} incliné à $45^{\circ}$ qui dévie le faisceau vers la paroi latérale, et l'\cle{oculaire} placé sur le côté du tube, dans lequel l'observateur plonge son œil. Il reste maintenant à répondre à la question fondamentale : \emph{comment ces trois éléments collaborent-ils pour transformer la lumière d'un astre lointain en une image observable confortablement ?} C'est l'objet de cette section. Nous allons suivre la lumière pas à pas, depuis l'astre jusqu'à la rétine, en justifiant chaque étape.

\courssub{La lumière d'un astre : un faisceau parallèle incliné}

\begin{definition}[Astre à l'infini et diamètre apparent]
Un astre observé depuis la Terre (étoile, planète, Lune) est situé à une distance immense devant la taille de l'instrument : on dit qu'il est \cle{à l'infini}. Tous les rayons lumineux issus d'un même point $B$ de l'astre arrivent alors sur le télescope sous forme d'un \cle{faisceau de rayons parallèles} entre eux. Deux points extrêmes $A$ et $B$ de l'astre sont vus depuis la Terre sous un angle $\theta$ appelé \cle{diamètre apparent} de l'astre, exprimé en radians (rad).
\end{definition}

Concrètement, la Lune est vue sous un diamètre apparent $\theta \approx 9{,}3 \times 10^{-3}\ \text{rad}$ (environ $0{,}53^{\circ}$), Jupiter sous un diamètre apparent de l'ordre de $10^{-4}\ \text{rad}$ lors de ses approches favorables. Ces angles sont \emph{très petits} : c'est toute la difficulté de l'astronomie, et toute la raison d'être du télescope.

\begin{attention}[Deux faisceaux, pas un seul]
Une erreur fréquente consiste à ne dessiner qu'un seul faisceau parallèle. En réalité, chaque point de l'astre envoie \emph{son propre} faisceau parallèle : le point $A$ (bord inférieur de l'astre) envoie un faisceau parallèle à l'axe optique si l'astre est centré, et le point $B$ (bord supérieur) envoie un faisceau \emph{incliné de l'angle $\theta$} par rapport au premier. C'est l'écart angulaire entre ces deux faisceaux qui contient toute l'information sur la taille de l'astre.
\end{attention}

\courssub{Le miroir primaire : formation de l'image réelle intermédiaire}

\begin{propriete}[Focalisation par le miroir concave — admise]
Pour un miroir sphérique concave utilisé dans les conditions de Gauss (rayons proches de l'axe et peu inclinés), tout faisceau de rayons parallèles à l'axe optique converge, après réflexion, en un point $F$ appelé \cle{foyer} du miroir, situé sur l'axe à la distance
\[ f = \dfrac{R}{2} \]
du sommet $S$ du miroir, où $R$ est le rayon de courbure de la calotte sphérique. Plus généralement, un faisceau parallèle \emph{incliné} d'un angle $\theta$ par rapport à l'axe converge en un point $B'$ situé dans le \cle{plan focal} (plan perpendiculaire à l'axe passant par $F$), hors de l'axe. Cette propriété est admise en Première ; elle se démontre en Terminale à partir des lois de la réflexion.
\end{propriete}

Appliquons cette propriété à notre télescope. Le faisceau parallèle issu du point $B$ de l'astre, incliné de l'angle $\theta$, frappe le grand miroir concave et converge après réflexion en un point $B'$ du plan focal. Le faisceau issu du point $A$ (sur l'axe) converge en $A' = F$. L'ensemble des points images ainsi formés constitue une \cle{image réelle intermédiaire} $A'B'$ de l'astre.

\begin{aretenir}[Nature de l'image intermédiaire]
L'image formée par le miroir primaire dans son plan focal est :
\begin{itemize}
\item \cle{réelle} : les rayons convergent effectivement en $A'B'$ ; on pourrait la recueillir sur un écran placé dans le plan focal ;
\item \cle{renversée} : le point $B$ « en haut » de l'astre donne un point image $B'$ \emph{en dessous} de l'axe ;
\item de \cle{taille très petite} : quelques millimètres seulement, comme nous allons le calculer.
\end{itemize}
\end{aretenir}

\begin{propriete}[Taille de l'image intermédiaire — démonstration guidée]
Pour un astre de diamètre apparent $\theta$ (en radians), la taille de l'image intermédiaire dans le plan focal d'un miroir de distance focale $f$ est :
\[ A'B' = f\,\tan\theta \approx f\,\theta \qquad (\theta \text{ en rad, petit angle}). \]
\emph{Démonstration guidée.} Considérons le rayon issu de $B$ qui arrive au sommet $S$ du miroir : il s'y réfléchit symétriquement par rapport à l'axe (loi de la réflexion : angle d'incidence égal à l'angle de réflexion) et coupe le plan focal en $B'$. Dans le triangle $SFB'$, rectangle en $F$ :
\begin{enumerate}
\item le côté $SF$ a pour longueur la distance focale $f$ ;
\item l'angle en $S$ vaut $\theta$ (le rayon incident fait l'angle $\theta$ avec l'axe, et la réflexion en $S$ conserve cet angle par symétrie) ;
\item le côté opposé à cet angle est précisément $FB' = A'B'$.
\end{enumerate}
La trigonométrie dans ce triangle donne $\tan\theta = \dfrac{A'B'}{f}$, d'où $A'B' = f\tan\theta$. Comme les diamètres apparents des astres sont très petits, on utilise l'approximation des petits angles $\tan\theta \approx \theta$ (avec $\theta$ en radians), d'où le résultat. \emph{Cette démonstration est exigible : sache la refaire sur une figure propre.}
\end{propriete}

\begin{exemplebox}[Image de la Lune dans le plan focal]
Un télescope amateur possède un miroir primaire de distance focale $f = \SI{1,20}{m}$. On observe la pleine Lune, de diamètre apparent $\theta = 9{,}3 \times 10^{-3}\ \text{rad}$. Taille de l'image intermédiaire :
\[ A'B' = f\,\theta = 1{,}20 \times 9{,}3 \times 10^{-3} = 1{,}1 \times 10^{-2}\ \text{m} \approx \SI{11}{mm}. \]
Toute la Lune tient dans une image réelle d'un centimètre ! Vérification dimensionnelle : $f$ est en mètres, $\theta$ est sans dimension (le radian est un rapport de longueurs), donc $f\theta$ est bien une longueur. C'est cette petite image lumineuse que l'oculaire va ensuite « examiner » comme une loupe.
\end{exemplebox}

\courssub{Le miroir plan secondaire : déport du plan focal}

Si l'on plaçait l'œil directement dans le plan focal du miroir primaire, la tête de l'observateur bloquerait la lumière incidente : c'est précisément le problème que Newton a résolu en 1668. Le \cle{miroir plan secondaire}, petit miroir incliné à $45^{\circ}$ sur l'axe, intercepte le faisceau convergent \emph{avant} le foyer et le renvoie perpendiculairement à l'axe du tube, vers la paroi latérale. Tout se passe comme si le plan focal était « déporté » sur le côté du tube : c'est là que se forme effectivement l'image intermédiaire $A'B'$, et c'est là que l'on place l'oculaire.

\begin{attention}[Le miroir plan est neutre optiquement]
Le miroir plan ne modifie \emph{ni la taille, ni la position relative} de l'image intermédiaire : il se contente de « plier » le trajet de la lumière de $90^{\circ}$. La distance du miroir plan au point où se forme $A'B'$ est simplement reportée latéralement. Ne dis jamais que le miroir plan « grossit » ou « reforme » l'image : c'est un simple déviateur. (Il introduit seulement une petite obstruction au centre du faisceau incident, sans conséquence notable sur l'image.)
\end{attention}

\courssub{L'oculaire : de l'image intermédiaire à l'image définitive à l'infini}

L'image intermédiaire $A'B'$ est réelle mais minuscule : l'œil nu ne pourrait l'examiner en détail qu'en l'approchant à moins de la distance minimale de vision distincte, ce qui est impossible ici. On interpose donc l'\cle{oculaire}, une lentille convergente de courte distance focale $f'_{oc}$ (ou un ensemble de lentilles jouant ce rôle), qui joue le rôle de \emph{loupe} pour observer $A'B'$.

\begin{definition}[Réglage afocal]
Le télescope est réglé de sorte que l'image intermédiaire $A'B'$ se forme dans le \cle{plan focal objet de l'oculaire} : le foyer objet $F_{oc}$ de l'oculaire coïncide avec $A'B'$. D'après les propriétés des lentilles convergentes, tout rayon issu d'un point du plan focal objet émerge de la lentille \emph{parallèle} aux autres rayons issus du même point : l'oculaire renvoie donc l'image définitive \cle{à l'infini}. Le télescope, qui reçoit un objet à l'infini et donne une image à l'infini, est un \cle{système afocal}.
\end{definition}

L'intérêt pratique est capital : un œil normal (emmetrope) qui regarde un objet à l'infini \emph{n'accommode pas}, c'est-à-dire que son cristallin est au repos. L'observation à travers le télescope est donc \cle{confortable et sans fatigue}, même pendant de longues séances d'observation du ciel étoilé.

\begin{aretenir}[Nature de l'image définitive]
L'image définitive donnée par le télescope de Newton réglé en afocal est :
\begin{itemize}
\item \cle{virtuelle} : elle ne peut pas être recueillie sur un écran, elle est seulement « vue » par l'œil ;
\item située \cle{à l'infini} ;
\item \cle{renversée} par rapport à l'astre (le miroir primaire renverse une fois, l'oculaire ne redresse pas). En astronomie, ce renversement est \emph{sans conséquence} : il n'y a ni haut ni bas dans l'espace ! Pour l'observation terrestre, on ajouterait un système redresseur.
\end{itemize}
\end{aretenir}

\begin{methode}[Analyser le fonctionnement en deux étapes]
Face à toute question sur le télescope de Newton, décompose systématiquement :
\begin{enumerate}
\item \textbf{Étape 1 — le miroir primaire} : l'astre à l'infini donne une image \emph{réelle, renversée}, de taille $f\theta$, dans le plan focal du miroir (déporté latéralement par le miroir plan) ;
\item \textbf{Étape 2 — l'oculaire} : cette image intermédiaire, placée au foyer objet de l'oculaire, est transformée en faisceaux \emph{parallèles} sortants : l'image définitive est rejetée à l'infini, où l'œil la reçoit sans accommoder.
\end{enumerate}
Résumé de la chaîne complète :
\[ \text{astre }(\infty) \xrightarrow{\ \text{miroir concave}\ } \text{image réelle } A'B' \xrightarrow{\ \text{miroir plan}\ } \text{renvoi latéral} \xrightarrow{\ \text{oculaire}\ } \text{image }(\infty) \xrightarrow{\ \text{œil}\ } \text{rétine}. \]
\end{methode}

\courssub{Tracé de la marche des rayons}

\begin{popfigure}
\begin{tikzpicture}[scale=1.0,>=Stealth]
% Tube (indication)
\draw[popmuted,dashed] (-0.2,1.6) -- (8.6,1.6);
\draw[popmuted,dashed] (-0.2,-1.6) -- (8.6,-1.6);
% Miroir primaire concave (arc vu par la concavité vers la gauche)
\draw[popblue,very thick] (8.2,1.5) .. controls (7.55,0.6) and (7.55,-0.6) .. (8.2,-1.5);
\draw[popblue] (8.2,1.5) -- (8.35,1.62) (8.2,-1.5) -- (8.35,-1.62);
\node[popblue] at (8.75,0) {\small primaire};
% Axe optique
\draw[popmuted] (-0.2,0) -- (8.8,0);
% Miroir plan secondaire à 45°
\draw[poporange,very thick] (4.6,-0.55) -- (5.7,0.55);
\node[poporange,anchor=south east] at (4.6,0.6) {\small miroir plan $45^{\circ}$};
% Faisceau parallèle incliné arrivant de la gauche
\draw[->,popgreen,thick] (0,1.05) -- (7.72,0.62);
\draw[->,popgreen,thick] (0,0.35) -- (7.80,0.20);
\draw[->,popgreen,thick] (0,-0.35) -- (7.80,-0.22);
% Rayons réfléchis convergents vers le miroir plan
\draw[->,popgreen,thick] (7.72,0.62) -- (5.35,0.10);
\draw[->,popgreen,thick] (7.80,0.20) -- (5.15,-0.06);
\draw[->,popgreen,thick] (7.80,-0.22) -- (4.95,-0.22);
% Rayons renvoyés vers le bas (latéral) par le miroir plan
\draw[->,popgreen,thick] (5.35,0.10) -- (5.55,-1.15);
\draw[->,popgreen,thick] (5.15,-0.06) -- (5.15,-1.15);
\draw[->,popgreen,thick] (4.95,-0.22) -- (4.75,-1.15);
% Image intermédiaire (plan focal déporté)
\draw[popink,thick] (4.55,-1.15) -- (5.75,-1.15);
\node[popink,anchor=north] at (5.15,-1.2) {\small $A'B'$};
% Oculaire (lentille biconvexe) en bas
\draw[poppurple,very thick] (4.55,-2.0) .. controls (4.9,-1.75) and (5.4,-1.75) .. (5.75,-2.0);
\draw[poppurple,very thick] (4.55,-2.0) .. controls (4.9,-2.25) and (5.4,-2.25) .. (5.75,-2.0);
\node[poppurple,anchor=west] at (5.85,-2.0) {\small oculaire};
% Faisceau parallèle sortant vers l'œil
\draw[->,popgreen,thick] (5.55,-2.05) -- (5.85,-3.3);
\draw[->,popgreen,thick] (5.15,-2.0) -- (5.15,-3.3);
\draw[->,popgreen,thick] (4.75,-2.05) -- (4.45,-3.3);
% Œil
\draw[popdark,thick] (4.6,-3.5) .. controls (5.15,-3.85) .. (5.7,-3.5);
\draw[popdark,thick] (4.6,-3.5) .. controls (5.15,-3.2) .. (5.7,-3.5);
\fill[popdark] (5.15,-3.5) circle (0.07);
\node[popdark,anchor=west] at (5.8,-3.5) {\small œil};
% Angle theta incident
\draw[->,popgreen] (1.9,0) arc[start angle=0,end angle=17,radius=1.9];
\node[popgreen] at (2.35,0.32) {$\theta$};
\node[popgreen,anchor=south] at (0.1,1.15) {\small faisceau de l'astre ($\infty$)};
\end{tikzpicture}
\legende{Marche d'un faisceau lumineux dans le télescope de Newton : le faisceau parallèle incliné de $\theta$ converge après réflexion sur le miroir primaire ; le miroir plan à $45^{\circ}$ le renvoie latéralement où se forme l'image intermédiaire $A'B'$ ; l'oculaire la transforme en faisceau parallèle reçu par l'œil.}
\end{popfigure}

Pour établir les relations de grossissement (section suivante), il est commode de « déplier » le schéma, c'est-à-dire d'ignorer le miroir plan (qui ne change rien aux propriétés de l'image) et de représenter le miroir primaire comme une lentille convergente équivalente de distance focale $f$ :

\begin{popfigure}
\begin{tikzpicture}[scale=1.0,>=Stealth]
% Axe
\draw[popmuted,->] (-0.3,0) -- (11,0);
% Lentille équivalente (miroir primaire déplié)
\draw[popblue,very thick,<->] (3,-1.6) -- (3,1.6);
\node[popblue] at (3,2.0) {\small primaire ($f$)};
% Foyer F
\fill[popblue] (6.2,0) circle (0.06); \node[popblue,anchor=north] at (6.2,-0.05) {\small $F=A'$};
% Image intermédiaire B'
\draw[->,popink,thick] (6.2,0) -- (6.2,-0.75); \node[popink,anchor=west] at (6.3,-0.6) {\small $B'$};
% Rayon incliné theta passant par le centre
\draw[->,popgreen,thick] (0,1.5) -- (6.2,-0.75);
\draw[popgreen,thick,dashed] (6.2,-0.75) -- (8.5,-1.18);
\draw[->,popgreen] (1.6,0) arc[start angle=0,end angle=-26.5,radius=1.6];
\node[popgreen] at (2.0,-0.42) {$\theta$};
% Oculaire
\draw[poppurple,very thick,<->] (7.4,-1.6) -- (7.4,1.6);
\node[poppurple] at (7.4,2.0) {\small oculaire ($f'_{oc}$)};
% Faisceau parallèle sortant incliné theta'
\draw[->,popgreen,thick] (7.4,-0.62) -- (10.5,0.9);
\draw[->,popgreen,thick] (7.4,-0.3) -- (10.5,1.22);
\draw[->,popgreen,thick] (7.4,-0.94) -- (10.5,0.58);
\draw[->,popgreen] (9.3,0) arc[start angle=0,end angle=26.5,radius=1.9];
\node[popgreen] at (9.75,0.5) {$\theta'$};
% Distances focales
\draw[<->,popmuted] (3,-1.85) -- (6.2,-1.85); \node[popmuted] at (4.6,-2.15) {\small $f$};
\draw[<->,popmuted] (6.2,-1.85) -- (7.4,-1.85); \node[popmuted] at (6.8,-2.15) {\small $f'_{oc}$};
\end{tikzpicture}
\legende{Schéma « déplié » du télescope : le miroir primaire est remplacé par une lentille équivalente de focale $f$. Le faisceau incident fait l'angle $\theta$ avec l'axe ; le faisceau émergent fait l'angle $\theta' > \theta$ : l'astre apparaît grossi. L'image intermédiaire $A'B'$ est au foyer objet de l'oculaire (montage afocal).}
\end{popfigure}

Sur ce schéma déplié, on lit directement la double action de l'instrument : le miroir primaire \emph{matérialise} l'angle $\theta$ sous forme d'une petite image réelle $A'B' = f\theta$, puis l'oculaire \emph{reconvertit} cette image en un faisceau parallèle sortant sous un angle $\theta'$ plus grand. L'œil reçoit donc l'astre sous un angle apparent augmenté : c'est le grossissement, que nous quantifierons dans la section 4.

\begin{attention}[Sens de lecture du schéma déplié]
Sur le schéma déplié, la lumière semble « traverser » le miroir primaire représenté comme une lentille : c'est une \emph{astuce de représentation} rendue possible parce que le miroir concave et la lentille convergente ont le même effet focalisateur. Dans l'instrument réel, la lumière fait un aller-retour dans le tube. Ne confonds pas les deux représentations dans ta copie : annonce toujours clairement « schéma déplié équivalent » si tu l'utilises.
\end{attention}

\begin{exoresolu}[Image intermédiaire de Jupiter]
Énoncé. Un télescope de Newton scolaire possède un miroir primaire de rayon de courbure $R = \SI{1,60}{m}$. On observe Jupiter, de diamètre apparent $\theta = 1{,}2 \times 10^{-4}\ \text{rad}$. (a) Déterminer la distance focale du miroir primaire. (b) Calculer la taille de l'image intermédiaire de Jupiter. (c) L'oculaire a une distance focale $f'_{oc} = \SI{20}{mm}$. À quelle distance de l'oculaire doit se former l'image intermédiaire pour une observation sans accommodation ?
\tcblower
Solution détaillée. (a) Pour un miroir sphérique concave, $f = \dfrac{R}{2} = \dfrac{1{,}60}{2} = \SI{0,80}{m}$.
(b) Taille de l'image intermédiaire, avec l'approximation des petits angles (justifiée car $\theta \ll 1\ \text{rad}$) :
\[ A'B' = f\,\theta = 0{,}80 \times 1{,}2 \times 10^{-4} = 9{,}6 \times 10^{-5}\ \text{m} \approx \SI{0,10}{mm}. \]
L'image réelle de Jupiter ne mesure qu'un dixième de millimètre : impossible à examiner à l'œil nu, d'où le rôle indispensable de l'oculaire.
(c) Pour une observation sans accommodation, le système doit être afocal : l'image intermédiaire doit coïncider avec le foyer objet de l'oculaire, donc se former à la distance $f'_{oc} = \SI{20}{mm}$ en amont de l'oculaire. C'est exactement ce réglage que l'on réalise lors de la \emph{mise au point} (section 5) en faisant coulisser le porte-oculaire.
\end{exoresolu}

\begin{savaistu}[Pourquoi Newton a-t-il choisi un miroir ?]
À l'époque de Newton (1668), les lunettes astronomiques à lentilles souffraient d'un défaut redoutable, les \emph{aberrations chromatiques} : la lentille se comporte comme un prisme et décompose la lumière blanche, entourant les images d'irisations colorées. Newton crut (à tort, on le saura plus tard) ce défaut irrémédiable et conçut un instrument \emph{à réflexion} : un miroir réfléchit toutes les couleurs de la même façon, donc sans aucune aberration chromatique. C'est l'\emph{intérêt du télescope sur les instruments à lentilles}, que nous détaillerons en section 6. Aujourd'hui encore, tous les grands télescopes professionnels — y compris le télescope spatial James Webb — sont des instruments à miroirs, héritiers directs du télescope de Newton.
\end{savaistu}

\coursec{Tracé de la marche d'un faisceau lumineux}

Dans la section précédente, nous avons établi le principe de fonctionnement : le miroir primaire concave forme une image réelle, inversée et réduite de l'astre au sein de son plan focal ; le miroir plan renvoie cette image vers le côté du tube ; l'oculaire (une lentille convergente) fonctionne alors comme une loupe afocale, restituant un faisceau parallèle pour l'œil détendu. Nous allons maintenant maîtriser le \cle{tracé géométrique rigoureux} de la marche des rayons lumineux, compétence incontournable au Probatoire. Ce tracé se décompose en étapes logiques, que nous allons illustrer pas à pas.

\courssub{Rappel des règles de tracé pour le miroir concave}

Avant d'aborder l'instrument complet, rappelons les règles fondamentales pour la construction de l'image par un miroir concave de sommet $S_1$, de centre $C_1$ et de foyer $F_1$ (distance focale $F = \overline{S_1F_1} > 0$ en géométrie réelle ; $f' = -F$ en convention algébrique de Descartes). Nous travaillons ici en \textbf{géométrie réelle} (grandeurs positives, angles mesurés au rapporteur), cadre usuel des tracés au Probatoire.

\begin{propriete}[Règles de tracé pour un miroir concave (Géométrie de Gauss)]
Pour un miroir concave de sommet $S_1$, de foyer $F_1$ et de centre $C_1$ ($S_1C_1 = 2F$) :
\begin{enumerate}
    \item \textbf{Rayon incident parallèle à l'axe optique principal $\Delta_1$ :} Le rayon réfléchi passe par le foyer $F_1$.
    \item \textbf{Rayon incident passant par le foyer $F_1$ :} Le rayon réfléchi est parallèle à l'axe optique principal $\Delta_1$.
    \item \textbf{Rayon incident passant par le centre de courbure $C_1$ :} Le rayon réfléchi revient sur lui-même (incidence normale, la normale au miroir en ce point est la droite $(C_1I)$).
    \item \textbf{Rayon incident passant par le sommet $S_1$ :} Le rayon réfléchi est symétrique du rayon incident par rapport à l'axe optique $\Delta_1$ (loi de Descartes : $i = r$, la normale en $S_1$ est l'axe $\Delta_1$).
\end{enumerate}
Ces règles découlent de la loi de la réflexion ($i = r$) et de l'approximation de Gauss (miroir de faible ouverture, surface sphérique confondue avec son plan tangent au sommet pour les rayons paraxiaux).
\end{propriete}

\begin{attention}[Convention de signe et géométrie réelle]
Au Probatoire, les tracés sont réalisés en \textbf{géométrie réelle} : on dessine les grandeurs positives à l'échelle. Le foyer $F_1$ est placé \emph{devant} le miroir (côté objet), à une distance $S_1F_1 = F$. On ne mélange pas le tracé géométrique (figure à l'échelle) et les formules algébriques (relation de Descartes $\frac{1}{p'} + \frac{1}{p} = \frac{1}{f'}$ avec $f'<0$).
\end{attention}

\courssub{Construction géométrique pas à pas du tracé complet (Schéma réel plié)}

Considérons un télescope de Newton de focale primaire $F$ (miroir $M_1$) et d'oculaire de focale $f'$ (lentille $L_{oc}$). L'astre est à l'infini ; un point $B$ de cet astre (haut du champ) émet un faisceau de rayons parallèles entre eux, inclinés d'un angle $\theta$ par rapport à l'axe optique principal $\Delta_1$.

\begin{methode}[Tracé complet en 4 étapes (Schéma plié)]
\textbf{Matériel :} Règle graduée, rapporteur, compas, papier millimétré (format A3 recommandé). \textbf{Échelle :} Choisie pour que le tube tienne sur la feuille (ex. 1/10, 1/20, 1/50).
\begin{enumerate}
    \item \textbf{Repérage de l'optique primaire (Miroir $M_1$) :}
        \begin{itemize}
            \item Tracer l'axe optique principal $\Delta_1$ (horizontal).
            \item Placer le sommet $S_1$ du miroir primaire. Tracer le miroir concave (arc de cercle de centre $C_1$, rayon $R_1 = 2F$). En approximation de Gauss, on représente souvent seulement la portion utile au sommet (diamètre $D$) et on trace les normales (rayons $C_1I$) aux points d'incidence.
            \item Repérer le foyer $F_1$ sur $\Delta_1$ tel que $S_1F_1 = F$. Tracer le plan focal $\mathcal{P}_{F_1}$ (perpendiculaire à $\Delta_1$ en $F_1$).
        \end{itemize}
    \item \textbf{Formation de l'image intermédiaire $B'$ (Faisceau incident incliné $\theta$) :}
        \begin{itemize}
            \item Depuis l'infini (bord gauche de la feuille), tracer \textbf{deux rayons incidents parallèles} inclinés de l'angle $\theta$ (ex. $\theta = 0,5^\circ$ ou $1^\circ$) : 
                \begin{itemize}
                    \item \textbf{Rayon principal (ou chef de faisceau) :} Passe par le sommet $S_1$ (hauteur 0).
                    \item \textbf{Rayon marginal :} Passe par le bord supérieur du miroir (hauteur $y = D/2$).
                \end{itemize}
            \item \textbf{Construction des réflexions sur $M_1$ :}
                \begin{itemize}
                    \item \textit{Rayon principal :} Passe par $S_1$. La normale est l'axe $\Delta_1$. Angle d'incidence $i = \theta$. Rayon réfléchi symétrique : angle $-\theta$ par rapport à $\Delta_1$ (descend vers l'axe).
                    \item \textit{Rayon marginal :} Point d'incidence $I$ sur le miroir. Tracer la normale $(C_1I)$. Mesurer l'angle d'incidence $i$ entre le rayon incident et la normale. Tracer le rayon réfléchi de l'autre côté de la normale tel que $r = i$.
                \end{itemize}
            \item \textbf{Image $B'$ :} Le point d'intersection des deux rayons réfléchis définit l'image réelle $B'$. \textbf{Vérification :} $B'$ se situe sur le plan focal $\mathcal{P}_{F_1}$. Comme le rayon principal se réfléchit vers le bas (angle $-\theta$), $B'$ est du côté opposé à l'axe par rapport au rayon incident $\rightarrow$ \textbf{image inversée}.
        \end{itemize}
    \item \textbf{Réflexion sur le miroir plan $M_2$ (Déviation à $90^\circ$) :}
        \begin{itemize}
            \item Placer le miroir plan $M_2$ (segment incliné à $45^\circ$ par rapport à $\Delta_1$) \textbf{avant} le foyer $F_1$ (distance $S_1M_2 \approx F - \delta$, $\delta$ quelques cm). Son plan coupe l'axe $\Delta_1$ en $O$.
            \item Pour \textbf{chaque rayon} arrivant sur $M_2$ (les deux rayons réfléchis par $M_1$ convergent vers $B'$) :
                \begin{itemize}
                    \item Tracer la normale au miroir plan (perpendiculaire à $M_2$, donc à $45^\circ$ ou $135^\circ$ selon l'orientation).
                    \item Mesurer l'angle d'incidence $i$ ; tracer le rayon réfléchi tel que $r = i$ de l'autre côté de la normale.
                \end{itemize}
            \item Le faisceau convergent est dévié de \textbf{$90^\circ$} vers le haut (nouvel axe optique $\Delta_2$, vertical). L'image $B'$ devient l'objet réel $B''$ pour l'oculaire, situé symétriquement de $B'$ par rapport au plan du miroir $M_2$.
        \end{itemize}
    \item \textbf{Émergence par l'oculaire $L_{oc}$ (Système afocal) :}
        \begin{itemize}
            \item Placer la lentille mince convergente (oculaire) de centre optique $O_{oc}$ sur l'axe $\Delta_2$ de sorte que son foyer objet $F_{oc}$ \textbf{coïncide} avec l'objet $B''$ (donc distance $O_{oc}B'' = O_{oc}F_{oc} = f'$).
            \item Tracer les rayons incidents sur l'oculaire (ceux issus de $M_2$ convergent vers $B''$) :
                \begin{itemize}
                    \item Rayon passant par $O_{oc}$ : non dévié.
                    \item Rayon parallèle à $\Delta_2$ (ou tout autre rayon) : réfracté selon les règles de la lentille mince (passe par $F'_{oc}$ si incident // $\Delta_2$).
                \end{itemize}
            \item Les rayons émergents sont \textbf{parallèles entre eux}, inclinés d'un angle $\theta'$ par rapport à $\Delta_2$. L'instrument est \textbf{afocal} (entrée parallèle, sortie parallèle). L'œil placé derrière l'oculaire reçoit un faisceau parallèle : observation sans accommodation (œil détendu).
        \end{itemize}
\end{enumerate}
\end{methode}

\begin{popfigure}
\begin{tikzpicture}[scale=0.75, >=Stealth]
% Styles
\tikzset{
    mirror/.style={thick, popblue},
    ray/.style={poporange, thick, ->},
    rayconstr/.style={popgreen, thin, dashed},
    axis/.style={popdark, thin, dashed},
    normal/.style={poppurple, thin, dotted},
    point/.style={fill=popdark, circle, inner sep=1.5pt},
    labelstyle/.style={font=\small, popdark}
}
% --- FIGURE (a) : Miroir primaire - Formation image B' ---
\begin{scope}[xshift=0cm]
% Axe optique
\draw[axis] (-1.5,0) -- (9,0) node[right, labelstyle] {$\Delta_1$};
% Miroir concave : Sommet S1(0,0), Centre C1(4,0), Rayon R=4, Focale F=2.
% Arc supérieur : (0,0) arc (180:90:4) -> arrive en (4,4). Centre calcule: (0-4*cos180, 0-4*sin180) = (4,0). Correct.
% Arc inférieur : (0,0) arc (180:270:4) -> arrive en (4,-4).
\draw[mirror] (0,0) arc (180:90:4);
\draw[mirror] (0,0) arc (180:270:4);
% Points caractéristiques
\node[point, label=below left:{$S_1$}] at (0,0) {};
\node[point, label=below:{$F_1$}] at (2,0) {};
\node[point, label=above left:{$C_1$}] at (4,0) {};
% Diamètre utile D = 3 (hauteur 1.5)
\def\DemiD{1.5}
% Rayons incidents (angle theta = 15 deg pour visibilite, {tan(15)=0.268})
% Hauteur rayon principal a S1 = 0.
% Hauteur rayon marginal a S1 = 1.5 (bord miroir).
% Direction theta = 15 deg. Vecteur (cos15, sin15) ~ (0.966, 0.259).
% Rayon principal : passe par S1(0,0). Incident depuis gauche.
\draw[ray] (-1.5, -0.39) -- (0,0) node[midway, below, labelstyle, sloped] {$\theta$};
% Rayon marginal : passe par (0, 1.5). Incident depuis gauche.
\draw[ray] (-1.5, 1.11) -- (0, 1.5) node[midway, above, labelstyle, sloped] {$\theta$};
% Normale pour rayon marginal (C1I). I = (0, 1.5) sur miroir? Non, miroir est arc de cercle.
% Point d'incidence I sur miroir a hauteur y=1.5. Equation miroir (x-4)^2 + y^2 = 16.
% (x-4)^2 = 16 - 2.25 = 13.75 -> x = 4 - sqrt(13.75) = 4 - 3.71 = 0.29.
% Pour simplifier le dessin (approximation Gauss), on considere l'incidence au plan tangent en S1 (x=0).
% Mais pour rigueur "Probatoire", on trace la normale au point d'incidence reel.
% Ici, on utilise l'approximation Gauss standard cours lycee: Miroir ~ Plan tangent en S1 pour les hauteurs.
% Donc incidence a (0, 1.5). Normale = Axe Delta1 (horizontal).
% Rayon incident angle 15 deg. Normale angle 0 deg. i = 15 deg.
% Reflechi angle -15 deg. Passe par F1(2,0)?
% Equation reflechi: y - 1.5 = tan(-15)(x - 0) -> y = 1.5 - 0.268x.
% A x=2, y = 1.5 - 0.536 = 0.964. Ne passe pas par F1(2,0).
% ERREUR PEDAGOGIQUE CLASSIQUE : Le rayon marginal // a l'axe passe par F1. Le rayon marginal INCLINE ne passe pas par F1.
% METHODE CORRECTE POUR LE COURS :
% 1. Tracer rayon principal (par S1) -> reflechi symetrique (angle -theta). Intersection avec plan focal x=2 -> B'.
%    Hauteur B' = -2 * tan(15) = -0.536.
% 2. Tracer rayon marginal (hauteur 1.5, angle theta). Construire sa reflection par la normale (C1I).
%    On trace la normale C1I. On mesure i. On trace r=i. On verifie qu'il passe par B'.
% C'est ce que je vais dessiner.

% Coordonnees reelles pour TikZ (R=4, F=2, C1=(4,0))
% Point incidence I sur miroir pour y=1.5 : x_I = 4 - sqrt(16 - 2.25) = 0.29.
% Normale = droite C1(4,0) -> I(0.29, 1.5). Pente = (1.5-0)/(0.29-4) = 1.5/-3.71 = -0.404. Angle normale = -22 deg (ou 158 deg).
% Rayon incident : angle 15 deg.
% Angle incidence i = 15 - (-22) = 37 deg? Non, angles orientes.
% Normal vector N = I - C1 = (-3.71, 1.5). Angle N = 158 deg.
% Incident vector v_i = (cos15, sin15) = (0.966, 0.259). Angle v_i = 15 deg.
% Angle entre v_i et N = 158 - 15 = 143 deg. Angle aigu = 37 deg. i = 37 deg.
% Reflechi : symetrique de v_i par rapport a N. Angle v_r = 158 + 37 = 195 deg (ou -165 deg).
% Pente tan(-165) = tan(15) = 0.268. Direction (-0.966, -0.259).
% Equation depuis I(0.29, 1.5) : y - 1.5 = 0.268(x - 0.29).
% A x=2 (plan focal), y = 1.5 + 0.268(1.71) = 1.5 + 0.46 = 1.96. Ne passe pas par B'(2, -0.54).
% PROBLEME : Ouverture trop grande (D=3, F=2 -> f/1.3). Approximation Gauss non valide pour rayon marginal.
% REQUIS PROGRAMME : "Approximation de Gauss : miroir de faible ouverture".
% Il FAUT choisir D << F. Ex: F=8cm (echelle), D=2cm.
% Je vais redessiner avec des valeurs coherentes pour l'approximation de Gauss (paraxial).
% F=4cm (sur dessin). D=1cm. Theta=10 deg.
% S1=(0,0), F1=(4,0), C1=(8,0). Miroir arc (180:90:8) etc. Trop grand pour scope.
% SIMPLIFICATION PEDAGOGIQUE STANDARD (Manuels Cameroun) :
% On trace le miroir comme une ligne verticale legere ou un arc tres plat au sommet.
% On trace les rayons incidents sur le PLAN TANGENT en S1 (x=0).
% 1. Rayon principal (0,0) angle theta -> reflechi angle -theta.
% 2. Rayon marginal (0, h) angle theta. 
%    Regle "Miroir plan tangent" : Reflection symetrique par rapport a l'axe -> angle -theta. 
%    Les deux rayons reflechis sont PARALLELES (angle -theta). Ils ne se coupent pas -> Image a l'infini. FAUX.
% IL FAUT LA COURBURE.
% La construction standard "Bac/Probatoire" pour miroir concave objet a l'infini incline :
% 1. Tracer l'axe, le miroir (arc), F, C.
% 2. Tracer le rayon PARALLELE A L'AXE (objet AXE) -> Reflechi par F. Cela place F et le plan focal.
% 3. Pour l'objet INCLINE (angle theta) :
%    - Rayon Principal (par S1) -> Reflechi symetrique (angle -theta). Coupe le plan focal en B'.
%    - Rayon passant par C (Centre) -> Reflechi sur lui-meme. 
%      MAIS un rayon venant de l'infini incline ne passe pas par C (sur l'axe).
%    - Rayon passant par F -> Reflechi // axe. 
%      Un rayon venant de l'infini incline ne passe pas par F (sur l'axe).
% CONCLUSION : En geometrie de Gauss (paraxial), pour un objet a l'infini INCLINE, 
% SEUL le rayon principal (par S1) est traceable simplement par une regle simple.
% Les autres rayons necessitent la construction par la NORMALE (C1I).
% AU PROBATOIRE : On trace le rayon principal pour trouver B' sur le plan focal.
% Puis on trace le rayon marginal EN CONSTRUISANT SA NORMALE (C1I) et en appliquant i=r.
% On verifie qu'il passe par B' (stigmatisme approché).

% DESSIN TIKZ CORRIGE (Approximation visuelle pour le cours) :
% On dessine le miroir peu ouvert. 
% Rayon principal : Par S1, reflechi symetrique -> B' sur plan focal.
% Rayon marginal : Incident a hauteur h (bord miroir), angle theta. 
% On trace la normale (vers C1). On trace le reflechi (i=r). 
% On montre qu'il coupe le rayon principal en B' sur le plan focal.

% Parametres dessin : F=4cm (unite tikz). C1=(8,0). Miroir arc (180:90:8) trop grand.
% On zoom sur le sommet. Miroir approx parabolique y^2 = 2px = 4Fx = 16x. x = y^2/16.
% Pour y=1 (D=2), x=1/16=0.06. Tres plat.
% On va dessiner un miroir "schema" : trait vertical epais a x=0 pour le sommet, avec flèches concaves.
% ET on trace les normales radiales pour le rayon marginal.

% NOUVELLE STRATEGIE TIKZ : Schéma pédagogique clair, pas forcement a l'echelle metrique exacte, mais geometriquement correct en logique.
% S1=(0,0). F1=(3,0). C1=(6,0). (F=3, R=6).
% Miroir : Arc de cercle centre (6,0) rayon 6. De y=-1.5 a y=1.5. x = 6 - sqrt(36-y^2).
% y=1.5 -> x = 6 - sqrt(33.75) = 6 - 5.81 = 0.19.
% Theta = 10 deg. tan(10)=0.176.
% Rayon Principal : (0,0) angle 10 deg incident. Reflechi angle -10 deg. 
%   Eq: y = -0.176 x. A x=3 (F1), y = -0.53. -> B'(3, -0.53).
% Rayon Marginal : Incident sur I(0.19, 1.5). Direction angle 10 deg.
%   Normale : C1(6,0) -> I(0.19, 1.5). Vecteur (-5.81, 1.5). Angle normale = 165.5 deg.
%   Incident angle = 10 deg.
%   Difference = 155.5 deg. i = 180-155.5 = 24.5 deg.
%   Reflechi angle = 165.5 + 24.5 = 190 deg (-170 deg). Pente tan(-170)=tan(10)=0.176.
%   Eq depuis I: y - 1.5 = 0.176(x - 0.19).
%   A x=3: y = 1.5 + 0.176(2.81) = 1.5 + 0.49 = 1.99. LOIN de B'(-0.53).
%   Ouverture encore trop grande (h/F = 1.5/3 = 0.5). Il faut h/F < 0.1.
%   F=10, h=1. C1=(20,0). I: x=20-sqrt(400-1)=20-19.97=0.025. Presque plan.
%   Normale quasi horizontale. i ~ theta. Reflechi ~ -theta. Convergence tres lente.
%   B' a x=10, y=-1.76.
%   Rayon marginal reflechi: depuis (0.025, 1) pente -0.176. A x=10, y = 1 - 0.176*9.97 = -0.75. Pas -1.76.
%   L'approximation de Gauss (stigmatisme parfait) n'est valable que pour les rayons PARAXIAUX (proches de l'axe, petits angles, petites hauteurs).
%   Le "rayon marginal" (bord du miroir) N'EST PAS un rayon paraxial si D/F n'est pas tres petit.
%   EN COURS / PROBATOIRE : On trace le rayon principal (paraxial) pour trouver B'.
%   On trace le rayon marginal EN APPROXIMATION GAUSSIENNE : on le fait passer par B' (stigmatisme impose).
%   Ou on trace la normale et on applique i=r, mais on accepte qu'il ne passe pas exactement par B' si l'ouverture est grande (aberration spherique).
%   Le programme dit "Tracer la marche d'un faisceau lumineux". 
%   La methode attendue : 
%   1. Tracer F1 avec rayon // axe.
%   2. Tracer plan focal.
%   3. Tracer rayon principal (par S1) -> B' sur plan focal.
%   4. Tracer le cone convergent : depuis les bords du miroir, tracer des rayons VERS B'.
%   C'est la representation standard du "faisceau convergent vers l'image reelle".

% J'ADOPTE LA METHODE PEDAGOGIQUE STANDARD (Cone convergent) :
% 1. Rayon // axe -> F1 (definit F1 et plan focal).
% 2. Rayon principal incline (par S1) -> reflechi symetrique -> coupe plan focal en B'.
% 3. Rayons marginaux (bords miroir) : on trace les rayons REFLECHIS convergent VERS B'.
%    (On ne trace pas l'incident du marginal incline, ou on le trace parallele au principal et on dit "par stigmatisme, le reflechi passe par B'").
%    Souvent on trace l'incident marginal parallele au principal, et le reflechi vers B', en notant que la loi i=r est verifiee approximativement (Gauss).

% TikZ Final pour Figure (a) :
% Miroir : Arc centre (6,0) rayon 6, de y=-1 a y=1 (D=2, F=3 -> f/1.5, encore grand mais visible).
% S1=(0,0), F1=(3,0), C1=(6,0).
% Plan focal x=3.
% Rayon axe (theta=0) : incident // axe hauteur 1 -> reflechi vers F1(3,0). (Trait fin, gris).
% Rayon principal (theta=15 deg) : Incident par S1(0,0) angle 15. Reflechi angle -15. Coupe x=3 en B'(3, -0.8).
% Rayon marginal haut : Incident hauteur 1 (sur miroir x~0.08), angle 15. 
%   Reflechi : trace VERS B' (stigmatisme Gauss).
%   Incident : trace depuis gauche vers point incidence.
% Normale au point incidence (vers C1).

\draw[axis] (-1,0) -- (7,0) node[right, labelstyle] {$\Delta_1$};
% Miroir (arc cercle centre 6,0 rayon 6)
\draw[mirror] (0.08, 1) arc (98.5:81.5:6); % y=1 -> angle asin(1/6)=9.6 deg. 90+9.6=99.6. 90-9.6=80.4. Centre (6,0).
% Correction arc TikZ: \draw (x,y) arc (start:end:radius). 
% Point haut (0.08, 1). Centre (6,0). Vecteur centre->point = (-5.92, 1). Angle = 170.4 deg.
% Point bas (0.08, -1). Angle = -170.4 deg (189.6).
% Arc de 170.4 a 189.6 ? Non, miroir concave face gauche. Points de y=-1 a y=1.
% Angle debut (haut) = 170.4. Angle fin (bas) = -170.4 (ou 189.6).
% \draw (0.08, 1) arc (170.4:189.6:6); 
% TikZ arc angles en degres. 
% Pour simplifier le code : \draw[mirror] (0,1) -- (0,-1); % Approximation plan tangent au sommet (Miroir plan) ? Non, doit etre concave.
% Utilisons une parabole approx : \draw[mirror] plot[domain=-1:1, samples=20] ({\x*\x/12}, \x); % F=3 -> p=6 -> y^2=12x -> x=y^2/12.
% y=1 -> x=1/12=0.083. y=0.5 -> x=0.02. Tres bon.
\draw[mirror, thick] plot[domain=-1:1, samples=30] ({\x*\x/12}, \x) node[above right, labelstyle, pos=0.8] {$M_1$};
% Points
\node[point, label=below left:{$S_1$}] at (0,0) {};
\node[point, label=below:{$F_1$}] at (3,0) {};
\node[point, label=above left:{$C_1$}] at (6,0) {};
% Plan focal
\draw[rayconstr] (3, -1.5) -- (3, 1.5) node[right, labelstyle] {Plan focal $\mathcal{P}_{F_1}$};
% Rayon axe (reference pour F1) - Gris
\draw[popmuted, thin, ->] (-1, 1) -- (0.083, 1) -- (3, 0) node[midway, above, sloped, labelstyle, popmuted] {Rayon axe $\to F_1$};
% Faisceau incline theta = 15 deg (tan=0.268)
% Rayon Principal (par S1)
\draw[ray] (-1, -0.268) -- (0,0) node[midway, below, labelstyle, sloped] {$\theta$};
\draw[ray] (0,0) -- (4, -0.7) node[midway, below right, labelstyle] {Rayon principal};
% Image B'
\node[point, label=right:{$B'$}] at (3, -0.8) {};
% Rayon Marginal Haut (Incident)
% Point incidence sur miroir pour y=1 : x=1/12=0.083.
\draw[ray] (-1, 0.732) -- (0.083, 1) node[midway, above, labelstyle, sloped] {$\theta$};
% Normale en I (vers C1(6,0))
\draw[normal] (6,0) -- (0.083, 1) node[midway, right, labelstyle, poppurple] {Normale $(C_1I)$};
% Rayon Marginal Reflechi (Vers B' par stigmatisme Gauss)
\draw[ray] (0.083, 1) -- (3, -0.8) node[midway, right, labelstyle] {Rayon marginal};
% Angle theta incident
\draw[popmuted, ->] (-0.5, 0.1) arc (0:15:0.5) node[right, labelstyle, popmuted] {$\theta$};
% Angle -theta reflechi principal
\draw[popmuted, ->] (0.5, -0.05) arc (0:-15:0.5) node[right, labelstyle, popmuted] {$-\theta$};

\node[labelstyle] at (3, 1.8) {\textbf{(a) Formation de l'image $B'$ au plan focal de $M_1$}};
\end{scope}

% --- FIGURE (b) : Miroir plan M2 ---
\begin{scope}[xshift=10cm]
% Axe incident horizontal Delta1, Axe emergent vertical Delta2
\draw[axis] (-1,0) -- (5,0) node[right, labelstyle] {$\Delta_1$};
\draw[axis] (3,-1) -- (3,5) node[above, labelstyle] {$\Delta_2$};
% Miroir plan M2 a 45 deg. Passe par O(3,0) ? Non, M2 coupe Delta1 en O.
% Place M2 avant F1. Disons F1 a x=3 (dans repere local). M2 a x=2.5.
% M2 : droite passant par (2.5, 0) angle 45 deg. Equation y = x - 2.5.
% Normale : pente -1. Angle 135 deg (vers haut-gauche) ou -45 deg (bas-droite).
% Normale vers l'incident (haut-gauche) : angle 135 deg.
\draw[mirror, thick] (1.5, -1) -- (4.5, 2) node[above right, labelstyle] {$M_2$};
% Point O intersection M2 / Delta1
\node[point, label=below:{$O$}] at (2.5, 0) {};
% Normale en O (pour rayon principal)
\draw[normal] (2.5, 0) -- (1, 1.5) node[above left, labelstyle, poppurple] {Normale};
% Rayons incidents (venant de fig a, converges vers B')
% B' etait en (3, -0.8) dans repere fig a (F1=3). 
% Ici F1 serait a x=3. M2 a x=2.5. B' a x=3 (apres M2).
% Les rayons arrivent sur M2 AVANT d'atteindre B'.
% Rayon Principal : vient de S1(0,0) angle -15 deg. Eq y = -0.268x.
%   Intersection avec M2 (y=x-2.5) : -0.268x = x - 2.5 -> 1.268x = 2.5 -> x=1.97. y=-0.53.
%   Point I1(1.97, -0.53).
% Rayon Marginal : vient de I_m(0.08, 1) vers B'(3, -0.8). Pente = (-1.8)/2.92 = -0.616.
%   Eq: y - 1 = -0.616(x - 0.08).
%   Intersection M2: x - 2.5 = 1 - 0.616x + 0.05 -> 1.616x = 3.55 -> x=2.20. y=-0.30.
%   Point I2(2.20, -0.30).

% On simplifie le dessin : On montre les deux rayons convergents coupant M2, et leurs reflechis.
% Rayon 1 (Principal) : Incident sur M2 en I1. Angle incidence ? 
%   Vecteur incident v_i = (1, -0.268) (direction vers droite-bas). Angle = -15 deg.
%   Normale N = (-1, 1) (vers haut-gauche). Angle = 135 deg.
%   Angle entre v_i (345 deg) et N (135 deg) = 210 deg -> 150 deg. i = 30 deg.
%   Reflechi : symetrique par rapport a N. Angle v_r = 135 + 30 = 165 deg (ou -195). 
%   Direction ~ (-0.966, 0.259). Monte vers gauche ? Non, doit aller vers HAUT (Delta2).
%   Erreur : Normale doit pointer vers le rayon incident. Rayon vient de gauche. M2 pente 45 (haut-droite). Normale vers haut-gauche (135 deg) est correcte.
%   v_i angle -15 deg (ou 345). N angle 135.
%   v_i est a "droite" de N ? Non. Produit vectoriel.
%   Mieux : Loi de reflection vectorielle. R = I - 2(I.N)N.
%   I = ({cos(-15)}, {sin(-15)}) = (0.966, -0.259).
%   N = (-cos45, sin45) = (-0.707, 0.707).
%   I.N = -0.683 - 0.183 = -0.866. cos(150) = -0.866. i = 150 deg ? Non, i est aigu. i = 30 deg.
%   R = I - 2(-0.866)N = I + 1.732 N = (0.966, -0.259) + 1.732(-0.707, 0.707) = (0.966 - 1.225, -0.259 + 1.225) = (-0.259, 0.966).
%   Angle R = arctan(0.966/-0.259) = 105 deg. (Vertical haut = 90 deg). Donc 15 deg par rapport a la verticale.
%   Le rayon reflechi part vers le HAUT (Delta2) en formant 15 deg a GAUCHE de la verticale.
%   Comme B' etait en BAS (y=-0.8), l'objet B'' pour l'oculaire sera en GAUCHE de Delta2.
%   Coherent : Image inversee.

% Dessin simplifie pour la figure (b) :
% On montre M2, la normale, le rayon principal incident et reflechi.
% Et le rayon marginal incident et reflechi.
% On note B'' l'intersection des reflechis (symetrique de B' par rapport a M2).

\draw[ray] (0.5, -0.13) -- (1.97, -0.53) node[midway, below, labelstyle] {Rayon 1 (Principal)};
\draw[ray] (1.97, -0.53) -- (1.97, 1.5) node[midway, right, labelstyle] {Rayon 1'} ; % Vers le haut (approx vertical)
% Correction direction reflechi : Angle 105 deg (15 deg a gauche de la verticale).
% Depuis I1(1.97, -0.53), direction (-0.26, 0.97).
\draw[ray] (1.97, -0.53) -- (1.0, 2.5) node[midway, left, labelstyle] {Rayon 1'};

% Rayon 2 (Marginal)
% Incident depuis gauche-haut vers I2 sur M2.
% Dans fig a, marginal part de (0.08, 1) vers B'(3, -0.8).
% Intersection M2 (y=x-2.5) : calcule plus haut I2(2.20, -0.30).
\draw[ray] (0.5, 0.8) -- (2.20, -0.30) node[midway, above, labelstyle] {Rayon 2 (Marginal)};
% Reflechi : Doit passer par B'' (symetrique de B' par M2).
% B'(3, -0.8). M2: y=x-2.5 -> x-y-2.5=0.
% Symetrique B'' : x' = x - 2a(ax+by+c)/(a^2+b^2). a=1, b=-1, c=-2.5.
% d = (3 - (-0.8) - 2.5)/2 = 1.3/2 = 0.65.
% x'' = 3 - 2*1*0.65 = 1.7. y'' = -0.8 - 2*(-1)*0.65 = 0.5.
% B''(1.7, 0.5). (A gauche de M2, au dessus de Delta1).
% Rayon reflechi 2 : de I2(2.20, -0.30) vers B''(1.7, 0.5).
\draw[ray] (2.20, -0.30) -- (1.7, 0.5) -- (1.0, 1.5) node[midway, above, labelstyle] {Rayon 2'};

% Objet B'' pour oculaire
\node[point, label=left:{$B''$}] at (1.7, 0.5) {};
\draw[rayconstr] (1.7, 0.5) -- (1.7, 2.5) node[midway, right, labelstyle] {Objet oculaire (sur $\Delta_2$)}; % Delta2 est x=3? Non.
% Delta2 est l'axe optique de l'oculaire. Il passe par B'' et est perpendiculaire a M2 ? 
% Non, l'oculaire est centre sur le faisceau reflechi. L'axe Delta2 passe par le centre de l'oculaire et est parallele au rayon principal reflechi (qui passe par le centre de l'oculaire).
% Le rayon principal reflechi (1') passe par I1(1.97, -0.53) et a pour direction (-0.26, 0.97).
% L'oculaire est place sur ce rayon. Son centre O_oc est sur ce rayon.
% L'axe Delta2 est la droite (I1, direction rayon principal reflechi).
% B'' est sur cet axe (par construction symetrie).
% Donc Delta2 passe par I1 et B''.
% Equation Delta2 : passe par (1.97, -0.53) et (1.7, 0.5). Pente = (1.03)/(-0.27) = -3.8. Angle ~ 105 deg. Correct.

% On ne trace pas l'oculaire ici (figure c).
\node[labelstyle] at (2.5, 3) {\textbf{(b) Déviation $90^\circ$ sur $M_2$ et formation de $B''$}};
\end{scope}

% --- FIGURE (c) : Oculaire ---
\begin{scope}[xshift=20cm]
% Axe Delta2 vertical (approx). On redresse pour clarite : Axe horizontal.
% On "deplie" localement l'oculaire.
% Objet B'' au foyer objet F_oc. Lentille L_oc. Rayons emergents paralleles.
\draw[axis] (-1,0) -- (7,0) node[right, labelstyle] {$\Delta_2$ (axe oculaire)};
% Lentille mince
\draw[thick, popblue] (2, -1.2) -- (2, 1.2) node[right, labelstyle] {$L_{oc}$};
\node[point, label=below:{$O_{oc}$}] at (2,0) {};
\node[point, label=below:{$F_{oc}$}] at (1,0) {}; % f' = 1cm a gauche
\node[point, label=below:{$F'_{oc}$}] at (3,0) {}; % f' = 1cm a droite
% Objet B'' au foyer objet
\node[point, label=above:{$B''$}] at (1, 0.6) {}; % Hauteur h'' = 0.6
% Rayons incidents (convergents vers B'')
% Rayon 1 (Principal) : passe par O_oc (centre lentille) -> non devié.
%   Doit venir de B''(1, 0.6) et passer par O_oc(2,0). Pente = -0.6.
\draw[ray] (0.5, 0.9) -- (1, 0.6) -- (2, 0) -- (3, -0.6) node[midway, below, labelstyle] {Rayon principal};
% Rayon 2 (Marginal) : Parallele a Delta2 (horizontal) -> Refracte par F'_oc.
%   Doit venir de B''(1, 0.6) horizontalement vers la lentille (x=2). Hauteur 0.6.
%   Refracte vers F'_oc(3,0). Pente = (0-0.6)/(3-2) = -0.6. Parallele au rayon 1 !
\draw[ray] (0.5, 0.6) -- (1, 0.6) -- (2, 0.6) -- (3, 0) -- (4, -0.6) node[midway, above, labelstyle] {Rayon // $\Delta_2$};
% Rayon 3 (Autre) : Depuis B'' vers haut lentille (2, 1.0) -> Refracte ?
%   Pour montrer parallele, deux rayons suffisent.
% Faisceau emergent : Rayons paralleles, pente -0.6. Angle theta' = arctan(0.6) = 31 deg.
% Theta etait 15 deg. Grossissement G = -theta'/theta = -31/15 = -2.06.
% Theorique G = -F/f' = -3/1 = -3. (Approximation petit angle : tan theta ~ theta).
% tan(15)=0.268. h' = -F tan theta = -0.8. h'' = 0.8 (symetrie M2 approx).
% theta' = h''/f' = 0.8/1 = 0.8 rad = 45 deg. 
% Mon dessin a h''=0.6, f'=1 -> theta'=31 deg. G=-2. 
% Je vais ajuster h'' pour correspondre a G=-3. h'' = 0.8. f'=1. theta'=38.6 deg.
% Theta=15 deg. G=-2.5. Proche.
% Je mettrai des valeurs coherentes dans le texte, le dessin est qualitatif.

\draw[ray, poporange!80!black] (2, -0.6) -- (4, -0.6) node[midway, below, labelstyle] {$\theta'$};
\draw[ray, poporange!80!black] (2, 0.6) -- (4, 0.6) node[midway, above, labelstyle] {Faisceau émergent parallèle};

\node[labelstyle] at (2, 1.8) {\textbf{(c) Émergence afocale par l'oculaire $L_{oc}$}};
\end{scope}
\end{tikzpicture}
\legende{Tracé géométrique complet du télescope de Newton en trois temps : (a) Formation de l'image $B'$ au plan focal du miroir primaire $M_1$ (rayon principal par $S_1$, rayon marginal par construction de la normale $C_1I$ et stigmatisme de Gauss vers $B'$) ; (b) Réflexion à $45^\circ$ sur le miroir plan $M_2$ (loi $i=r$ pour chaque rayon) formant l'objet réel $B''$ pour l'oculaire ; (c) Lentille oculaire $L_{oc}$ placée de sorte que $B''$ soit à son foyer objet $F_{oc}$, donnant un faisceau émergent parallèle (instrument afocal).}
\end{popfigure}

\courssub{La technique du schéma déplié (système afocal équivalent)}

Pour simplifier les calculs de grossissement et les tracés schématiques (notamment en examen où le temps est limité), on utilise le \textbf{schéma déplié}. On « déplie » le trajet optique en supprimant le miroir plan $M_2$ et en alignant l'axe du miroir primaire $\Delta_1$ avec l'axe de l'oculaire $\Delta_2$. Le système devient optiquement équivalent à une \textbf{lunette astronomique afocale de Kepler} (car les deux foyers se confondent) : un objectif (le miroir primaire, puissance $V_{obj} = 1/F$) et un oculaire (puissance $V_{oc} = 1/f'$) séparés par la distance $d = F + f'$.

\begin{propriete}[Équivalence optique : Télescope de Newton $\equiv$ Lunette de Kepler afocale]
Dans le schéma déplié :
\begin{itemize}
    \item L'objectif est le miroir primaire (focale $F$, diamètre $D$).
    \item L'oculaire est la lentille convergente (focale $f'$).
    \item La distance entre le sommet du miroir $S_1$ et le centre optique de l'oculaire $O_{oc}$ est $L = F + f'$.
    \item L'image intermédiaire $B'$ (réelle, inversée) est commune au foyer image de l'objectif et au foyer objet de l'oculaire ($F_1 \equiv F_{oc}$).
\end{itemize}
Les formules de grossissement et de puissance sont identiques à celles de la lunette afocale.
\end{propriete}

\begin{popfigure}
\begin{tikzpicture}[scale=0.9, >=Stealth]
\tikzset{
    lens/.style={thick, popblue},
    mirror/.style={thick, popblue, decorate, decoration={zigzag, amplitude=1.5pt, segment length=3pt}}, 
    ray/.style={poporange, thick, ->},
    rayconstr/.style={popgreen, thin, dashed},
    axis/.style={popdark, thin, dashed},
    point/.style={fill=popdark, circle, inner sep=1.5pt},
    labelstyle/.style={font=\small, popdark}
}
% Axes alignés horizontalement
\draw[axis] (-1,0) -- (11,0) node[right, labelstyle] {$\Delta$ (axe déplié)};

% Miroir primaire (représenté comme un objectif pour le schéma déplié)
% Zigzag pour miroir
\draw[mirror] (0, -1.2) -- (0, 1.2) node[above, labelstyle] {$M_1$ (Objectif)};
\node[point, label=below:{$S_1$}] at (0,0) {};
\node[point, label=below:{$F_1$}] at (4,0) {};
\draw[rayconstr] (4, -1.5) -- (4, 1.5) node[right, labelstyle] {Plan focal commun ($F_1 \equiv F_{oc}$)};

% Oculaire
\draw[lens] (5, -1) -- (5, 1) node[above, labelstyle] {$L_{oc}$ (Oculaire)};
\node[point, label=below:{$O_{oc}$}] at (5,0) {};
\node[point, label=below:{$F'_{oc}$}] at (6,0) {};

% Distances
\draw[<->, popmuted] (0, -1.8) -- (4, -1.8) node[midway, below, labelstyle] {$F$};
\draw[<->, popmuted] (4, -1.8) -- (5, -1.8) node[midway, below, labelstyle] {$f'$};
\draw[<->, popmuted] (0, -2.5) -- (5, -2.5) node[midway, below, labelstyle] {$L = F + f'$};

% Rayons pour un point B à l'infini (angle theta)
% Paramètres : F=4, f'=1. Theta = 10 deg (tan=0.176). Hauteur rayon principal à S1 = 0.
% Hauteur rayon marginal à S1 = 1 (D=2).
% Rayon principal (par S1) : Incident angle theta. Reflechi par M1 (objectif) vers F1.
%   En schema déplié, M1 se comporte comme une lentille convergente de focale F.
%   Rayon incident par centre optique (S1) -> non dévié ? Non, miroir : rayon par sommet -> reflechi symetrique.
%   Dans le schema déplié "lunette de Kepler", l'objectif est une lentille.
%   On trace comme une lunette : Rayon incident // axe -> F1. Rayon incident par centre O -> non dévié.
%   Ici "centre" de l'objectif miroir est S1.
%   Pour un objet à l'infini incliné (angle theta) :
%   1. Rayon principal (chef de faisceau) : Passe par le centre de l'objectif (S1). Non dévié (lentille) OU reflechi symetrique (miroir).
%      Dans le schema déplié standard (lunette), on considère l'objectif comme une lentille mince.
%      Le rayon principal passe par le centre optique S1 et n'est pas dévié.
%      Il arrive avec angle theta, sort avec angle theta. Il coupe le plan focal F1 en B' de hauteur h' = -F tan(theta).
%   2. Rayon marginal (parallele à l'axe incident) : Refracté par l'objectif vers F1.
%      Mais ici l'objet est incliné. Le rayon marginal n'est pas parallele à l'axe.
%      Methode standard lunette afocale :
%      - Faisceau incident parallèle incliné theta.
%      - Objectif (focale F) : forme image B' en son plan focal. Hauteur h' = -F * theta (petits angles).
%      - Oculaire (focale f') : place son foyer objet sur B'. Faisceau emergent parallèle incliné theta' = -h'/f' = (F/f') theta.
%      Grossissement G = theta'/theta = -F/f'.

% Dessin des rayons :
% Faisceau incident : deux rayons paralleles angle theta (10 deg).
% Hauteurs à S1 (x=0) : y1 = 0 (principal), y2 = 1 (marginal haut).
% Direction : (cos10, sin10) ~ (0.985, 0.174).
% Rayon 1 (Principal) : Passe par S1(0,0). Non dévié par objectif (schema lentille).
\draw[ray] (-1, -0.174) -- (0,0) -- (4, 0.696) -- (5, 0.696) -- (7, 0.696) node[midway, above, labelstyle] {$\theta$};
% Rayon 2 (Marginal) : Passe par (0, 1). Incident // au rayon 1.
\draw[ray] (-1, 1.174) -- (0, 1) -- (4, 0.696) -- (5, 0.696) -- (7, 0.696);
% Explication : Rayon 2 refracté par objectif vers F1(4,0) ? Non, vers B'(4, 0.696).
% B' est à x=4, y = h' = F * tan(theta) = 4 * 0.176 = 0.704. Proche de 0.696.
% Rayon 2 : Incident en (0,1). Doit passer par B'(4, 0.7).
% Pente = (0.7 - 1)/4 = -0.075. Angle = -4.3 deg.
% Rayon emergent : Depuis B'(4, 0.7) vers Oculaire (5,0). Distance 1 = f'.
% Rayon passe par B' et Ooc(5,0) ? Non, B' est au foyer objet. Rayon passant par Ooc n'est pas devié.
% Rayon de B' vers Ooc(5,0) : Pente = (0 - 0.7)/1 = -0.7. Angle = -35 deg.
% Rayon parallele à l'axe entrant dans oculaire (hauteur 0.7) -> sort par F'oc(6,0). Pente = -0.7. Parallele.
% Donc faisceau emergent angle theta' = -35 deg (tan = 0.7). Theta = 10 deg. G = -3.5. Theorique -4. Approximation tan.

% TikZ Drawing:
% Incident rays
\draw[ray] (-1, -0.174) -- (0,0) node[midway, below, labelstyle, sloped] {$\theta$};
\draw[ray] (-1, 1.174) -- (0, 1) node[midway, above, labelstyle, sloped] {$\theta$};
% Objective (Mirror M1) at x=0. 
% Ray 1 (through center S1) : undeviated.
\draw[ray] (0,0) -- (4, 0.7);
% Ray 2 (marginal) : refracted towards B'(4, 0.7).
\draw[ray] (0, 1) -- (4, 0.7);
% Image B'
\node[point, label=above:{$B'$}] at (4, 0.7) {};
% Ocular at x=5.
% Ray 1 continues through Ooc(5,0)? No, Ray 1 is at y=0.7 at x=4. At x=5, y = 0.7 + 0.7 = 1.4? No.
% Ray 1 goes from S1(0,0) to B'(4,0.7). Equation y = 0.175x.
% At x=5 (Ooc), y = 0.875.
% This ray goes through Ooc? No, Ooc is at (5,0). Ray 1 is at y=0.875.
% Principal ray for ocular: The ray that passes through Ooc.
% Where does the ray through Ooc come from? It comes from B' direction.
% Ray from B'(4, 0.7) to Ooc(5,0). Slope -0.7.
% This ray enters ocular at height 0. Emerges undeviated (slope -0.7).
% Another ray from B' parallel to axis (height 0.7) -> enters ocular at (5, 0.7) -> refracts through F'oc(6,0). Slope -0.7. Parallel.
% So emergent beam is parallel, slope -0.7. Angle theta' = atan(-0.7) ~ -35 deg.

% Draw emergent rays:
% Ray through Ooc (Principal ray for ocular):
\draw[ray] (4, 0.7) -- (5, 0) -- (7, -1.4) node[midway, below, labelstyle] {$\theta'$};
% Ray parallel to axis (Marginal for ocular):
\draw[ray] (4, 0.7) -- (5, 0.7) -- (6, 0) -- (7, -0.7);

% F'oc
\node[point, label=below:{$F'_{oc}$}] at (6,0) {};

\node[labelstyle] at (5, 2) {\textbf{Schéma déplié : Équivalent lunette de Kepler afocale}};
\end{tikzpicture}
\legende{Schéma déplié du télescope de Newton : équivalent à une lunette de Kepler afocale. L'axe optique est redressé. Le miroir primaire $M_1$ (traité comme un objectif de focale $F$) forme l'image $B'$ en son foyer $F_1$. L'oculaire $L_{oc}$ (focale $f'$) est placé de sorte que son foyer objet $F_{oc}$ coïncide avec $F_1$. Distance $S_1O_{oc} = F + f'$. Les rayons incidents parallèles (angle $\theta$) émergent parallèles (angle $\theta'$). Grossissement $G = -\frac{F}{f'}$.}
\end{popfigure}

\courssub{Application numérique : tracé à l'échelle imposée}

Au Probatoire, on te demande souvent de réaliser un tracé \textbf{à l'échelle} (ex. 1/10, 1/20, 1/50) sur papier millimétré. Voici la démarche complète sur un exemple réaliste.

\begin{exemplebox}[Tracé à l'échelle 1/20 d'un télescope de Newton]
\textbf{Données :} Miroir primaire $F = 800\,\text{mm}$ ; Oculaire $f' = 20\,\text{mm}$ ; Angle incident $\theta = 0,5^\circ$. \textbf{Échelle :} 1/20 (1 cm sur le papier = 20 mm réel).
\textbf{Calculs préliminaires (brouillon) :}
\begin{align*}
    \text{Longueur tube (approx)} & : L \approx F = 800\,\text{mm} \rightarrow 40\,\text{cm sur papier.} \\
    \text{Distance } S_1O_{oc} & = F + f' = 820\,\text{mm} \rightarrow 41\,\text{cm.} \\
    \text{Hauteur image } B' & : h' = F \tan\theta \approx 800 \times \tan(0,5^\circ) \approx 800 \times 0,00873 \approx 6,98\,\text{mm} \rightarrow 0,35\,\text{cm.} \\
    \text{Angle émergent } \theta' & : G = -F/f' = -40 \Rightarrow \theta' = 40 \times 0,5^\circ = 20^\circ.
\end{align*}
\textbf{Protocole de tracé (sur feuille A3 ou deux A4 joints) :}
\begin{enumerate}
    \item Tracer l'axe $\Delta$ horizontal (41 cm). Origine $S_1$ à gauche.
    \item Graduation : tous les 2 cm = 40 mm réel. Repérer $F_1$ à 40 cm (800 mm), $O_{oc}$ à 41 cm (820 mm), $F'_{oc}$ à 42 cm (840 mm).
    \item \textbf{Miroir $M_1$ :} Tracer un arc de cercle de rayon $R=1600\,\text{mm}=80\,\text{cm}$ (trop grand pour la feuille) $\rightarrow$ \textbf{Astuce :} On trace seulement la portion utile au sommet (approximation parabolique $\approx$ sphérique) ou on utilise la règle des tangentes. Au Probatoire, on représente souvent le miroir par son diamètre $D$ (ex. $D=150\,\text{mm} \to 7,5\,\text{cm}$) centré sur $S_1$, et on trace les normales aux points d'incidence.
    \item \textbf{Faisceau incident :} Depuis la gauche, tracer deux rayons parallèles formant $\theta=0,5^\circ$ avec $\Delta$. L'un passe par $S_1$ (hauteur 0), l'autre à hauteur $y = D/2 = 7,5\,\text{cm}$ (bord du miroir).
    \item \textbf{Réflexion $M_1$ :}
        \item Rayon par $S_1$ : réfléchi symétrique (angle $-0,5^\circ$).
        \item Rayon marginal (hauteur 7,5 cm, // $\Delta$) : réfléchi vers $F_1$ (à 40 cm). Tracer la droite joignant point d'incidence $(0, 7,5)$ à $F_1(40, 0)$.
    \item \textbf{Image $B'$ :} Intersection des deux rayons réfléchis. Vérifier qu'elle tombe sur la verticale $x=40\,\text{cm}$ (plan focal). Hauteur mesurée $\approx -0,35\,\text{cm}$ (inversée).
    \item \textbf{Miroir plan $M_2$ (schéma réel plié) :} Si tracé plié demandé : placer $M_2$ à 45° vers $x \approx 38\,\text{cm}$ (avant $F_1$). Reporter les rayons, tracer normales, angles $i=r=45^\circ \pm \theta/2$. Déviation $90^\circ$.
    \item \textbf{Oculaire (schéma déplié) :} Au point $O_{oc}$ (41 cm), tracer lentille. Vérifier $B'$ est à $f'=1\,\text{cm}$ à gauche de $O_{oc}$ (oui, $41-40=1\,\text{cm}$).
    \item \textbf{Émergence :} Rayon par $O_{oc}$ direct. Rayon // $\Delta$ (venant de $B'$) $\to$ réfracté par $F'_{oc}$ (42 cm). Les deux rayons émergents sont parallèles, pente $\tan\theta' \approx \tan 20^\circ \approx 0,36$.
\end{enumerate}
\textbf{Résultat visuel :} Un faisceau entrant « fin » ($0,5^\circ$) sort « large » ($20^\circ$), illustrant le grossissement $G=-40$.
\end{exemplebox}

\begin{attention}[Pièges classiques au Probatoire (Tracé)]
\begin{itemize}
    \item \textbf{Miroir plan APRÈS le foyer :} Si tu places $M_2$ après $F_1$, les rayons se croisent en $B'$ puis divergent. Le miroir plan renverrait un faisceau divergent $\rightarrow$ l'oculaire ne pourrait pas rendre le faisceau parallèle (image floue, mise au point impossible). \textbf{Règle : $M_2$ toujours avant $F_1$.}
    \item \textbf{Angles non mesurés sur $M_2$ :} Dessiner le miroir plan à « peu près 45° » et les rayons réfléchis « à peu près à 90° » est \textbf{sanctionné}. Tu dois : (1) tracer la normale (équerre), (2) mesurer $i$ au rapporteur, (3) tracer $r=i$ de l'autre côté.
    \item \textbf{Oculaire mal positionné :} Le centre optique $O_{oc}$ doit être à distance $f'$ de l'image $B''$ (ou $B'$ en schéma déplié). Si $O_{oc}$ est trop loin ou trop près, les rayons émergents ne sont pas parallèles (instrument non afocal, fatigue oculaire).
    \item \textbf{Sens des flèches :} Les flèches indiquent le sens de propagation de la lumière (Entrée $\to$ $M_1$ $\to$ $M_2$ $\to$ $L_{oc}$ $\to$ Sortie). Oublier les flèches ou les mettre en sens inverse = perte de points.
    \item \textbf{Échelle non respectée :} Si l'énoncé impose « Échelle 1/10 », tes mesures au rapporteur/règle sur la copie doivent correspondre. Ne pas écrire « non contractuel » si une échelle est demandée.
\end{itemize}
\end{attention}

\courssub{Exercice résumé : Tracé complet et détermination graphique du grossissement}

\begin{exoresolu}[Tracé et mesure de $G$ sur figure]
\textbf{Énoncé :} On considère un télescope de Newton de focale primaire $F = 600\,\text{mm}$ et d'oculaire $f' = 15\,\text{mm}$. Le diamètre du miroir est $D = 120\,\text{mm}$. Un astre est observé sous un angle $\theta = 1^\circ$. \textbf{Échelle imposée : 1/10.}
\begin{enumerate}
    \item Réaliser le tracé géométrique complet (schéma plié avec miroir plan) du faisceau issu du point $B$ de l'astre (haut du champ).
    \item Mesurer graphiquement l'angle émergent $\theta'$ et en déduire le grossissement $G_{\text{graph}}$.
    \item Comparer à la valeur théorique $G_{\text{th}}$.
\end{enumerate}

\tcblower
\textbf{Solution détaillée :}

\textbf{1. Préparation de l'échelle 1/10 :}
$1\,\text{cm} \leftrightarrow 10\,\text{mm}$.
$F = 60\,\text{cm}$ ; $f' = 1,5\,\text{cm}$ ; $D = 12\,\text{cm}$ ; $L \approx 61,5\,\text{cm}$ (nécessite grande feuille ou raccourci pédagogique : on trace seulement la zone utile près du foyer).
\textit{Astuce d'examen :} On ne trace pas 60 cm de tube vide. On trace le miroir $M_1$ au sommet $S_1$, on indique « 60 cm » sur l'axe, on place $F_1$ à 60 cm. On place $M_2$ à 58 cm (2 cm avant $F_1$). On trace l'oculaire au-dessus de $M_2$ à distance $f'=1,5\,\text{cm}$ de l'image $B''$.

\textbf{2. Tracé pas à pas (Description pour l'élève : tu le fais sur ta copie) :}
\begin{itemize}
    \item \textbf{Axe $\Delta_1$ horizontal.} $S_1$ origine. $F_1$ à 60 cm. $C_1$ à 120 cm (hors feuille).
    \item \textbf{Miroir $M_1$ :} Arc de cercle centre $C_1$ (ou portion droite de 12 cm centrée sur $S_1$ avec normales radiales).
    \item \textbf{Rayons incidents (angle $\theta=1^\circ$) :}
        \item Rayon 1 (principal) : passe par $S_1$.
        \item Rayon 2 (marginal haut) : parallèle à $\Delta_1$, hauteur $y = +6\,\text{cm}$ (demi-diamètre).
    \item \textbf{Réflexion sur $M_1$ :}
        \item Rayon 1 : réfléchi symétrique $\to$ angle $-1^\circ$ par rapport à $\Delta_1$.
        \item Rayon 2 : incident // $\Delta_1$ $\to$ réfléchi passant par $F_1$. Tracer droite $( \text{Point incidence}, F_1)$.
    \item \textbf{Image $B'$ :} Intersection des deux rayons réfléchis. Elle se trouve sur la verticale $x=60\,\text{cm}$. Mesure hauteur $h' \approx -1,05\,\text{cm}$ (théorique : $600 \tan 1^\circ \approx 10,47\,\text{mm} \to 1,05\,\text{cm}$). \textbf{Image inversée (hauteur négative).}
    \item \textbf{Miroir plan $M_2$ :} Placé à $x \approx 58\,\text{cm}$ (2 cm avant $F_1$), incliné $45^\circ$ (montant vers la droite).
        \item Tracer la normale en chaque point d'incidence (perpendiculaire à $M_2$, donc angle $-45^\circ$ par rapport à $\Delta_1$).
        \item Pour Rayon 1 (arrive avec angle $-1^\circ$ par rapport à $\Delta_1$, donc angle d'incidence $i_1 = 45^\circ - 1^\circ = 44^\circ$ par rapport à la normale). Rayon réfléchi : $r_1 = 44^\circ$ de l'autre côté de la normale $\to$ angle $90^\circ - 44^\circ = 46^\circ$ par rapport à $\Delta_1$ (axe vertical $\Delta_2$). \textit{Note : déviation totale $90^\circ$}.
        \item Pour Rayon 2 (arrive vers $F_1$, calculer son angle d'arrivée sur $M_2$ au rapporteur, appliquer $i=r$).
    \item \textbf{Objet $B''$ pour l'oculaire :} Intersection des rayons réfléchis par $M_2$ (ou symétrique de $B'$ par rapport à $M_2$).
    \item \textbf{Oculaire $L_{oc}$ :} Centre $O_{oc}$ sur $\Delta_2$ (vertical) tel que $O_{oc}B'' = f' = 1,5\,\text{cm}$ ($B''$ est au foyer objet).
    \item \textbf{Émergence :}
        \item Rayon passant par $O_{oc}$ : non dévié.
        \item Rayon parallèle à $\Delta_2$ (venant de $B''$) : réfracté vers $F'_{oc}$ (à 1,5 cm au-dessus de $O_{oc}$).
        \item Les deux rayons émergents sont parallèles. Mesurer leur angle $\theta'$ par rapport à $\Delta_2$.
\end{itemize}

\textbf{3. Mesure graphique et calcul :}
Sur le tracé soigné (rapporteur précis) :
$\theta'_{\text{mesuré}} \approx 40^\circ$ (légèrement moins à cause des approximations de tracé).
$G_{\text{graph}} = -\frac{\theta'}{\theta} \approx -\frac{40^\circ}{1^\circ} = -40$.
\textbf{Valeur théorique :} $G_{\text{th}} = -\frac{F}{f'} = -\frac{600}{15} = -40$.
\textbf{Conclusion :} L'accord est excellent ($< 2\%$ d'erreur de lecture). Le signe $-$ confirme l'image finale inversée par rapport à l'astre (deux inversions : $M_1$ et $L_{oc}$ ; $M_2$ ne change pas le sens haut/bas par rapport à l'axe plié, mais renverse gauche/droite).
\end{exoresolu}

\begin{aretenir}[Résumé des points clés pour le tracé]
\begin{itemize}
    \item \textbf{4 étapes :} (1) Faisceau // incident $\theta$ sur $M_1$ $\to$ (2) Convergence en $B'$ (plan focal $F_1$) via lois réflexion (rayon // $\to F_1$, rayon $S_1$ $\to$ symétrique) $\to$ (3) Miroir plan $M_2$ à $45^\circ$ \textbf{avant} $F_1$ : loi $i=r$ pour \textbf{chaque rayon} $\to$ déviation $90^\circ$, objet $B''$ pour oculaire $\to$ (4) Oculaire $L_{oc}$ : $B''$ au foyer objet $F_{oc}$ $\to$ faisceau émergent // angle $\theta'$.
    \item \textbf{Schéma déplié :} Aligner $M_1$ et $L_{oc}$ sur même axe, distance $F+f'$. Équivalent à lunette de Kepler afocale. Permet calculs rapides ($G=-F/f'$).
    \item \textbf{Échelle :} Calculer les cotes réduites \textbf{avant} de tracer. Choisir échelle pour tenir sur feuille (1/10 à 1/50).
    \item \textbf{Rigueur Probatoire :} Rapporteur obligatoire pour $M_2$ (normale, $i$, $r$). Flèches sens propagation. Légendes ($F_1, F_{oc}, B', \theta, \theta'$).
\end{itemize}
\end{aretenir}

\begin{savaistu}[Le plus grand télescope newtonien amateur au Cameroun ?]
De nombreux clubs d'astronomie au Cameroun (comme l'Association Camerounaise d'Astronomie - ACA) construisent des télescopes de type Dobson (monture altazimutale simple pour tube newtonien). Un miroir primaire de $300\,\text{mm}$ ($F \approx 1500\,\text{mm}$) permet d'observer les bandes de Jupiter, les anneaux de Saturne et quelques nébuleuses brillantes (M42 à Orion) depuis les collines de Yaoundé ou les plaines du Nord, loin de la pollution lumineuse des villes. Le tracé que tu maîtrises ici est exactement ce qui permet aux opticiens amateurs de calculer la position du porte-oculaire (focuseur) pour que l'image tombe pile au bon endroit, quel que soit l'oculaire utilisé (Plössl 25 mm, 10 mm, 6 mm...).
\end{savaistu}

\coursec{Grossissement et puissance}

Après avoir tracé la marche des rayons lumineux à travers le télescope de Newton et compris comment se forme l'image finale à l'infini, nous devons maintenant quantifier l'efficacité de l'instrument. \emph{De combien l'angle sous lequel nous voyons l'astre est-il agrandi ?} C'est la question du \cle{grossissement}. Nous verrons aussi comment caractériser l'oculaire par sa \cle{puissance}, une grandeur bien commode pour l'astronome praticien.

\courssub{Définition du grossissement d'un télescope afocal}

Rappelons la situation : le télescope est réglé en \emph{position afocale} (image finale à l'infini). L'œil, placé derrière l'oculaire, reçoit un faisceau de rayons parallèles. L'astre, situé à l'infini, est vu sous un certain angle par l'œil nu.

\begin{definition}[Grossissement d'un instrument d'optique afocal]
Le \textbf{grossissement} $G$ (ou $G_a$ pour grossissement angulaire) d'un télescope réglé à l'infini est le rapport de l'angle $\theta'$ sous lequel l'image finale est vue à travers l'instrument, à l'angle $\theta$ sous lequel l'objet (l'astre) est vu à l'œil nu :
\[
G = \frac{\theta'}{\theta}
\]
$G$ est un nombre \textbf{sans unité}. Si $G > 1$, l'instrument grossit ; si $G < 1$, il réduit (cas des jumelles de théâtre, non étudiées ici).
\end{definition}

\emph{Intuition :} Imaginez la Lune. À l'œil nu, elle semble petite (environ $\num{0,5}^\circ$). Dans le télescope, elle semble énorme, occupant une grande partie du champ de vision. Le grossissement dit exactement « combien de fois plus large » elle paraît.

\courssub{Établissement de la formule $G = \dfrac{F}{f'_{oc}}$ — Démonstration exigible au Probatoire}

C'est une démonstration classique, courte mais rigoureuse, qu'il faut savoir reproduire intégralement. Elle repose sur la géométrie de l'image intermédiaire $A'B'$ formée au foyer commun du miroir primaire et de l'oculaire.

\begin{popfigure}
\begin{tikzpicture}[scale=1.2, >=Stealth]
\pgfmathsetmacro{\F}{4.5}
\pgfmathsetmacro{\foc}{1.5}
\pgfmathsetmacro{\h}{0.8}
\pgfmathsetmacro{\theta}{atan(\h/\F)}
\pgfmathsetmacro{\thetap}{atan(\h/\foc)}

% Optical axis
\draw[thin, popmuted] (-0.5,0) -- (\F+\foc+0.5,0) node[right] {$\Delta$ (axe optique)};

% Primary Mirror at x=0
\draw[popblue, very thick, decorate, decoration={snake, amplitude=2pt, segment length=10pt}] (0, -\h*1.5) -- (0, \h*1.5) node[left] {Miroir primaire};

% Focal plane of mirror at x = F
\draw[popblue, thick] (\F,-\h*1.5) -- (\F,\h*1.5) node[above right] {Plan focal du miroir ($F$)};
\draw[dashed, popblue] (\F,0) node[below] {$F$} -- (\F,\h);

% Eyepiece lens at x = F + f'oc
\draw[poporange, thick] (\F+\foc,-\h*1.5) -- (\F+\foc,\h*1.5) node[above right] {Oculaire (foyer objet $F'_{oc}$)};
\draw[dashed, poporange] (\F+\foc,0) node[below] {$F'_{oc}$} -- (\F+\foc,\h);

% Intermediate Image A'B'
\draw[popdark, very thick] (\F,0) -- (\F,\h) node[midway, left] {$A'B'$};
\node[popdark] at (\F, -\h*0.3) {$A'$};
\node[popdark] at (\F, \h*1.1) {$B'$};

% Chief ray from mirror vertex (0,0) to B' (F, h) -> defines theta
\draw[popblue, thick, ->] (0,0) -- (\F,\h);
% Angle theta at mirror vertex
\draw[popblue, <->] (0.3,0) arc (0:\theta:0.3) node[midway, right=2pt, popblue] {$\theta$};

% Ray through eyepiece optical center O (F+foc, 0) to B' (F, h) and beyond
\draw[popgreen, thick, ->] (\F,\h) -- (\F+\foc,0) -- ++(\foc*1.5, -\h*1.5) node[right] {$\theta'$};
% Angle theta' at eyepiece (emergent ray down-right)
\draw[poporange, <->] (\F+\foc+0.3,0) arc (0:-\thetap:0.3) node[midway, below=2pt, poporange] {$\theta'$};

% Dimensions labels
\draw[<->, popmuted] (0, -\h*1.8) -- (\F, -\h*1.8) node[midway, below] {$F = \overline{SF}$};
\draw[<->, popmuted] (\F, -\h*1.8) -- (\F+\foc, -\h*1.8) node[midway, below] {$f'_{oc} = \overline{F'_{oc}A'}$};

% Eye symbol
\draw[popdark] (\F+\foc+2.5, -0.3) ellipse (0.3 and 0.15);
\draw[popdark] (\F+\foc+2.5, 0) circle (0.1);
\node at (\F+\foc+2.5, -0.7) {Œil};

\end{tikzpicture}
\legende{Schéma géométrique pour l'établissement de la formule du grossissement $G = F/f'_{oc}$. L'image intermédiaire $A'B'$ (hauteur $h$) est commune au plan focal du miroir (foyer image $F$) et au foyer objet de l'oculaire $F'_{oc}$. Le rayon principal passant par le sommet du miroir $S$ et le centre optique de l'oculaire $O$ définit les angles $\theta$ et $\theta'$.}
\end{popfigure}

\begin{propriete}[Formule du grossissement d'un télescope de Newton afocal]
Pour un télescope de Newton réglé à l'infini (position afocale), de focale objectif $F$ (focale du miroir primaire) et de focale oculaire $f'_{oc}$, le grossissement vaut :
\[
\boxed{G = \frac{F}{f'_{oc}}}
\]
\emph{Analyse dimensionnelle :} $[F] = \text{m}$, $[f'_{oc}] = \text{m}$ $\Rightarrow$ $[G] = 1$ (sans unité). Cohérent.
\end{propriete}

\begin{experience}[Démonstration guidée — À savoir reproduire]
\emph{Étape 1 : Angle au miroir (objet à l'infini).}
L'astre est à l'infini. Les rayons incidents sont parallèles et forment un angle $\theta$ avec l'axe optique. Ils convergent au foyer $F$ du miroir. L'image intermédiaire $A'B'$ (hauteur $h = A'B'$) se forme dans le plan focal du miroir. La relation géométrique (petits angles) donne :
\[
\tan\theta = \frac{h}{F} \quad \xrightarrow[\theta \ll 1]{\text{petits angles}} \quad \theta \approx \frac{h}{F} \qquad (1)
\]

\emph{Étape 2 : Angle à l'oculaire (image finale à l'infini).}
Le télescope est afocal : l'image finale est à l'infini. Cela impose que l'image intermédiaire $A'B'$ se trouve \emph{exactement} au foyer objet $F'_{oc}$ de l'oculaire. L'oculaire fonctionne comme une loupe : l'objet ($A'B'$) est à son foyer, l'image est à l'infini. Les rayons émergents sont parallèles et forment un angle $\theta'$ avec l'axe. Dans le triangle rectangle au centre optique de l'oculaire $O$ :
\[
\tan\theta' = \frac{h}{f'_{oc}} \quad \xrightarrow[\theta' \ll 1]{\text{petits angles}} \quad \theta' \approx \frac{h}{f'_{oc}} \qquad (2)
\]

\emph{Étape 3 : Rapport.}
En divisant (2) par (1), la hauteur $h = A'B'$ se simplifie :
\[
G = \frac{\theta'}{\theta} = \frac{h/f'_{oc}}{h/F} = \frac{F}{f'_{oc}}
\]
\emph{Fin de la démonstration.}
\end{experience}

\begin{attention}[Conditions d'application et pièges]
\begin{itemize}
\item \textbf{Position afocale obligatoire :} La formule $G = F/f'_{oc}$ n'est valable que si le télescope est réglé pour donner une image à l'infini (œil au repos, accommodation nulle). Si on fait la mise au point pour un œil myope ou hypermétrope (image à distance finie), la formule change légèrement.
\item \textbf{Unités cohérentes :} $F$ et $f'_{oc}$ doivent être exprimés dans la \textbf{même unité} (mm ou m). $G$ est sans unité.
\item \textbf{Signe du grossissement :} Par convention, pour un télescope de Newton, l'image finale est \textbf{renversée} (tête en bas) par rapport à l'objet. En optique signée (algébrique), le grossissement est $G = -F/f'_{oc}$. Au Probatoire, on donne souvent la \textbf{valeur absolue} $|G| = F/f'_{oc}$ et on précise « image renversée ». Vérifiez l'énoncé.
\item \textbf{Ne pas confondre $G$ et $\gamma$ :} Le \textbf{grossissement $G$} est un rapport d'\textbf{angles} (sans unité). Le \textbf{grandissement $\gamma$} (ou $g$) est un rapport de \textbf{tailles d'images réelles} (taille image / taille objet), utilisé pour les lentilles formant une image sur écran. $\gamma$ est sans unité aussi, mais son sens physique est différent. Confondre les deux est une erreur classique au Probatoire.
\end{itemize}
\end{attention}

\courssub{Conséquences pratiques : choix de l'oculaire}

La formule $G = F/f'_{oc}$ dicte le choix de l'instrument et des accessoires.

\begin{aretenir}[Règles pratiques]
\begin{itemize}
\item \textbf{$G$ augmente avec $F$ :} Un miroir à longue focale (ex: $F = \SI{1200}{\mm}$) donne de forts grossissements même avec des oculaires standards.
\item \textbf{$G$ diminue avec $f'_{oc}$ :} Pour « zoomer », on met un oculaire de \textbf{courte focale} (ex: $\SI{4}{\mm}$, $\SI{6}{\mm}$). Pour un champ large (nébuleuses), on met un oculaire de \textbf{longue focale} (ex: $\SI{25}{\mm}$, $\SI{32}{\mm}$).
\item \textbf{Choix de l'oculaire pour un $G$ visé :} $f'_{oc} = \dfrac{F}{G}$.
\end{itemize}
\end{aretenir}

\begin{exemplebox}[Exemple chiffré : Télescope de Newton $F = \SI{900}{\mm}$]
Un astronome amateur possède un télescope de Newton de focale $F = \SI{900}{\mm}$ (soit $\SI{0,90}{\m}$). Il dispose de trois oculaires : $f'_{oc1} = \SI{25}{\mm}$, $f'_{oc2} = \SI{10}{\mm}$, $f'_{oc3} = \SI{4}{\mm}$.
Calculons les grossissements :
\begin{align*}
G_1 &= \frac{900}{25} = 36 \quad \text{(faible grossissement, champ large, idéal pour la Voie Lactée)} \\
G_2 &= \frac{900}{10} = 90 \quad \text{(grossissement moyen, bon pour les planètes)} \\
G_3 &= \frac{900}{4} = 225 \quad \text{(fort grossissement, détails planétaires, mais image plus sombre)}
\end{align*}
Observation de la Lune : $\theta_{\text{Lune}} \approx \num{0,5}^\circ = \frac{\pi}{360}\,\text{rad} \approx \num{8,73e-3}\,\text{rad}$.
Avec l'oculaire $\SI{10}{\mm}$ ($G=90$) : $\theta' = G \times \theta \approx 90 \times \num{0,5}^\circ = \num{45}^\circ$. La Lune remplit presque tout le champ apparent de l'oculaire (souvent $\num{50}^\circ$ à $\num{60}^\circ$).
\end{exemplebox}

\courssub{La puissance d'un oculaire et d'un télescope}

En optique géométrique, on caractérise souvent un système par sa \textbf{puissance} (ou \cle{vergences}), mesurée en \textbf{dioptries ($\delta$)}.

\begin{definition}[Puissance d'un oculaire / Puissance de grossissement]
La \textbf{puissance} $P$ d'un instrument d'optique afocal (ou de son oculaire seul) est définie comme le quotient de l'angle $\theta'$ (en radians) sous lequel l'image finale est vue, par la taille linéaire $h = A'B'$ de l'objet intermédiaire (l'image réelle au foyer) :
\[
P = \frac{\theta'}{h}
\]
L'unité est la \textbf{dioptrie ($\delta$)} : $1\,\delta = 1\,\text{m}^{-1}$.
\end{definition}

Pour l'oculaire utilisé comme loupe (objet à son foyer $f'_{oc}$, image à l'infini), on a vu que $\theta' \approx h/f'_{oc}$. Donc :
\[
\boxed{P = \frac{1}{f'_{oc}}}
\]
avec $f'_{oc}$ en \textbf{mètres} pour obtenir $P$ en dioptries.

\begin{propriete}[Lien entre Puissance $P$ et Grossissement $G$]
En combinant $G = F/f'_{oc}$ et $P = 1/f'_{oc}$ (avec $F$ en mètres), on obtient la relation fondamentale :
\[
\boxed{P = \frac{G}{F} \quad \text{ou} \quad G = P \times F}
\]
Cette formule est très utile : connaissant la focale du télescope $F$ (fixe) et la puissance de l'oculaire $P$ (gravée sur le fût, ex: $\SI{25}{\delta}$ pour $f'_{oc}=\SI{40}{\mm}$), on trouve $G$ instantanément.
\end{propriete}

\begin{exemplebox}[Puissance et grossissement]
Reprenons le télescope $F = \SI{900}{\mm} = \SI{0,90}{\m}$.
\begin{itemize}
\item Oculaire $f'_{oc} = \SI{10}{\mm} = \SI{0,010}{\m}$ $\Rightarrow$ $P = 1/0,010 = \SI{100}{\per\meter}$ (soit $\SI{100}{\delta}$).
\item Grossissement $G = P \times F = 100 \times 0,90 = 90$. (Retrouve $900/10 = 90$).
\end{itemize}
Sur le fût de l'oculaire, on lit souvent : « $\SI{10}{\mm}$ » (focale) ou « $\SI{100}{\delta}$ » (puissance). Les deux sont équivalents.
\end{exemplebox}

\courssub{Limites pratiques du grossissement : la barre des $2 \times D$}

On pourrait croire qu'il suffit de prendre un oculaire de focale très courte (ex: $\SI{2}{\mm}$) pour obtenir un grossissement énorme ($G=450$ pour $F=\SI{900}{\mm}$). La physique impose des limites.

\begin{savaistu}[Pourquoi ne pas grossir indéfiniment ? — Ouverture et Diffraction]
Le miroir primaire a un diamètre $D$ (l'\cle{ouverture}). C'est lui qui collecte la lumière.
\begin{enumerate}
\item \textbf{Luminosité :} La lumière collectée est proportionnelle à $D^2$. Si on grossit $G$ fois, l'image s'étale sur une surface angulaire $G^2$ fois plus grande. La luminance (brillance surfacique) chute comme $1/G^2$. Au-delà d'un certain $G$, l'image devient trop sombre pour voir des détails.
\item \textbf{Résolution (Diffraction) :} L'œil a une résolution angulaire d'environ $1'$ (minute d'arc) = $1/60^\circ \approx \num{3e-4}\,\text{rad}$. Le télescope a une résolution théorique (critère de Rayleigh) $\theta_R \approx 1,22 \lambda / D$. Grossir au-delà de ce qui permet à l'œil de résoudre $\theta_R$ ne montre que du flou (« grossissement vide »).
\end{enumerate}
Une règle empirique universelle en astronomie amateur donne le \textbf{grossissement maximal utile} :
\[
\boxed{G_{\text{max}} \approx 2 \times D \quad (D \text{ en millimètres})}
\]
\end{savaistu}

\begin{exemplebox}[Application : Télescope $D = \SI{150}{\mm}$, $F = \SI{900}{\mm}$]
\begin{itemize}
\item $G_{\text{max}} \approx 2 \times 150 = 300\times$.
\item Oculaire minimal requis : $f'_{oc,\text{min}} = F / G_{\text{max}} = 900 / 300 = \SI{3}{\mm}$.
\item Un oculaire de $\SI{2}{\mm}$ ($G=450$) donnera une image plus grande mais floue et sombre : du « grossissement vide ».
\item Grossissement minimal utile : $\approx D/6 \approx 25\times$ (pupille de sortie $\approx \SI{6}{\mm}$, pupille de l'œil dilatée). En dessous, on perd de la lumière (pupille de sortie $> \SI{7}{\mm}$).
\end{itemize}
\end{exemplebox}

\courssub{Accessoires : La lentille de Barlow (Hors programme, culture générale)}

Pour éviter d'acheter une multitude d'oculaires à courtes focales (souvent inconfortables : faible dégagement oculaire), on utilise une \textbf{lentille de Barlow}. C'est un système divergent (lentille concave) placé avant le foyer du miroir, qui \emph{allonge artificiellement la focale effective du télescope}.

\begin{methode}[Effet d'une Barlow $\times 2$]
\begin{enumerate}
\item On insère la Barlow dans le porte-oculaire, puis l'oculaire dans la Barlow.
\item La focale effective du télescope devient $F_{\text{eff}} = k \times F$ ($k=2$ pour une Barlow $\times 2$, $k=3$ pour $\times 3$, etc.).
\item Le grossissement devient $G_{\text{Barlow}} = k \times G_{\text{sans Barlow}}$.
\item \emph{Exemple :} Oculaire $\SI{10}{\mm}$ sur $F=\SI{900}{\mm}$ $\to G=90$. Avec Barlow $\times 2$ $\to G=180$ (équivalent oculaire $\SI{5}{\mm}$ mais avec le confort de l'oculaire $\SI{10}{\mm}$).
\end{enumerate}
\end{methode}

\begin{attention}[Barlow et qualité]
Une Barlow bon marché dégrade l'image (aberrations chromatiques, perte de contraste). Une bonne Barlow (apochromatique, 3-4 lentilles) est un investissement judicieux. Au Probatoire, on ne demande pas de calculer l'effet d'une Barlow, mais savoir qu'elle « multiplie la focale » est un plus.
\end{attention}

\courssub{Synthèse visuelle : $G$ en fonction de $f'_{oc}$}

\begin{popfigure}
\begin{tikzpicture}
\begin{axis}[
    width=\linewidth,
    height=0.55\linewidth,
    xlabel={Focale de l'oculaire $f'_{oc}$ (mm)},
    ylabel={Grossissement $G$},
    title={Grossissement $G = 900 / f'_{oc}$ pour $F = \SI{900}{\mm}$},
    domain=2:40,
    samples=100,
    xmin=0, xmax=42,
    ymin=0, ymax=500,
    grid=both,
    minor tick num=1,
    major grid style={popline},
    minor grid style={popline!50},
    axis lines=middle,
    enlargelimits={abs=0.5},
    clip=false
]
\addplot[popcurve, thick, domain=2.5:40] {900/x};
% Markers for typical eyepieces
\addplot[only marks, mark=*, mark size=2.5, poporange] coordinates {
    (4, 225)
    (6, 150)
    (10, 90)
    (15, 60)
    (25, 36)
    (32, 28)
};
% Dashed lines for Gmax example D=150mm -> Gmax=300
\addplot[dashed, popred, thick] coordinates {(0,300) (3,300)};
\addplot[dashed, popred, thick] coordinates {(3,0) (3,300)};
\node[popred, anchor=south west] at (3.2, 300) {$G_{\text{max}} \approx 300$ ($D=\SI{150}{\mm}$)};
% Labels for markers
\node[poporange, anchor=south, font=\small] at (4, 235) {$\SI{4}{\mm}$};
\node[poporange, anchor=south, font=\small] at (10, 100) {$\SI{10}{\mm}$};
\node[poporange, anchor=south, font=\small] at (25, 46) {$\SI{25}{\mm}$};
\end{axis}
\end{tikzpicture}
\legende{Évolution hyperbolique du grossissement $G$ en fonction de la focale de l'oculaire $f'_{oc}$ pour un télescope de focale $F = \SI{900}{\mm}$. La zone au-dessus de $G_{\text{max}}$ (ligne pointillée rouge) correspond au « grossissement vide ».}
\end{popfigure}

\begin{exoresolu}[Exercice type Probatoire : Choix d'oculaire et puissance]
\textbf{Énoncé :} Un télescope de Newton a un miroir primaire de diamètre $D = \SI{200}{\mm}$ et de focale $F = \SI{1000}{\mm}$. L'astronome souhaite observer Jupiter.
\begin{enumerate}
\item Calculer le grossissement maximal utile $G_{\text{max}}$.
\item Déduire la focale minimale $f'_{oc,\text{min}}$ de l'oculaire à utiliser.
\item L'astronome possède un oculaire gravé « $\SI{12,5}{\mm}$ — $\SI{80}{\delta}$ ». Vérifier la cohérence de ce marquage.
\item Calculer le grossissement obtenu avec cet oculaire. Est-il adapté pour Jupiter ?
\item Quelle puissance $P$ (en $\delta$) doit avoir un oculaire pour atteindre exactement $G_{\text{max}}$ ?
\end{enumerate}

\tcblower
\textbf{Solution détaillée :}
\begin{enumerate}
\item \textbf{Grossissement maximal utile :} $G_{\text{max}} \approx 2 \times D = 2 \times 200 = \boxed{400\times}$.
\item \textbf{Focale minimale de l'oculaire :} $f'_{oc,\text{min}} = \dfrac{F}{G_{\text{max}}} = \dfrac{\SI{1000}{\mm}}{400} = \boxed{\SI{2,5}{\mm}}$. Un oculaire de focale inférieure donnerait du grossissement vide.
\item \textbf{Vérification du marquage :} $f'_{oc} = \SI{12,5}{\mm} = \SI{0,0125}{\m}$. Puissance théorique $P = \dfrac{1}{f'_{oc}} = \dfrac{1}{0,0125} = \SI{80}{\per\meter}$ (soit $\SI{80}{\delta}$). Le marquage « $\SI{80}{\delta}$ » est \textbf{cohérent}.
\item \textbf{Grossissement avec cet oculaire :} $G = \dfrac{F}{f'_{oc}} = \dfrac{1000}{12,5} = \boxed{80\times}$. (Ou $G = P \times F = 80 \times 1,0 = 80$).
\textit{Appréciation :} $80\times$ est un grossissement \textbf{faible à modéré} pour une planète. Jupiter (diamètre apparent $\approx 40''$ à $50''$) sera vue sous $\approx 3600'' = \num{1}^\circ$. C'est petit pour voir les bandes nuageuses et la Grande Tache Rouge. Il faudrait un grossissement plus fort (ex: $150\times$ à $250\times$), donc un oculaire de $f'_{oc} \approx \SI{4}{\mm}$ à $\SI{6,5}{\mm}$.
\item \textbf{Puissance pour $G_{\text{max}} = 400$ :} $P = \dfrac{G_{\text{max}}}{F} = \dfrac{400}{1,000} = \boxed{\SI{400}{\per\meter}}$ (soit $\SI{400}{\delta}$). (Correspond à $f'_{oc} = 1/400 = \SI{0,0025}{\m} = \SI{2,5}{\mm}$).
\end{enumerate}
\end{exoresolu}

\begin{aretenir}[Bilan : Grossissement et Puissance]
\begin{center}
\begin{tabularx}{\linewidth}{|l|X|X|}\hline
\textbf{Notion} & \textbf{Formule / Définition} & \textbf{Unité / Remarque} \\ \hline
Grossissement $G$ & $G = \dfrac{\theta'}{\theta} = \dfrac{F}{f'_{oc}}$ & Sans unité. Image renversée ($G<0$ algébrique). \\ \hline
Puissance oculaire $P$ & $P = \dfrac{\theta'}{h} = \dfrac{1}{f'_{oc}}$ & Dioptrie ($\delta = \text{m}^{-1}$). $f'_{oc}$ en \textbf{mètres}. \\ \hline
Lien $G$ et $P$ & $G = P \times F$ ($F$ en m) & Très pratique : $G = P(\delta) \times F(\text{m})$. \\ \hline
Grossissement max utile & $G_{\text{max}} \approx 2 \times D(\text{mm})$ & Au-delà : image sombre et floue (diffraction). \\ \hline
Choix oculaire & $f'_{oc} = \dfrac{F}{G_{\text{visé}}}$ & Barlow $\times k$ $\Rightarrow$ $F_{\text{eff}} = kF$. \\ \hline
\end{tabularx}
\end{center}
\end{aretenir}

\courssub{Exercices d'entraînement}

\begin{exercice}[Entraînement 1 : Calculs directs]
Un télescope de Newton a $F = \SI{1200}{\mm}$ et $D = \SI{150}{\mm}$.
\begin{enumerate}
\item Calculer $G_{\text{max}}$.
\item Quel grossissement avec un oculaire $f'_{oc} = \SI{20}{\mm}$ ? Et avec $f'_{oc} = \SI{6}{\mm}$ ?
\item Quelle puissance (en $\delta$) pour un oculaire donnant $G = 200$ ?
\end{enumerate}
\end{exercice}

\begin{exercice}[Entraînement 2 : Analyse critique]
Un vendeur vante un télescope « $600\times$ » sur l'emballage. Le tube indique $D = \SI{60}{\mm}$, $F = \SI{700}{\mm}$. L'oculaire fourni est un $f'_{oc} = \SI{4}{\mm}$ (avec une Barlow $\times 3$ bas de gamme).
\begin{enumerate}
\item Calculer le grossissement réel avec l'oculaire seul, puis avec la Barlow.
\item Comparer à $G_{\text{max}}$. Quel est votre avis sur la mention « $600\times$ » ?
\end{enumerate}
\end{exercice}

\begin{corrige}[Exercice 1]
\begin{enumerate}
\item $G_{\text{max}} \approx 2 \times 150 = 300\times$.
\item $G(\SI{20}{\mm}) = 1200/20 = 60\times$. $G(\SI{6}{\mm}) = 1200/6 = 200\times$. Les deux sont $< 300\times$, utilisables.
\item $G = P \times F \Rightarrow P = G/F = 200 / 1,2 \approx \SI{167}{\delta}$. (Vérif : $f'_{oc} = 1/167 \approx \SI{0,006}{\m} = \SI{6}{\mm}$).
\end{enumerate}
\end{corrige}

\begin{corrige}[Exercice 2]
\begin{enumerate}
\item Oculaire seul : $G = 700/4 = 175\times$. Avec Barlow $\times 3$ : $G = 3 \times 175 = 525\times$.
\item $G_{\text{max}} \approx 2 \times 60 = 120\times$. Le grossissement $525\times$ est \textbf{bien au-delà} de $G_{\text{max}}$. L'image sera très sombre, floue, et vide de détails. La mention « $600\times$ » est un \textbf{argument marketing trompeur} (grossissement vide), classique sur les instruments d'entrée de gamme « jouets ». Un bon instrument de $\SI{60}{\mm}$ s'utilise entre $20\times$ et $120\times$.
\end{enumerate}
\end{corrige}

\coursec{Mise au point du télescope de Newton}

\courssub{Transition depuis le grossissement}
Dans la section précédente, nous avons établi que le grossissement angulaire $G = -\dfrac{f'_1}{f'_2}$ dépend uniquement des distances focales algébriques du miroir primaire et de l'oculaire. Pourtant, connaître cette formule ne suffit pas pour observer : il faut encore que l'instrument soit \textbf{réglé} pour que l'image finale soit nette et confortable. C'est tout l'objet de la \textbf{mise au point} (ou \emph{focalisation}), étape indispensable avant toute observation, qu'elle soit terrestre ou astronomique. Rappelons que le télescope de Newton fonctionne en \textbf{configuration afocale} : l'objet est à l'infini, l'image finale doit l'être aussi pour un œil emmétrope (normal) sans accommodation. Toute dérive par rapport à cette condition afocale se traduit par une image floue.

\courssub{Principe optique de la mise au point}

\begin{definition}[Condition de mise au point afocale]
Pour un télescope de Newton utilisé en vision normale (œil emmétrope, accommodation nulle), la mise au point est réalisée lorsque le \textbf{plan focal image du miroir primaire} (déporté latéralement par le miroir plan secondaire) \textbf{coïncide exactement avec le plan focal objet de l'oculaire}.
\end{definition}

\par
Rappelons la géométrie : le miroir primaire concave (focale algébrique $f'_1 < 0$) forme une image réelle, inversée et réduite d'un objet lointain (étoile, planète) dans son plan focal, situé à une distance $|f'_1|$ de son sommet. Le miroir plan secondaire, incliné à $45^\circ$, réfléchit ce faisceau convergent vers le côté du tube, sans modifier la position du plan focal par rapport au sommet du primaire (seule la direction change). L'oculaire (lentille convergente de focale $f'_2 > 0$) doit être placé de sorte que son premier foyer objet $F_2$ tombe \textbf{sur ce plan focal intermédiaire}. Alors, le faisceau émergent de l'oculaire est parallèle (sortie à l'infini) : l'observateur voit une image nette sans forcer son œil.

\begin{propriete}[Position de l'oculaire pour la mise au point à l'infini]
Soit $f'_1$ la distance focale algébrique du miroir primaire ($f'_1 < 0$) et $f'_2$ celle de l'oculaire ($f'_2 > 0$). Si l'on note $d$ la distance (toujours positive) entre le sommet du miroir primaire et le centre optique de l'oculaire \emph{mesurée le long du chemin optique plié}, la condition de mise au point sur un objet à l'infini s'écrit :
\[
d = |f'_1| + f'_2
\]
En pratique, $d$ est fixé par la construction du tube ; c'est l'oculaire qui se déplace de quelques millimètres dans son porte-oculaire (crémaillère) pour satisfaire cette égalité.
\end{propriete}

\begin{attention}[Piège : signe des distances focales]
Le miroir primaire a une focale \textbf{négative} (miroir convergent en convention algébrique cartésienne où la lumière incidente vient de la gauche). L'oculaire a une focale \textbf{positive}. La formule $d = |f'_1| + f'_2$ utilise les \textbf{valeurs absolues} (longueurs physiques). Ne pas écrire $d = f'_1 + f'_2$ qui donnerait une distance absurde (négative ou trop courte).
\end{attention}

\courssub{Réglage pratique : la crémaillère de mise au point}

\begin{experience}[Observation de l'effet de la mise au point]
\textbf{Matériel :} Télescope de Newton sur monture stable, oculaire de focale moyenne (ex. $25\,\text{mm}$), cible terrestre lointaine (pylône, sommet d'arbre à $> 500\,\text{m}$) ou la Lune.
\textbf{Protocole :}
\begin{enumerate}
    \item Pointer l'instrument sur la cible à l'aide du chercheur (alignement grossier).
    \item Regarder dans l'oculaire. Tourner lentement la molette de mise au point (crémaillère) dans un sens, puis dans l'autre.
    \item Observer l'évolution de la netteté : il existe une position unique où l'image est \textbf{piquée} (détails nets, contours nets).
    \item Noter le sens de rotation pour \textbf{rentrer} l'oculaire (rapprocher l'oculaire du miroir primaire) et pour le \textbf{sortir}.
\end{enumerate}
\textbf{Interprétation :}
\begin{itemize}
    \item \textbf{Oculaire trop sorti} ($d > |f'_1| + f'_2$) : le plan focal intermédiaire se forme \emph{avant} le foyer objet de l'oculaire. Le faisceau incident sur l'oculaire est déjà divergent. L'oculaire renvoie un faisceau \textbf{convergent} : l'image virtuelle se forme à une distance finie devant l'œil. L'œil doit accommoder (fatigue) ou l'image reste floue si l'accommodation est insuffisante.
    \item \textbf{Oculaire trop rentré} ($d < |f'_1| + f'_2$) : le plan focal intermédiaire se forme \emph{après} le foyer objet de l'oculaire. Le faisceau incident est encore convergent. L'oculaire renvoie un faisceau \textbf{divergent} : l'image virtuelle se forme derrière l'œil (impossible à voir net sans instrument additionnel).
    \item \textbf{Réglage parfait} ($d = |f'_1| + f'_2$) : faisceau incident convergent vers $F_2$, faisceau émergent \textbf{parallèle} (afocal). Image nette, confortable, à l'infini.
\end{itemize}
\end{experience}

\begin{popfigure}
\begin{tikzpicture}[scale=0.9, >=Stealth]
    % Couleurs
    \colorlet{raycolor}{popblue}
    \colorlet{mirrorcolor}{popdark}
    \colorlet{lenscolor}{poppurple}
    \colorlet{eyecolor}{popgreen}
    
    % Styles
    \tikzset{
        mirror/.style={draw=mirrorcolor, very thick, fill=mirrorcolor!20},
        lens/.style={draw=lenscolor, very thick, fill=lenscolor!10},
        ray/.style={raycolor, thick},
        eye/.style={draw=eyecolor, thick, fill=eyecolor!10},
        axis/.style={dashed, thin, popmuted},
    }

    % --- Cas (a) : Oculaire trop sorti (faisceau convergent sortant) ---
    \begin{scope}[xshift=0cm]
        % Miroir primaire (représenté par son plan focal pour simplifier)
        \draw[mirror] (-0.5, -1.5) -- (-0.5, 1.5) node[midway, left=2mm] {Miroir primaire};
        % Plan focal intermédiaire
        \draw[axis] (3, -1.5) -- (3, 1.5) node[above] {$F'_1$ (plan focal intermédiaire)};
        % Miroir plan (schématisé par déviation)
        \draw[mirror] (2.5, 0) -- (3.5, 1) node[midway, above right] {Miroir plan};
        % Oculaire TROP SORTI (à droite de F'_1 + f'_2)
        \draw[lens] (5.5, -1) -- (5.5, 1) node[midway, right=2mm] {Oculaire};
        \draw[axis] (5.5 - 1.5, -1.5) -- (5.5 - 1.5, 1.5) node[above] {$F_2$};
        
        % Rayons
        \draw[ray] (-2, 1) -- (3, 1); % Incident parallèle
        \draw[ray] (3, 1) -- (3.5, 0.5); % Vers miroir plan
        \draw[ray] (3.5, 0.5) -- (5.5, 0.5); % Vers oculaire (hauteur 0.5)
        % Après oculaire : divergent car objet (F'_1) est avant F_2 -> image virtuelle à gauche
        \draw[ray] (5.5, 0.5) -- (7, 0.8);
        \draw[ray] (5.5, 0.5) -- (7, 0.2);
        
        % Rayon 2 : par le centre optique oculaire
        \draw[ray] (3.5, 0) -- (5.5, 0) -- (7, 0);
        
        % Œil
        \draw[eye] (7.5, -0.5) ellipse (0.4 and 0.8) node[above=4mm] {Œil accommode};
        
        \node[below=0.5cm] at (2.5, -1.8) {\textbf{(a) Oculaire trop sorti : faisceau émergent convergent}};
        \node[below=1.2cm] at (2.5, -1.8) {Image floue (ou fatigue oculaire)};
    \end{scope}

    % --- Cas (b) : Réglage afocal correct (faisceau parallèle sortant) ---
    \begin{scope}[xshift=11cm]
        % Miroir primaire
        \draw[mirror] (-0.5, -1.5) -- (-0.5, 1.5) node[midway, left=2mm] {Miroir primaire};
        % Plan focal intermédiaire
        \draw[axis] (3, -1.5) -- (3, 1.5) node[above] {$F'_1 = F_2$};
        % Miroir plan
        \draw[mirror] (2.5, 0) -- (3.5, 1) node[midway, above right] {Miroir plan};
        % Oculaire BIEN PLACÉ
        \draw[lens] (5.5, -1) -- (5.5, 1) node[midway, right=2mm] {Oculaire};
        % F_2 confondu avec F'_1 (à x=3)
        
        % Rayons
        \draw[ray] (-2, 1) -- (3, 1);
        \draw[ray] (3, 1) -- (3.5, 0.5);
        \draw[ray] (3.5, 0.5) -- (5.5, 0.5);
        % Après oculaire : PARALLÈLE car objet au foyer F_2
        \draw[ray] (5.5, 0.5) -- (7.5, 0.5);
        
        \draw[ray] (3.5, 0) -- (5.5, 0) -- (7.5, 0);
        
        % Œil
        \draw[eye] (8, -0.5) ellipse (0.4 and 0.8) node[above=4mm] {Œil relâché};
        
        \node[below=0.5cm] at (2.5, -1.8) {\textbf{(b) Réglage afocal correct : faisceau émergent parallèle}};
        \node[below=1.2cm] at (2.5, -1.8) {Image nette, confortable (vision à l'infini)};
    \end{scope}
\end{tikzpicture}
\legende{Effet de la position de l'oculaire sur la nature du faisceau émergent. (a) Oculaire trop éloigné du primaire : le foyer objet $F_2$ est en amont du plan focal intermédiaire $F'_1$, le faisceau incident sur l'oculaire est divergent $\rightarrow$ faisceau émergent convergent (œil doit accommoder). (b) Réglage parfait : $F_2$ coïncide avec $F'_1$, faisceau émergent parallèle (afocal), vision sans accommodation.}
\end{popfigure}

\courssub{Mise au point sur objet rapproché vs objet à l'infini}

\begin{propriete}[Décalage de mise au point pour un objet à distance finie]
Si l'objet observé n'est pas à l'infini (ex. : oiseau à $50\,\text{m}$, sommet de montagne à $5\,\text{km}$), l'image intermédiaire formée par le miroir primaire se décale \textbf{vers l'arrière} (s'éloigne du miroir) par rapport à son plan focal $F'_1$.
\par
D'après la relation de Descartes pour le miroir : $\dfrac{1}{p} + \dfrac{1}{p'} = \dfrac{1}{f'_1}$.
Pour un objet à l'infini ($p \to -\infty$), $p' = f'_1$.
Pour un objet à distance finie $p < 0$ (objet réel devant le miroir), $|p'| > |f'_1|$ : l'image réelle $A'B'$ se forme \textbf{plus loin} que $F'_1$.
\par
Pour retrouver la condition afocale ($F_2$ coïncidant avec cette nouvelle image intermédiaire), il faut \textbf{rentrer l'oculaire} (diminuer $d$) de quelques millimètres.
\end{propriete}

\begin{exemplebox}[Calcul du décalage de mise au point pour une observation terrestre]
Un télescope de Newton a un miroir primaire de focale $f'_1 = -1000\,\text{mm}$. On observe un pylône situé à $D = 200\,\text{m} = 200\,000\,\text{mm}$.
\begin{enumerate}
    \item Position de l'image intermédiaire pour l'objet à l'infini : $p'_\infty = f'_1 = -1000\,\text{mm}$ (distance $1000\,\text{mm}$ devant le miroir).
    \item Position pour l'objet à $D$ : $p = -200\,000\,\text{mm}$.
    \[
    \frac{1}{p'} = \frac{1}{f'_1} - \frac{1}{p} = \frac{1}{-1000} - \frac{1}{-200\,000} = -0,001 + 0,000005 = -0,000995\,\text{mm}^{-1}
    \]
    \[
    p' \approx -1005,03\,\text{mm}
    \]
    \item L'image intermédiaire recule de $\Delta p' \approx 5\,\text{mm}$ par rapport au plan focal.
    \item Il faut donc \textbf{rentrer l'oculaire d'environ $5\,\text{mm}$} pour compenser.
\end{enumerate}
\emph{Remarque :} Pour la Lune ($D \approx 384\,000\,\text{km}$), $\Delta p' \approx 0,0026\,\text{mm}$ : négligeable, on considère la mise au point « à l'infini ».
\end{exemplebox}

\courssub{La collimation : alignement des axes optiques}

Au-delà de la mise au point (translation de l'oculaire), la qualité de l'image exige un \textbf{alignement angulaire parfait} des trois éléments : miroir primaire, miroir plan secondaire, oculaire. C'est la \textbf{collimation}.

\begin{definition}[Collimation d'un télescope de Newton]
La collimation consiste à ajuster les orientations relatives des miroirs pour que :
\begin{enumerate}
    \item L'axe optique du miroir primaire vise le centre du miroir secondaire.
    \item Le miroir secondaire (plan, incliné à $45^\circ$) renvoie l'axe optique exactement vers le centre du porte-oculaire (axe de l'oculaire).
    \item Le centre du miroir secondaire est centré dans le tube (pas de vignettage asymétrique).
\end{enumerate}
Un télescope mal collimaté donne des images d'étoiles déformées (asymétriques, en « comète » ou « queue d'hirondelle ») même avec une mise au point parfaite.
\end{definition}

\begin{popfigure}
\begin{tikzpicture}[scale=0.85, >=Stealth]
    % Couleurs
    \colorlet{tubecolor}{popdark}
    \colorlet{primarycolor}{popblue}
    \colorlet{secondarycolor}{poporange}
    \colorlet{ocularcolor}{poppurple}
    \colorlet{axiscolor}{popgreen}
    
    \tikzset{
        tube/.style={draw=tubecolor, very thick},
        mirror/.style={draw=primarycolor, very thick, fill=primarycolor!10},
        secmirror/.style={draw=secondarycolor, very thick, fill=secondarycolor!10},
        ocular/.style={draw=ocularcolor, very thick, fill=ocularcolor!10},
        axis/.style={axiscolor, thick, ->},
        dasaxis/.style={axiscolor, dashed, thick},
        good/.style={axiscolor, thick},
        bad/.style={poppink, thick, dashed},
    }

    % Tube (carré pour simplicité)
    \draw[tube] (0,0) rectangle (10, 6);
    % Miroir primaire (bas)
    \draw[mirror] (1, 0.2) -- (9, 0.2) node[midway, below=2mm] {Miroir primaire (concave)};
    % Axe primaire
    \draw[good] (5, 0.2) -- (5, 3) node[midway, right] {Axe primaire};
    
    % Miroir secondaire (centré, 45°)
    \draw[secmirror] (3.5, 4.5) -- (6.5, 4.5) node[midway, above=2mm] {Miroir plan (45°)};
    % Axe réfléchi (horizontal vers oculaire)
    \draw[good] (5, 4.5) -- (5, 5.5) node[midway, right] {Axe oculaire};
    
    % Oculaire (côté droit)
    \draw[ocular] (10, 4) -- (10, 5) node[midway, right=2mm] {Oculaire};
    
    % Indications d'alignement
    \node[align=center, font=\small] at (11.5, 3) {
        \textbf{Collimation correcte} \\
        Axes alignés \\
        Image stellaire : disque d'Airy net
    };
    
    % --- Mauvais alignement (décalé à droite) ---
    \begin{scope}[xshift=14cm]
        \draw[tube] (0,0) rectangle (10, 6);
        \draw[mirror] (1, 0.2) -- (9, 0.2) node[midway, below=2mm] {Miroir primaire};
        % Axe primaire décentré
        \draw[bad] (4, 0.2) -- (4, 3) node[midway, left] {Axe primaire décentré};
        
        \draw[secmirror] (3.5, 4.5) -- (6.5, 4.5) node[midway, above=2mm] {Miroir plan};
        % Axe réfléchi dévié
        \draw[bad] (5, 4.5) -- (5, 5.5) node[midway, right] {Axe oculaire};
        
        \draw[ocular] (10, 4) -- (10, 5) node[midway, right=2mm] {Oculaire};
        
        % Rayon "comatique" schématique
        \draw[poppink, thick, decorate, decoration={snake, amplitude=1mm, segment length=4mm}] (4, 3) -- (5, 4.5) -- (5, 5.5);
        
        \node[align=center, font=\small] at (11.5, 3) {
            \textbf{Mauvaise collimation} \\
            Axes désalignés \\
            Image stellaire : \textbf{comète} (coma)
        };
    \end{scope}
\end{tikzpicture}
\legende{Géométrie de la collimation. Gauche : alignement parfait des trois axes (primaire $\rightarrow$ secondaire $\rightarrow$ oculaire). Droite : décentrage du miroir primaire (vis de bascule mal réglées) ou inclinaison du secondaire $\rightarrow$ aberration de coma (étoiles en forme de comètes).}
\end{popfigure}

\begin{methode}[Procédure de collimation (méthode « Cheshire » ou laser, simplifiée)]
\begin{enumerate}
    \item \textbf{Centrage du secondaire :} Regarder par le porte-oculaire (sans oculaire). Ajuster la vis centrale et les trois vis de maintien de la « araignée » (support du secondaire) pour que le miroir secondaire apparaisse \textbf{parfaitement centré} dans le tube et que le reflet du miroir primaire remplisse le secondaire.
    \item \textbf{Alignement du secondaire (tilt) :} Toujours sans oculaire, agir sur les trois vis de bascule du secondaire pour que le reflet du porte-oculaire (ou de l'œil) apparaisse \textbf{centré} dans le reflet du miroir primaire (vu dans le secondaire).
    \item \textbf{Alignement du primaire (tilt) :} Agir sur les trois vis de bascule au fond du tube (souvent plus grosses, à l'arrière du miroir) pour que le reflet du secondaire (et du porte-oculaire) revienne \textbf{exactement au centre} du miroir primaire. L'image de l'œil (ou du laser) doit revenir sur lui-même.
    \item \textbf{Vérification à l'étoile :} Pointer une étoile brillante (Polaris, Véga). Défocaliser légèrement (sortir l'oculaire). L'anneau de diffraction (disque d'Airy élargi) doit être \textbf{concentrique}. Si l'ombre du secondaire est décentrée dans l'anneau $\rightarrow$ retoucher le primaire. Si l'anneau est ovale $\rightarrow$ astigmatisme ou pincement du miroir.
\end{enumerate}
\end{methode}

\begin{attention}[Erreurs fréquentes en collimation]
\begin{itemize}
    \item \textbf{Confondre mise au point et collimation :} Tourner la molette de mise au point ne corrige \textbf{jamais} une image en « comète ». La mise au point agit sur la \emph{distance} (translation), la collimation sur les \emph{angles} (rotation).
    \item \textbf{Forcer les vis :} Les vis de collimation (surtout du primaire) sont sensibles. Tourner par \textbf{petits coups} ($1/8$ de tour), vérifier, recommencer. Bloquer les contre-écrous (si présents) \emph{après} réglage final.
    \item \textbf{Négliger le centrage du secondaire :} Un secondaire décentré vignette le champ de manière asymétrique et dégrade la résolution.
    \item \textbf{Collimer en plein jour sur le Soleil :} \textbf{INTERDIT}. Danger immédiat de cécité. Collimer sur une étoile (nuit) ou un reflet de lampe lointain (jour, sans Soleil dans le champ).
\end{itemize}
\end{attention}

\courssub{Ordre des opérations en séance d'observation}

\begin{methode}[Mise au point complète en 4 étapes (routine d'observation)]
\begin{enumerate}
    \item \textbf{Viser avec le chercheur :} Pointez le tube grossièrement vers la cible. Regardez dans le chercheur (petite lunette à faible grossissement, champ large, réticule). Centrez l'objet sur le réticule. Verrouillez les axes de la monture (si équatoriale) ou serrez les freins (azimutale).
    \item \textbf{Placer l'oculaire de plus grande focale ($f'_2$ max) :} Ex. $32\,\text{mm}$ ou $25\,\text{mm}$. Cela donne le \textbf{plus faible grossissement} ($|G|$ min), le \textbf{plus grand champ réel}, la \textbf{plus grande pupille de sortie} (image plus lumineuse, tolérance de mise au point plus large). C'est le plus facile pour « accrocher » l'objet et faire la mise au point grossière.
    \item \textbf{Mise au point précise :} Tournez la molette de crémaillère jusqu'à obtenir une image \textbf{parfaitement piquée} (détails nets, absence de halo). Pour une étoile : le disque d'Airy doit être le plus petit et rond possible. Pour la Lune/planète : contours nets, détails fins visibles. \emph{Astuce :} Faites un aller-retour rapide autour du point net pour bien identifier le maximum de netteté.
    \item \textbf{Changer d'oculaire pour grossir (si besoin) :} Remplacez par un oculaire de focale plus courte (ex. $10\,\text{mm}$, $6\,\text{mm}$). \textbf{Refaites la mise au point} (léger ajustement, car le plan focal objet $F_2$ change de position selon $f'_2$). Vérifiez que la turbulence atmosphérique (seeing) le permet : grossissement utile max $\approx 2 \times D(\text{mm})$ (ex. $200\times$ pour $D=100\,\text{mm}$).
\end{enumerate}
\end{methode}

\begin{propriete}[Pourquoi commencer par la grande focale ?]
\begin{itemize}
    \item \textbf{Champ plus grand :} On retrouve l'objet plus facilement dans le champ de l'oculaire (le chercheur n'est pas parfaitement parallélisé avec le tube principal).
    \item \textbf{Profondeur de champ plus grande :} La tolérance sur la position de l'oculaire $\Delta d$ pour une netteté acceptable est proportionnelle à $(f'_2)^2$ (via le nombre d'ouverture $N = |f'_1|/D$ et la formule de profondeur de champ). Un $f'_2$ grand pardonne un réglage imparfait.
    \item \textbf{Pupille de sortie plus grande :} $D_{\text{sortie}} = D_{\text{entrée}} / |G| = D \cdot f'_2 / |f'_1|$. Plus lumineux, plus confortable pour l'œil.
\end{itemize}
\end{propriete}

\begin{attention}[Piège classique de l'examen / TP]
\begin{itemize}
    \item \textbf{Énoncé :} « L'élève règle la mise au point avec un oculaire $f'_2 = 5\,\text{mm}$, n'y arrive pas, puis met un $25\,\text{mm}$ et y arrive. Expliquer. »
    \item \textbf{Réponse attendue :} Avec $f'_2 = 5\,\text{mm}$, grossissement fort ($|G| = 200$ pour $f'_1=-1000$), champ très petit (objet facile à rater), profondeur de champ très faible (réglage hyper-sensible), pupille de sortie minuscule ($0,5\,\text{mm}$ pour $D=100\,\text{mm}$) : image sombre, œil mal centré, diffraction visible. Avec $f'_2 = 25\,\text{mm}$, situation inverse : facile, lumineux, tolérant. \textbf{Toujours commencer par la plus grande focale disponible.}
\end{itemize}
\end{attention}

\courssub{Exemple résolu : Mise au point et changement d'oculaire}

\begin{exoresolu}[Mise au point successive avec deux oculaires]
Un télescope de Newton a un miroir primaire de diamètre $D = 150\,\text{mm}$ et de focale $f'_1 = -1200\,\text{mm}$. Il est équipé de deux oculaires : $O_1$ ($f'_{2,1} = 25\,\text{mm}$) et $O_2$ ($f'_{2,2} = 6\,\text{mm}$). On observe la Lune (distance $\approx \infty$).

\textbf{1. Grossissements et pupilles de sortie.}
Le grossissement angulaire algébrique est $G = -\dfrac{f'_1}{f'_2}$. Comme $f'_1 < 0$ et $f'_2 > 0$, $G > 0$. L'image finale est toutefois \textbf{inversée} par rapport à l'objet (le miroir primaire inverse l'image, le miroir plan ne la redresse pas). On retient généralement la valeur absolue $|G| = \dfrac{|f'_1|}{f'_2}$.
\[
G_1 = -\frac{f'_1}{f'_{2,1}} = -\frac{-1200}{25} = +48 \quad (|G_1| = 48)
\]
\[
G_2 = -\frac{f'_1}{f'_{2,2}} = -\frac{-1200}{6} = +200 \quad (|G_2| = 200)
\]
Pupilles de sortie :
\[
D_{s1} = \frac{D}{|G_1|} = \frac{150}{48} = 3,1\,\text{mm} \quad \text{(confortable, $< 7\,\text{mm}$ pupille œil)}
\]
\[
D_{s2} = \frac{D}{|G_2|} = \frac{150}{200} = 0,75\,\text{mm} \quad \text{(petite, observation planétaire, sensible au centrage œil)}
\]

\textbf{2. Position théorique de l'oculaire pour mise au point à l'infini.}
Distance miroir primaire $\rightarrow$ oculaire (chemin plié) :
\[
d_1 = |f'_1| + f'_{2,1} = 1200 + 25 = 1225\,\text{mm}
\]
\[
d_2 = |f'_1| + f'_{2,2} = 1200 + 6 = 1206\,\text{mm}
\]

\textbf{3. Déplacement de la crémaillère lors du changement $O_1 \to O_2$.}
L'oculaire $O_2$ a une focale plus courte. Son foyer objet $F_2$ est \textbf{plus proche} de sa lentille. Pour que $F_2$ coïncide avec le plan focal fixe du primaire ($F'_1$), il faut \textbf{rentrer l'oculaire} (le rapprocher du miroir primaire) de :
\[
\Delta d = d_1 - d_2 = 1225 - 1206 = 19\,\text{mm}
\]
\emph{En pratique :} On fait la mise au point avec $O_1$ (étape 3 de la méthode). On change pour $O_2$. On \textbf{rentre} la crémaillère d'environ $19\,\text{mm}$ (souvent un tour complet ou deux selon le pas de vis) et on affine.

\textbf{4. Grossissement utile maximum ?}
$G_{\text{max utile}} \approx 2 \times D(\text{mm}) = 300\times$. $G_2 = 200\times$ est \textbf{acceptable} (si seeing bon). $G_1 = 48\times$ est \textbf{très confortable} (champ large, nébuleuses, amas).
\end{exoresolu}

\courssub{Bilan et savoir-faire à retenir}

\begin{aretenir}[Check-list « Mise au point Newton » — Probatoire]
\begin{itemize}
    \item \textbf{Principe :} Confondre plan focal primaire (déporté) et foyer objet oculaire $\rightarrow$ faisceau émergent parallèle (afocal).
    \item \textbf{Geste :} Tourner molette crémaillère. \textbf{Rentrer} oculaire pour objet rapproché (ou oculaire courte focale) ; \textbf{Sortir} pour objet à l'infini (ou oculaire longue focale).
    \item \textbf{Ordre :} Chercheur $\rightarrow$ Oculaire $f'_2$ \textbf{max} $\rightarrow$ Mise au point $\rightarrow$ Oculaire $f'_2$ plus court $\rightarrow$ Re-mise au point (rentrer).
    \item \textbf{Collimation :} Aligner axes (primaire $\rightarrow$ secondaire $\rightarrow$ oculaire). Vis de bascule primaire (fond tube) + vis secondaire (araignée). Vérification : étoile défocalisée = anneaux concentriques.
    \item \textbf{Erreur type :} Mise au point directe avec courte focale $\rightarrow$ échec (champ trop petit, sensibilité trop forte).
\end{itemize}
\end{aretenir}

\begin{savaistu}[Collimation active sur les géants professionnels]
Les très grands télescopes (VLT de l'ESO au Chili, $D=8,2\,\text{m}$ ; future ELT $D=39\,\text{m}$) n'utilisent pas de miroirs primaires monolithiques rigides (trop lourds, se déforment sous leur poids). Ils ont des miroirs \textbf{actifs} (minces, $< 20\,\text{cm}$ d'épaisseur) soutenus par des centaines d'actionneurs (vérins pneumatiques ou piezoélectriques) qui corrigent la forme en temps réel (optique active). La collimation (alignement secondaire/primaire) et la mise au point sont pilotées par ordinateur via des capteurs de front d'onde (analyseurs de Shack-Hartmann) : le télescope \textbf{se recollimate tout seul} toutes les minutes pour compenser la gravité, le vent, la température. L'amateur, lui, a ses trois petites vis... et sa patience !
\end{savaistu}

\begin{exercice}[Entraînement : Mise au point et collimation]
\begin{enumerate}
    \item Un télescope Newton $f'_1 = -800\,\text{mm}$ est mis au point sur la Lune avec un oculaire $f'_2 = 20\,\text{mm}$. L'observateur veut regarder un oiseau perché à $30\,\text{m}$. Dans quel sens tourner la molette (rentrer ou sortir l'oculaire) ? Justifier par le signe de $\Delta p'$.
    \item Lors de la vérification à l'étoile (Véga, décalée), l'anneau de diffraction défocalisé montre l'ombre du miroir secondaire décentrée vers le haut. Quelle vis du miroir primaire faut-il toucher (haut/bas) et dans quel sens (serrer/desserer) pour recentrer l'ombre ?
    \item Calculer le déplacement de l'oculaire nécessaire pour passer d'un oculaire $26\,\text{mm}$ à un $9\,\text{mm}$ sur un télescope de focale $1000\,\text{mm}$.
\end{enumerate}
\end{exercice}

\begin{corrige}[Exercice n]
\begin{enumerate}
    \item \textbf{Rentrer l'oculaire.} Objet à distance finie ($p = -30\,000\,\text{mm}$) $\rightarrow$ image intermédiaire $p'$ plus négative que $f'_1$ ($|p'| > 800\,\text{mm}$). L'image recule vers l'arrière. Pour la rattraper avec $F_2$, il faut rapprocher l'oculaire du miroir (diminuer $d$).
    \item \textbf{Vis du bas (ou vis du haut selon montage), serrer la vis du bas / desserrer la vis du haut.} L'ombre du secondaire est haute $\rightarrow$ l'axe primaire vise trop bas $\rightarrow$ il faut relever l'axe primaire $\rightarrow$ pousser le bas du miroir vers l'avant (serrer vis bas) ou tirer le haut (desserrer vis haut). \emph{Note :} Sur la plupart des Newton, 3 vis à 120° ; on agit sur la vis opposée à la direction du décentrage de l'ombre.
    \item $\Delta d = f'_{2,\text{old}} - f'_{2,\text{new}} = 26 - 9 = 17\,\text{mm}$. Il faut \textbf{rentrer l'oculaire de $17\,\text{mm}$}.
\end{enumerate}
\end{corrige}

\coursec{Intérêt du télescope par rapport aux instruments à lentilles}

\courssub{Rappel sur la lunette astronomique et ses limites}

Avant d'analyser les atouts du télescope de Newton, rappelons le principe de son principal concurrent historique : la \cle{lunette astronomique} (ou lunette de Kepler). Elle utilise un \cle{objectif} constitué d'une \textbf{lentille convergente} de grande focale $f'_o$ qui forme une image réelle, inversée et réduite d'un objet lointain (étoile, planète) dans son plan focal. Un second système de lentilles, l'\cle{oculaire}, joue le rôle d'une loupe pour observer cette image intermédiaire.

\begin{attention}[Limites fondamentales de la lentille objectif]
Si la lunette astronomique a permis les premières grandes découvertes (lunes de Jupiter, phases de Vénus), elle bute sur trois obstacles physiques majeurs dès qu'on cherche à augmenter son diamètre $D$ pour collecter plus de lumière :
\begin{enumerate}
    \item \textbf{Aberration chromatique :} L'indice de réfraction $n$ du verre dépend de la longueur d'onde $\lambda$ (dispersion). La focale $f' = \frac{R}{n-1}$ (pour une lentille mince sphérique) varie donc avec la couleur. Les rayons bleus sont plus déviés que les rayons rouges : l'image présente des franges colorées. Corriger ce défaut exige des objectifs \emph{achromatiques} (doublet courbe-flint/courbe-crown) coûteux et lourds.
    \item \textbf{Aberration sphérique :} Les lentilles à surfaces sphériques ne convergent pas tous les rayons parallèles en un point unique. La correction impose des surfaces asphériques, difficiles à polir avec précision pour de grands diamètres.
    \item \textbf{Déformation mécanique et coût :} Une lentille ne peut être maintenue que par son bord (monture). Sous son propre poids, une grande lentille se déforme (affaissement central), ruinant la qualité optique. De plus, le verre optique de grande pureté doit être homogène \emph{dans tout son volume} (pas de bulles, pas de stries), ce qui rend la fabrication de diamètres $> 1\,\text{m}$ quasi impossible et astronomiquement chère.
\end{enumerate}
\end{attention}

\begin{savaistu}[La grande lunette de l'Observatoire de Paris (1886)]
La plus grande lunette jamais construite en France a un objectif de $D = 1,20\,\text{m}$ (focale $23\,\text{m}$). Elle pèse plus de 2 tonnes. Au-delà, le verre casse sous son poids ou se déforme trop. C'est cette impasse technologique qui a poussé les astronomes vers les miroirs.
\end{savaistu}

\courssub{Les quatre avantages majeurs du miroir newtonien}

Le télescope de Newton remplace l'objectif lentille par un \cle{miroir concave} (sphérique pour les petits diamètres, \textbf{parabolique} pour les instruments de qualité). La réflexion obéit à la loi de Descartes ($i = r$) qui \textbf{ne dépend pas de la longueur d'onde}. C'est la clé de tous les avantages.

\begin{propriete}[Absence d'aberration chromatique — Propriété admise au Probatoire]
La réflexion sur un miroir (métallique ou diélectrique) est indépendante de l'indice de réfraction du milieu traversant. Toutes les longueurs d'onde incidentes selon une même direction sont réfléchies selon la même direction. \textbf{Le foyer est unique pour toutes les couleurs.} Il n'y a \emph{aucune dispersion} à la réflexion.
\end{propriete}

\begin{propriete}[Fabricabilité de très grands diamètres]
Un miroir est un bloc de verre (ou vitro-céramique, ex. Zerodur) dont \textbf{seule la surface avant} doit être polie avec une précision nanométrique. Le dos du miroir peut être \textbf{soutenu par une cellule active} (vérins pneumatiques, aimants) qui compense la gravité et les déformations thermiques. Le substrat n'a pas besoin d'être optiquement homogène en volume.
\begin{itemize}
    \item Cela permet des miroirs monolithiques jusqu'à $8,4\,\text{m}$ (VLT, Subaru, Gemini).
    \item Au-delà, on utilise des \textbf{miroirs segmentés} (hexagones de $1,4\,\text{m}$ assemblés) : Keck ($10\,\text{m}$), GranTeCan ($10,4\,\text{m}$), TMT ($30\,\text{m}$), ELT ($39\,\text{m}$).
\end{itemize}
\end{propriete}

\begin{propriete}[Luminosité et pouvoir collecteur proportionnels à $D^2$]
La quantité d'énergie lumineuse collectée par seconde (flux lumineux $\Phi$) est proportionnelle à la surface d'entrée $S = \frac{\pi D^2}{4}$.
\[
\frac{\Phi_{\text{téléscope}}}{\Phi_{\text{lunette}}} = \left(\frac{D_{\text{téléscope}}}{D_{\text{lunette}}}\right)^2
\]
Un télescope newtonien amateur de $D = 200\,\text{mm}$ collecte $\left(\frac{200}{60}\right)^2 \approx 11$ fois plus de lumière qu'une lunette de $60\,\text{mm}$. Il révèle des objets 11 fois plus faibles (magnitude limite supérieure d'environ $2,7$ magnitudes).
\end{propriete}

\begin{propriete}[Grossissement du télescope de Newton]
Comme pour la lunette astronomique, le grossissement $G$ (ou puissance) est le rapport de la focale de l'objectif (ici le miroir primaire) sur la focale de l'oculaire :
\[
G = \frac{f'_p}{f'_e}
\]
où $f'_p$ est la focale du miroir primaire (positive pour un miroir concave) et $f'_e$ la focale de l'oculaire (positive pour un oculaire convergent). L'image finale est inversée ($G < 0$ par convention algébrique, mais on retient la valeur absolue $|G|$).
\end{propriete}

\begin{exemplebox}[Comparaison chiffrée : Lunette $60\,\text{mm}$ vs Newton $200\,\text{mm}$]
\begin{itemize}
    \item \textbf{Surface collectrice :} $S_{60} = \pi \times 30^2 \approx 2\,827\,\text{mm}^2$ ; $S_{200} = \pi \times 100^2 \approx 31\,416\,\text{mm}^2$. Rapport $\approx 11,1$.
    \item \textbf{Magnitude limite visuelle (ciel noir, œil dilaté $7\,\text{mm}$) :} $m_{\text{lim}} \approx 2,7 + 5\log_{10}(D/\text{mm})$.
        \begin{align*}
            m_{\text{lim}}(60) &\approx 2,7 + 5\log_{10}(60) \approx 11,6 \\
            m_{\text{lim}}(200) &\approx 2,7 + 5\log_{10}(200) \approx 14,2
        \end{align*}
        Le télescope montre des étoiles $\approx 2,6$ magnitudes plus faibles (facteur $10^{2,6/2,5} \approx 11$ en flux).
    \item \textbf{Prix indicatif (neuf, entrée de gamme) :} Lunette achromatique $60\,\text{mm} \approx 150\,000\,\text{FCFA}$ ; Newton $200\,\text{mm}/1000\,\text{mm} \approx 250\,000\,\text{FCFA}$. Le télescope est \emph{moins cher au mm² de surface collectrice}.
\end{itemize}
\end{exemplebox}

\begin{propriete}[Pouvoir de résolution (pouvoir séparateur) — Complément utile au Probatoire]
La diffraction limite la capacité à distinguer deux sources ponctuelles proches. L'angle minimal résoluble (critère de Rayleigh) pour une pupille circulaire de diamètre $D$ éclairée en lumière de longueur d'onde $\lambda$ est :
\[
\theta_{\min} \approx 1,22\,\frac{\lambda}{D} \quad (\text{en radians})
\]
\begin{aretenir}[Formule à connaître]
$\theta_{\min} (\text{rad}) \approx 1,22\,\dfrac{\lambda}{D}$ \quad avec $\lambda$ et $D$ dans la même unité (mètres).
Plus $D$ est grand, plus $\theta_{\min}$ est petit : \textbf{le télescope distingue des détails plus fins}.
\end{aretenir}
\begin{exemplebox}[Résolution : Lunette $60\,\text{mm}$ vs Newton $200\,\text{mm}$ ($\lambda = 550\,\text{nm}$)]
\begin{align*}
\theta_{\min}(60) &\approx 1,22 \times \frac{550 \times 10^{-9}}{0,060} \approx 1,12 \times 10^{-5}\,\text{rad} \approx 2,3'' \\
\theta_{\min}(200) &\approx 1,22 \times \frac{550 \times 10^{-9}}{0,200} \approx 3,35 \times 10^{-6}\,\text{rad} \approx 0,69''
\end{align*}
Le télescope $200\,\text{mm}$ sépare théoriquement des détails \textbf{3,3 fois plus fins} (ex. cratères lunaires plus petits, séparation d'étoiles doubles plus serrées).
\end{exemplebox}
\end{propriete}

\begin{propriete}[Suppression de l'aberration sphérique par le miroir parabolique]
Un miroir sphérique souffre d'aberration sphérique (les rayons marginaux coupent l'axe plus près que les rayons paraxiaux). Un miroir \textbf{parabolique} de foyer $F$ a pour propriété géométrique : tout rayon incident parallèle à l'axe optique converge \textbf{exactement} en $F$, sans approximation de Gauss. C'est la forme idéale pour un télescope newtonien (rapport $f/D$ généralement $< 6$). La fabrication parabolique est standardisée (test de Foucault, interférométrie).
\end{propriete}

\courssub{Inconvénients à nuancer : une vision équilibrée}

\begin{attention}[Ne pas écrire : « Le télescope de Newton n'a aucun défaut »]
Tout instrument a des compromis. Pour une rédaction de Probatoire équilibrée, cite ces points :
\begin{enumerate}
    \item \textbf{Obstruction centrale :} Le miroir secondaire plan (ou convexe pour les Cassegrain) et son support (araignée) masquent une partie de l'entrée. Perte de lumière $\approx (d/D)^2$ (ex. $d=50\,\text{mm}, D=200\,\text{mm} \to 6,25\%$). Cela réduit légèrement le contraste (diffraction aux bords de l'obstacle : pics de diffraction en croix sur les étoiles brillantes).
    \item \textbf{Collimation (alignement) :} Les miroirs primaire et secondaire doivent être parfaitement alignés (axes confondus). Un choc, une dilatation thermique différentielle désalignent l'instrument. \textbf{La collimation est une maintenance régulière} (vis de réglage au dos du primaire, vis sur le porte-secondaire). Une lunette, une fois collimée en usine, ne bouge pratiquement plus.
    \item \textbf{Image renversée :} Comme la lunette astronomique, l'image finale est inversée (rotation $180^\circ$). Pas de redressement par prisme (perte de lumière, aberrations) en astronomie ; on s'habitue à la carte du ciel « tête en bas ».
    \item \textbf{Tube ouvert :} Courants d'air (turbulence interne), poussière, rosée sur le miroir primaire. Nécessite un pare-buée (chauffage résistance) et un nettoyage annuel délicat (miroir aluminisé fragile).
    \item \textbf{Coma (aberration hors axe) :} Pour les objets hors axe, les rayons forment une tache en « comète ». Défaut dominant pour les rapports $f/D$ courts ($< 5$), limitant le champ net.
\end{enumerate}
\end{attention}

\courssub{Bilan comparatif synthétique}

Le tableau suivant résume les critères de choix pour un élève ou un astronome amateur. Il est structuré pour une réponse de type « comparer » au Probatoire.

\begin{popfigure}
\begin{center}
\begin{tabularx}{\linewidth}{|l|X|X|X|}
\hline
\rowcolor{popblue}\color{white}\textbf{Critère} & \color{white}\textbf{Lunette astronomique (achromatique)} & \color{white}\textbf{Télescope de Newton (parabolique)} & \color{white}\textbf{Lunette de Galilée} \\
\hline
\rowcolor{popblueL!30}
Aberration chromatique & Résiduelle (doublet), franges violettes sur objets brillants & \textbf{Absente} (réflexion achromatique) & Forte (objectif simple) \\
\hline
\rowcolor{poporangeL!20}
Aberration sphérique & Corrigée par formes asphériques complexes & \textbf{Nulle} si miroir parabolique bien fait & Non corrigée (objectif simple) \\
\hline
\rowcolor{popblueL!30}
Diamètre max pratique (amateur) & $\sim 150\,\text{mm}$ (coût $\uparrow\uparrow$, poids) & $\sim 400\,\text{mm}$ accessible, $> 500\,\text{mm}$ possible & $< 60\,\text{mm}$ (champ étroit) \\
\hline
\rowcolor{poporangeL!20}
Luminosité (collecte de lumière) & Limitée par $D_{\max}$ & \textbf{Supérieure} ($D$ grand facile) & Très faible \\
\hline
\rowcolor{popblueL!30}
Pouvoir de résolution $\theta_{\min}$ & Modéré ($\propto 1/D$) & \textbf{Meilleur} ($D$ grand) & Faible \\
\hline
\rowcolor{poporangeL!20}
Coût au $\text{cm}^2$ de surface & Élevé (verre optique volume, doublet) & \textbf{Faible} (surface seule polie, substrat standard) & Faible mais $D$ petit \\
\hline
\rowcolor{popblueL!30}
Entretien / Robustesse & Faible (collimation usine stable) & \textbf{Collimation régulière nécessaire}, tube ouvert & Faible (alignement fixe) \\
\hline
\rowcolor{poporangeL!20}
Image finale & Renversée ($180^\circ$) & Renversée ($180^\circ$) & Droite (terrestre) \\
\hline
\rowcolor{popblueL!30}
Encombrement (longueur tube) & Long ($L \approx f'_o + f'_e$) & Court ($L \approx f'_p$, replié) & Court ($L \approx f'_o - |f'_e|$) \\
\hline
\end{tabularx}
\end{center}
\legende{Tableau comparatif : Lunette astronomique vs Télescope de Newton vs Lunette de Galilée. Les lignes bleues claires indiquent l'avantage au Newton.}
\end{popfigure}

\courssub{Méthode de rédaction pour une question « Comparer » au Probatoire}

\begin{methode}[Structurer une comparaison d'instruments d'optique]
Pour répondre à une consigne du type « Comparer le télescope de Newton et la lunette astronomique » ou « Justifier le choix d'un télescope pour l'observation du ciel profond », suis ce plan en 3 paragraphes distincts :
\begin{enumerate}
    \item \textbf{Qualité d'image (Aberrations) :} Cite la \textbf{chromatique} (absente au Newton, résiduelle à la lunette) et la \textbf{sphérique} (corrigée par parabole au Newton, doublet asphérique cher à la lunette). Conclusion : image plus pure, piquée, sans franges colorées pour le Newton.
    \item \textbf{Luminosité et Résolution (Diamètre) :} Rappelle $\Phi \propto D^2$ et $\theta_{\min} \approx 1,22\lambda/D$. Explique que le miroir se soutient par l'arrière $\Rightarrow$ $D$ grand accessible $\Rightarrow$ \textbf{voir plus faible} (nébuleuses, galaxies) et \textbf{voir plus fin} (détails planétaires, étoiles doubles).
    \item \textbf{Praticité et Coût :} Coût de fabrication $\propto D^3$ pour lentille (volume) vs $\propto D^2$ pour miroir (surface). Tube plus court (longueur $\approx f'$ vs $f'_o+f'_e$). Mentionne l'entretien (collimation) comme seul inconvénient réel, mais gérable.
\end{enumerate}
\textbf{Astuce :} Utilise des connecteurs logiques : \emph{« En revanche », « Par conséquent », « Ce qui permet », « Toutefois »}.
\end{methode}

\courssub{Applications modernes : l'héritage de Newton}

\begin{savaistu}[Du miroir de 3 cm (1668) aux géants segmentés de 39 m]
\begin{itemize}
    \item \textbf{1668 :} Newton polit un miroir d'alliage cuivre-étain (spéculos) de $D \approx 3\,\text{cm}$, $f' \approx 15\,\text{cm}$. Il observe les lunes de Jupiter.
    \item \textbf{1789 :} William Herschel construit un miroir de $1,2\,\text{m}$ (spéculos) — plus grand que toute lunette de l'époque.
    \item \textbf{1917 :} Télescope Hooker ($2,5\,\text{m}$, verre Pyrex) — Edwin Hubble découvre l'expansion de l'univers.
    \item \textbf{1990 :} \textbf{Hubble Space Telescope} — Miroir primaire $D = 2,4\,\text{m}$, $f/24$ (Ritchey-Chrétien, dérivé Cassegrain). Pas d'atmosphère $\Rightarrow$ résolution diffraction-limited ($0,05''$).
    \item \textbf{1998-- :} \textbf{VLT (ESO, Chili)} — 4 miroirs unitaires $D = 8,2\,\text{m}$ (monolithiques, actifs, $f/1,75$), combinables en interférométrie.
    \item \textbf{2021 :} \textbf{James Webb Space Telescope (JWST)} — Miroir primaire \textbf{segmenté} $D = 6,5\,\text{m}$ (18 hexagones de béryllium plaqués or), déployé en orbite L2. Héritier direct : \emph{réflexion, grand diamètre, segmentation, refroidissement passif}.
    \item \textbf{2028+ :} \textbf{ELT (ESO)} — $D = 39\,\text{m}$ (798 segments), $f/1,0$. Le plus grand œil du monde sur le ciel.
\end{itemize}
Au Cameroun, l'astronomie amateur se développe (clubs à Yaoundé, Douala, Ngaoundéré). Un Newton $150/750$ ou $200/1000$ sur monture Dobson est l'instrument « meilleur rapport performance/prix » pour débuter l'observation du ciel profond (nébuleuse d'Orion, amas d'Hercule, galaxies du Vierge) depuis nos latitudes tropicales.
\end{savaistu}

\begin{exoresolu}[Choix d'instrument pour un lycée camerounais]
\textbf{Énoncé :} Un club d'astronomie d'un lycée de Yaoundé dispose d'un budget de $500\,000\,\text{FCFA}$. Il hésite entre une lunette achromatique $80\,\text{mm}/900\,\text{mm}$ ($450\,000\,\text{FCFA}$) et un télescope Newton $150\,\text{mm}/750\,\text{mm}$ sur monture Dobson ($480\,000\,\text{FCFA}$). Le but est l'initiation au ciel profond (nébuleuses, galaxies) et planétaire. Conseille le club en justifiant par trois arguments physiques chiffrés.

\textbf{Solution détaillée :}
\begin{enumerate}
    \item \textbf{Luminosité (collecte de lumière) :} Rapport des surfaces $S_{150}/S_{80} = (150/80)^2 = (1,875)^2 \approx 3,5$. Le Newton collecte \textbf{3,5 fois plus de photons}. Pour des objets diffus faiblement lumineux (nébuleuses), c'est déterminant : la magnitude limite gagne $\Delta m = 2,5\log_{10}(3,5) \approx 1,3$ mag. Beaucoup d'objets invisibles dans la lunette deviennent détectables.
    \item \textbf{Résolution (détails planétaires) :} $\theta_{\min} \propto 1/D$. $\theta_{\min}(150)/\theta_{\min}(80) = 80/150 \approx 0,53$. Le Newton résout des détails \textbf{presque deux fois plus fins} sur Jupiter (bandes, Grande Tache) ou Saturne (division de Cassini nette).
    \item \textbf{Aberration chromatique :} La lunette $80\,\text{mm}$ achromatique montrera un liseré violet sur Jupiter et les étoiles brillantes (dispersion résiduelle). Le Newton, par réflexion, donne une image \textbf{strictement achromatique}, plus contrastée.
    \item \textbf{Coût/Encombrement :} Les deux rentrent dans le budget. Le Newton Dobson est plus court ($\sim 80\,\text{cm}$ vs $\sim 1\,\text{m}$), plus stable (centre de gravité bas), et ne nécessite pas de trépied équatorial coûteux.
\end{enumerate}
\textbf{Conclusion :} Le \textbf{télescope Newton $150\,\text{mm}$} est le choix pédagogique optimal. Il illustre parfaitement les avantages physiques du miroir (chromatisme nul, grand $D$ accessible) pour un budget maîtrisé.
\end{exoresolu}

\begin{exercice}[Entraînement Probatoire]
Un élève affirme : « Puisque le télescope de Newton a un miroir parabolique, il n'a aucune aberration et donne une image parfaite. »
\begin{enumerate}
    \item En te basant sur la physique des miroirs, explique pourquoi l'aberration sphérique est corrigée par la parabole.
    \item Cite \textbf{deux} défauts réels subsistant dans un télescope newtonien bien construit et bien collimaté, et explique leur origine physique.
    \item Pourquoi les grands télescopes professionnels (VLT, JWST, ELT) utilisent-ils tous des miroirs et non des lentilles ? Donne l'argument principal lié à la mécanique du solide.
\end{enumerate}
\end{exercice}

\begin{corrige}[Exercice Entraînement Probatoire]
\begin{enumerate}
    \item \textbf{Aberration sphérique et parabole :} Pour un miroir sphérique, les rayons marginaux (loin de l'axe) coupent l'axe plus près du miroir que les rayons paraxiaux (proches de l'axe) : c'est l'aberration sphérique. Un miroir \textbf{parabolique} a pour propriété géométrique que \textbf{tout rayon incident parallèle à l'axe optique est réfléchi vers le foyer unique $F$}, quelle que soit sa hauteur d'incidence. La parabole est la courbe isochrone pour la réflexion : chemin optique constant. Donc \textbf{aberration sphérique nulle pour les objets à l'infini sur l'axe}.
    \item \textbf{Défauts réels (hors axe et obstruction) :}
        \begin{itemize}
            \item \textbf{Coma (aberration de coma) :} Pour les objets \emph{hors axe} (champ large), les rayons incidents parallèles mais inclinés ne convergent plus en un point unique mais en une tache en forme de « comète » (queue dirigée vers l'axe ou l'extérieur). C'est l'aberration dominante du Newton à courte focale ($f/D < 5$). Elle limite le champ utile net.
            \item \textbf{Obstruction centrale et diffraction :} Le miroir secondaire plan (et son araignée à 3 ou 4 branches) masque le centre du primaire. Cela réduit le contraste (énergie diffusée dans les anneaux de diffraction) et crée des pics de diffraction (croix) sur les étoiles brillantes.
        \end{itemize}
    \item \textbf{Argument mécanique principal :} Une lentille ne peut être maintenue que par son \textbf{bord} (anneau). Son poids crée une flexion (flèche) proportionnelle à $D^4 / \text{épaisseur}^3$. Pour $D > 1\,\text{m}$, la déformation optique devient inacceptable sauf épaisseur massive (poids $\propto D^3$, coût $\propto D^3$, risque de rupture). Un miroir est maintenu par sa \textbf{face arrière} (cellule à vérins/aimants) sur toute sa surface. La flexion est contrôlée activement (optique active). Le substrat n'a pas besoin de transparence volumique. Cela autorise $D = 8\,\text{m}$ (monolithe) ou $39\,\text{m}$ (segmenté).
\end{enumerate}
\end{corrige}

\newpage

\coursec{Activité d'intégration}

\coursec{Le télescope de Newton}

\courssub{Objectifs de l'activité}

À l'issue de cette activité d'intégration, tu seras capable de :
\begin{itemize}
\item décrire et schématiser un télescope de Newton (miroir primaire concave, miroir secondaire plan, oculaire) en respectant les proportions réelles ;
\item tracer la marche d'un faisceau lumineux issu d'un objet à l'infini à travers l'instrument ;
\item calculer le grossissement $G$ et la puissance $P$ d'un télescope pour différents oculaires ;
\item choisir un oculaire en appliquant le critère du grossissement maximal utile ;
\item calculer la taille d'une image intermédiaire dans le plan focal du miroir ;
\item comparer objectivement un télescope à miroir et une lunette à lentilles, et argumenter un choix d'instrument.
\end{itemize}

\courssub{Situation}

Le club scientifique du Lycée de Ngaoundéré prépare la « Nuit des étoiles » organisée chaque année au mois de janvier, période où le ciel de l'Adamaoua est particulièrement pur et dégagé, loin de la pollution lumineuse des grandes villes. Le proviseur a débloqué un budget modeste et le club a pu acquérir le matériel suivant :

\begin{itemize}
\item un \cle{miroir primaire concave} sphérique de distance focale $F = \SI{1000}{mm}$ et de diamètre $D = \SI{130}{mm}$ ;
\item un \cle{miroir secondaire plan} de petit axe $d = \SI{35}{mm}$, destiné à renvoyer la lumière sur le côté du tube ;
\item trois \cle{oculaires} interchangeables assimilés à des lentilles convergentes de distances focales $f'_1 = \SI{25}{mm}$, $f'_2 = \SI{12,5}{mm}$ et $f'_3 = \SI{6}{mm}$ ;
\item un tube en PVC, une crémaillère de mise au point et une monture équatoriale fabriqués à l'atelier du lycée.
\end{itemize}

Le club souhaite observer en priorité :
\begin{itemize}
\item les \cle{cratères lunaires} : un cratère typique est vu depuis la Terre sous un diamètre apparent $\theta_c = \num{0,01}^{\circ}$ ;
\item la \cle{Lune} dans son ensemble, vue sous un diamètre apparent $\theta_L = \num{0,5}^{\circ}$ ;
\item les \cle{satellites galiléens de Jupiter} (Io, Europe, Ganymède, Callisto), objets ponctuels très peu lumineux.
\end{itemize}

Le proviseur, avant d'autoriser l'achat d'accessoires supplémentaires, pose la question suivante au club : « Pourquoi ne pas acheter plutôt une lunette astronomique de même prix ? ». Le club te charge, toi et ton groupe, de dimensionner le télescope, de vérifier ses performances et de rédiger l'argumentaire de réponse.

\courssub{Consignes}

\begin{enumerate}
\item \textbf{Schéma de l'instrument.} Réalise un schéma du télescope de Newton à l'échelle $1/10$ (coupe du tube) : miroir primaire concave, miroir secondaire plan incliné à $45^{\circ}$, oculaire placé sur le côté. Indique les distances focales réelles sur le schéma.
\item \textbf{Marche des rayons.} Sur ce schéma, trace la marche d'un faisceau parallèle issu d'un cratère lunaire, incliné de $\theta_c = \num{0,01}^{\circ}$ par rapport à l'axe optique, depuis son entrée dans le tube jusqu'à la sortie de l'oculaire (instrument réglé en vision à l'infini, c'est-à-dire en système \cle{afocal}). Où se forme l'image intermédiaire ?
\item \textbf{Grossissements et puissances.} Calcule le grossissement $G$ et la puissance $P$ de l'instrument pour chacun des trois oculaires. Présente les résultats dans un tableau.
\item \textbf{Choix de l'oculaire.} Le grossissement maximal utile d'un instrument est donné par $G_{\max} \approx 2D$ avec $D$ en millimètres. Vérifie si chaque oculaire respecte ce critère et détermine celui qui donne le grossissement le plus élevé tout en restant utilisable.
\item \textbf{Image de la Lune.} Calcule la taille $y'$ de l'image intermédiaire de la Lune ($\theta_L = \num{0,5}^{\circ}$) formée dans le plan focal du miroir primaire. Cette image tient-elle dans le champ de l'oculaire ?
\item \textbf{Argumentaire.} Rédige en dix lignes maximum un texte destiné au proviseur expliquant pourquoi le télescope de Newton est préférable à une lunette de même prix pour observer les satellites de Jupiter (pense à : luminosité, aberrations chromatiques, coût du diamètre, obstruction).
\end{enumerate}

\courssub{Ressources à mobiliser}

\begin{aretenir}[Formulaire utile]
\begin{itemize}
\item Télescope de Newton : un miroir concave (objectif) forme dans son \cle{plan focal} une image réelle de l'objet à l'infini ; un miroir plan à $45^{\circ}$ renvoie cette image sur le côté ; l'\cle{oculaire} joue le rôle de loupe et donne une image finale à l'infini (réglage afocal : plan focal du miroir confondu avec le plan focal objet de l'oculaire).
\item Taille de l'image intermédiaire : $y' = F \tan\theta \approx F\,\theta$ ($\theta$ en radians).
\item Grossissement : $G = \dfrac{\theta'}{\theta} = \dfrac{F}{f'_{\text{oculaire}}}$.
\item Puissance : $P = \dfrac{1}{f'_{\text{oculaire}}}$ (en dioptries si $f'$ est en mètres).
\item Grossissement maximal utile : $G_{\max} \approx 2D$ ($D$ en mm).
\item Avantages du miroir : pas d'aberrations chromatiques, grand diamètre à coût modéré, donc grande \cle{luminosité} (flux collecté proportionnel à $D^2$).
\end{itemize}
\end{aretenir}

\courssub{Méthode : Tracer la marche des rayons dans un télescope de Newton}

\begin{methode}[Tracé géométrique rigoureux]
\begin{enumerate}
\item \textbf{Repères.} Trace l'axe optique principal. Place le miroir primaire concave (représenté par un arc de cercle ou une flèche convergente) au fond du tube. Son foyer $F$ est à distance $F$ du sommet.
\item \textbf{Miroir secondaire.} Place le miroir plan à $45^{\circ}$ sur l'axe, légèrement en avant du foyer $F$ (distance $< F$). Trace la normale au miroir (à $45^{\circ}$ de l'axe).
\item \textbf{Faisceau incident.} Dessine au moins trois rayons parallèles inclinés de l'angle $\theta$ (exagéré pour la lisibilité) : un passant par le centre du miroir primaire, un par le bord supérieur, un par le bord inférieur.
\item \textbf{Réflexion sur le primaire.} Chaque rayon incident se réfléchit en passant par le foyer $F$ (loi de Descartes pour miroir concave : rayon parallèle $\to$ passe par $F$). Les trois rayons convergent vers un point image $I$ situé dans le plan focal (perpendiculaire à l'axe en $F$).
\item \textbf{Réflexion sur le secondaire.} Les rayons convergent vers $I$ mais rencontrent le miroir plan. Applique la loi de la réflexion (angle incidence = angle réflexion) ou la construction de l'image par un miroir plan : l'image $I'$ de $I$ est symétrique par rapport au plan du miroir. Le faisceau réfléchi converge vers $I'$.
\item \textbf{Oculaire.} Place l'oculaire (lentille convergente) de sorte que son foyer objet $F'_o$ coïncide avec $I'$ (réglage afocal). Trace la sortie : les rayons émergents sont parallèles entre eux, inclinés de l'angle $\theta'$ tel que $\tan\theta' = \overline{I'I_{\text{oc}}}/f'$.
\end{enumerate}
\end{methode}

\courssub{Attention : Pièges classiques et points de vigilance}

\begin{attention}
\begin{itemize}
\item \textbf{Unités du grossissement :} $G$ est un nombre pur (rapport de deux angles ou de deux longueurs focales en \emph{même unité}). Ne pas écrire « $\times$ » comme unité, c'est un facteur.
\item \textbf{Puissance en dioptries :} La formule $P = 1/f'$ impose $f'$ en \textbf{mètres}. Erreur fréquente : utiliser $f'$ en mm et trouver $P$ en $\delta$ (ex: $1/25 = 0,04$ au lieu de $40$).
\item \textbf{Grossissement maximal utile :} $G_{\max} \approx 2D$ ($D$ en mm) est un critère empirique (diffraction + turbulence). Le dépasser donne une image « vide » (floue, sombre) sans détail supplémentaire.
\item \textbf{Sens de l'image :} Le télescope de Newton donne une image \textbf{renversée} (rotation de $180^{\circ}$) par rapport à l'objet. Pas de redressement par le miroir plan (seul un renvoi coudé ou une lentille de Barlow redressent).
\item \textbf{Obstruction :} Le miroir secondaire (et son porte-miroir) obstrue une partie de l'ouverture ($d/D \approx 27\%$ ici), réduisant légèrement le contraste et la résolution théorique (critère de Rayleigh modifié), mais le gain en diamètre compense largement.
\end{itemize}
\end{attention}

\courssub{Le saviez-vous ? : Newton, l'optique et le Cameroun}

\begin{savaistu}
Isaac Newton (1643--1727) a inventé ce télescope vers 1668 pour contourner l'aberration chromatique des lentilles (dispersion du verre), qu'il jugeait incurable. Il a poli lui-même ses miroirs en alliage de cuivre-étain (spéulum). Aujourd'hui, les miroirs sont en verre (Pyrex, Zerodur) aluminisés. Au Cameroun, l'Observatoire astronomique de l'Université de Yaoundé I et les clubs d'astronomie (Ngaoundéré, Douala, Maroua) utilisent des télescopes Newton (Dobson ou sur monture équatoriale) pour l'observation du ciel profond et planétaire. Le site de l'Adamaoua (altitude $\sim \SI{1100}{m}$, faible humidité, faible pollution lumineuse) offre un \emph{seeing} (qualité de la turbulence atmosphérique) souvent excellent ($< 2''$), idéal pour exploiter la résolution d'un $D = \SI{130}{mm}$ ($\theta_R \approx \SI{1,06}{''}$).
\end{savaistu}

\courssub{Démarche et corrigé commenté}

\begin{exoresolu}[Conception du télescope du club d'astronomie]
Énoncé : Dimensionner le télescope de Newton du club ($F = \SI{1000}{mm}$, $D = \SI{130}{mm}$, oculaires de $25$, $12,5$ et $\SI{6}{mm}$), tracer la marche des rayons, calculer grossissements et puissances, choisir l'oculaire optimal, calculer l'image de la Lune et argumenter le choix de l'instrument.
\tcblower

\textbf{1. Schéma à l'échelle 1/10 et marche des rayons.}
À l'échelle $1/10$, la distance focale du miroir est $F_{\text{dessin}} = \SI{100}{mm} = \SI{10}{cm}$ et le diamètre $D_{\text{dessin}} = \SI{13}{mm}$. Le tube fait $\sim \SI{11}{cm}$ de long. Le miroir secondaire (petit axe $\SI{3,5}{mm}$) est placé à $\SI{9}{cm}$ du primaire. L'oculaire ($f' = \SI{25}{mm} \to \SI{2,5}{cm}$) est centré sur l'axe réfléchi, son foyer objet à $\SI{1}{cm}$ du miroir secondaire (image intermédiaire).

\begin{popfigure}
\begin{tikzpicture}[scale=1.0, >=Stealth, thick]
% Couleurs
\colorlet{mirrorPri}{popblue}
\colorlet{mirrorSec}{poporange}
\colorlet{ocul}{popgreen}
\colorlet{rayCol}{poppink}
\colorlet{axisCol}{popmuted}
\colorlet{tubeCol}{popdark}

% --- Tube ---
\draw[tubeCol, thick] (0,-0.7) -- (11,-0.7) -- (11,0.7) -- (0,0.7) -- cycle;
% --- Axe optique principal ---
\draw[axisCol, dashed, thin] (0,0) -- (11,0);
\node[axisCol, below] at (0,-0.8) {Sommet miroir primaire};
\node[axisCol, below] at (10,-0.8) {Foyer F (10 cm)};

% --- Miroir Primaire (Concave) à x=0 ---
\draw[mirrorPri, very thick] (0,-0.65) .. controls (0.6,0) .. (0,0.65);
\draw[mirrorPri, thick] (-0.2,0) -- (0,0); % Support arrière
\node[mirrorPri, left] at (-0.3,0) {Miroir primaire $F=\SI{1000}{mm}$};

% --- Foyer F ---
\coordinate (F) at (10,0);
\draw[axisCol, thin] (F) ++(-0.1,-0.1) -- ++(0.2,0.2) (F) ++(0.1,-0.1) -- ++(-0.2,0.2);

% --- Miroir Secondaire (Plan 45°) centré en (9,0) ---
\coordinate (M2c) at (9,0);
\draw[mirrorSec, very thick] (M2c) ++(-0.5,-0.5) -- ++(1,1);
\draw[mirrorSec, thin, ->] (M2c) ++(0.35,-0.35) -- ++(0.3,0.3) node[above right, mirrorSec, font=\small] {Normale};
\node[mirrorSec, below left] at (8.3,-0.5) {Miroir secondaire plan $45^{\circ}$};

% --- Image intermédiaire I' (Foyer objet oculaire) ---
% Distance miroir secondaire -> image = distance objet -> miroir = 1 cm (car F à 10, M2 à 9)
\coordinate (Iprime) at (9,1);
\draw[rayCol, thick] (Iprime) ++(-0.1,0) -- ++(0.2,0) (Iprime) ++(0,-0.1) -- ++(0,0.2);
\node[rayCol, right] at (Iprime) {Image intermédiaire $I'$ (Foyer objet oculaire)};

% --- Oculaire (Lentille convergente) ---
% Foyer objet à I' (9,1). f' = 2.5 cm (pour oculaire 25mm). Centre lentille à y = 1 + 2.5 = 3.5.
\coordinate (OcCenter) at (9,3.5);
\coordinate (FocObj) at (9,1.0);
\coordinate (FocImg) at (9,6.0);
% Lentille
\draw[ocul, very thick] (OcCenter) ++(-0.5,0) to[out=90,in=180] ++(0.5,0.3) to[out=0,in=90] ++(0.5,-0.3) to[out=270,in=0] ++(-0.5,-0.3) to[out=180,in=270] ++(-0.5,0.3) -- cycle;
% Foyers oculaire
\draw[ocul, thin] (FocObj) ++(-0.15,0) -- ++(0.3,0) node[right, ocul, font=\small] {$F'_o$};
\draw[ocul, thin] (FocImg) ++(-0.15,0) -- ++(0.3,0) node[right, ocul, font=\small] {$F'_i$};
\node[ocul, right] at (OcCenter) {Oculaire $f'=\SI{25}{mm}$};

% --- RAYONS : Faisceau incident incliné (angle visuel exagéré ~10°) ---
\def\ang{10} % angle visuel en degrés pour le dessin
\def\h{0.65} % demi-hauteur miroir
% Rayon 1 : Bord supérieur miroir (y=+h)
\draw[rayCol, ->, thick] (-1, {\h + 1*tan(\ang)}) -- (0, \h);
\draw[rayCol, thick] (0, \h) -- (F); % Vers foyer
% Intersection avec miroir secondaire (droite y = x - 9)
% Rayon : y = h - (h/10)x. Inter : h - 0.1h x = x - 9 -> x(1+0.1h) = 9+h -> x approx 9.06
\coordinate (P1) at (9.06, 0.06);
\draw[rayCol, thick] (0, \h) -- (P1);
% Réflexion : Incident (1, -0.065). Normale (-1,1)/sqrt2. Refléchi vers le haut.
\draw[rayCol, thick, ->] (P1) -- (Iprime);

% Rayon 2 : Centre miroir (y=0)
\draw[rayCol, ->, thick] (-1, {0 + 1*tan(\ang)}) -- (0, 0);
\draw[rayCol, thick] (0, 0) -- (F);
\coordinate (P2) at (9,0);
\draw[rayCol, thick] (0,0) -- (P2);
\draw[rayCol, thick, ->] (P2) -- (Iprime); % Refléchi vertical

% Rayon 3 : Bord inférieur miroir (y=-h)
\draw[rayCol, ->, thick] (-1, {-\h + 1*tan(\ang)}) -- (0, -\h);
\draw[rayCol, thick] (0, -\h) -- (F); % y = -h + (h/10)x
% Intersection miroir secondaire : -h + 0.1h x = x - 9 -> x(1-0.1h)=9-h -> x approx 8.94
\coordinate (P3) at (8.94, -0.06);
\draw[rayCol, thick] (0, -\h) -- (P3);
\draw[rayCol, thick, ->] (P3) -- (Iprime);

% --- RAYONS ÉMERGENTS (Afocal) ---
% Depuis I' (objet au foyer objet) -> Lentille -> Parallèles
% Rayon passant par centre optique (non dévié)
\draw[rayCol, thick, ->] (Iprime) -- (OcCenter) -- ++(0, 2.5);
% Rayon passant par Foyer image (sort horizontal) -> non, objet au foyer objet -> sort parallèle à l'axe principal du rayon.
% Construction standard : Rayon depuis I' parallèle à l'axe oculaire -> passe par F'_i.
% Rayon depuis I' vers centre O -> sort même direction.
% Ici axe oculaire = vertical.
% Rayon 1 (depuis I' vers centre O) : déjà tracé.
% Rayon 2 (depuis I' horizontal) : 
\draw[rayCol, thick] (Iprime) -- ++(1,0) node[midway, above, font=\tiny, rayCol] {parallèle axe oculaire};
\draw[rayCol, thick, ->] (10,1) -- (10,6) -- ++(0, 1.5); % Passe par F'_i (9,6) -> attend x=9. Correction:
\draw[rayCol, thick] (Iprime) -- ++(-1,0) coordinate (tmp);
\draw[rayCol, thick, ->] (tmp) -- (FocImg) -- ++(0, 1.5);
% Rayon 3 (symétrique)
\draw[rayCol, thick] (Iprime) -- ++(0.5, -0.2) -- (OcCenter) -- ++(0.5, 1.5); % Approx

% Légendes angles
\draw[->, axisCol] (0.5, 0.2) arc (15:10:0.5) node[midway, right, axisCol, font=\small] {$\theta$};
\draw[->, axisCol] (9, 3.8) arc (-90:-80:0.5) node[midway, right, axisCol, font=\small] {$\theta'$};

\end{tikzpicture}
\legende{Schéma du télescope de Newton (échelle 1/10) : marche d'un faisceau issu d'un cratère lunaire (angle $\theta$ exagéré). L'image intermédiaire $I'$ se forme dans le plan focal objet de l'oculaire. L'instrument est afocal : rayons entrants parallèles, rayons sortants parallèles.}
\end{popfigure}

\textbf{2. Position de l'image intermédiaire.}
L'image intermédiaire $I'$ (réelle, renversée) se forme dans le \textbf{plan focal du miroir primaire}, interceptée par le miroir secondaire qui la renvoie latéralement. Elle coïncide avec le \textbf{foyer objet de l'oculaire} (réglage afocal).

\textbf{3. Grossissements et puissances.}
On applique $G = \dfrac{F}{f'}$ (avec $F$ et $f'$ en mm) et $P = \dfrac{1}{f'}$ (avec $f'$ en m).
\begin{align*}
G_1 &= \frac{1000}{25} = 40          & P_1 &= \frac{1}{0,025} = \SI{40}{\delta} \\
G_2 &= \frac{1000}{12,5} = 80       & P_2 &= \frac{1}{0,0125} = \SI{80}{\delta} \\
G_3 &= \frac{1000}{6} \approx 166,7 & P_3 &= \frac{1}{0,006} \approx \SI{166,7}{\delta}
\end{align*}

\begin{center}
\begin{tabular}{|c|c|c|c|}
\hline
\rowcolor{popblueL} Oculaire $f'$ (mm) & 25 & 12,5 & 6 \\ \hline
Grossissement $G$ & 40 & 80 & $\approx 167$ \\ \hline
Puissance $P$ ($\delta$) & 40 & 80 & $\approx 167$ \\ \hline
\end{tabular}
\end{center}
\textit{Vérification dimensionnelle :} $G = F/f'$ est un rapport de longueurs, sans dimension (cohérence : rapport d'angles). $P = 1/f'$ en $\text{m}^{-1} = \delta$ (dioptrie).

\textbf{4. Choix de l'oculaire.}
Grossissement maximal utile : $G_{\max} \approx 2D = 2 \times 130 = 260$.
\begin{itemize}
\item $G_1 = 40 < 260$ : \textbf{Respecté} (grand champ, repérage, Lune entière).
\item $G_2 = 80 < 260$ : \textbf{Respecté} (détails lunaires, planètes).
\item $G_3 \approx 167 < 260$ : \textbf{Respecté} (grossissement maximal utilisable pour ce diamètre, détails fins cratères, séparation satellites de Jupiter).
\end{itemize}
\textbf{Conclusion :} Les trois oculaires sont utilisables. L'oculaire de $\SI{6}{mm}$ donne le grossissement le plus élevé exploitable ($G \approx 167$). En pratique, on commence par le $\SI{25}{mm}$ (champ large, image lumineuse) pour pointer l'objet, puis on « zoome » avec le $\SI{6}{mm}$.

\textbf{5. Image intermédiaire de la Lune.}
Angle apparent $\theta_L = \num{0,5}^{\circ} = \num{0,5} \times \frac{\pi}{180} \approx \num{8,727e-3}\,\text{rad}$.
Taille dans le plan focal : $y' = F \tan\theta_L \approx F\,\theta_L$ (petits angles).
\[ y' = 1000 \times \tan(\num{0,5}^{\circ}) \approx 1000 \times \num{8,727e-3} \approx \SI{8,73}{mm}. \]
L'image réelle de la Lune pleine mesure $\approx \SI{8,7}{mm}$ de diamètre.
\textit{Champ de l'oculaire :} Un oculaire standard (type Plössl) de $\SI{25}{mm}$ a un diaphragme de champ (diamètre apparent) souvent $\sim \SI{20}{mm}$ à $\SI{25}{mm}$. L'image de la Lune ($\SI{8,7}{mm}$) tient donc \textbf{aisément} dans le champ (champ réel $\approx \frac{\text{champ apparent}}{G} \approx \frac{50^{\circ}}{40} \approx 1,25^{\circ} > 0,5^{\circ}$).
Pour le cratère ($\theta_c = \num{0,01}^{\circ}$) : $y'_c = 1000 \tan(\num{0,01}^{\circ}) \approx \SI{0,175}{mm}$. Ce détail est résolu par l'oculaire de $\SI{6}{mm}$ ($G=167$, taille angulaire finale $\approx 1,67^{\circ}$, bien visible).

\textbf{6. Argumentaire pour le proviseur (rédaction type).}
« Monsieur le Proviseur, pour un budget identique, un télescope de Newton de $\SI{130}{mm}$ de diamètre collecte $\left(\frac{130}{70}\right)^2 \approx 3,4$ fois plus de lumière qu'une lunette de $\SI{70}{mm}$ (prix équivalent), car le flux lumineux est proportionnel à $D^2$. Or, les satellites de Jupiter sont des sources ponctuelles de faible magnitude : leur détection dépend avant tout de la \textbf{luminosité} (ouverture), non du grossissement. De plus, le miroir primaire ne souffre d'\textbf{aucune aberration chromatique} (la réflexion est achromatique), contrairement à l'objectif d'une lunette achromatique bon marché qui présente un résidu coloré (franges violettes) gênant sur les planètes. Enfin, la fabrication d'un grand miroir unique est bien moins coûteuse et plus aisée qu'un doublet achromatique de grand diamètre. Le télescope de Newton est donc, à budget fixe, l'instrument \textbf{le plus performant} pour notre Nuit des étoiles. »
\end{exoresolu}

\courssub{Bilan de l'activité}

Cette activité t'a permis de mettre en évidence les points essentiels de la leçon :
\begin{itemize}
\item le télescope de Newton associe un \cle{miroir concave} (objectif), un \cle{miroir plan} de renvoi et un \cle{oculaire} ; l'image intermédiaire se forme dans le plan focal du miroir et l'image finale est rejetée à l'infini en réglage afocal ;
\item le \cle{grossissement} $G = F/f'$ augmente quand la focale de l'oculaire diminue, mais il est limité en pratique par le critère $G_{\max} \approx 2D$ : ici les trois oculaires sont utilisables, le meilleur compromis pour le détail étant l'oculaire de $\SI{6}{mm}$ ;
\item la taille de l'image intermédiaire ne dépend que de $F$ et de l'angle $\theta$ ($y' = F\tan\theta$) : la Lune donne une image de $\SI{8,7}{mm}$ dans le plan focal ;
\item l'intérêt majeur du télescope sur la lunette est le \cle{grand diamètre à faible coût} (luminosité en $D^2$) et l'\cle{absence d'aberrations chromatiques}, au prix d'une légère obstruction par le miroir secondaire et d'un entretien régulier (collimation).
\end{itemize}
Tu es maintenant capable de dimensionner, régler et justifier le choix d'un instrument d'observation astronomique : compétence directement mobilisable au Probatoire comme dans la vie du club scientifique de ton lycée.

\newpage

\coursec{Méthodes et exercices résolus}

\courssub{Méthodes et exercices résolus}

\begin{methode}[Expliquer le fonctionnement d'un télescope de Newton]
\textbf{Objectif :} Rédiger une explication claire, structurée et complète du principe de fonctionnement, attendue au Probatoire.
\begin{enumerate}
    \item \textbf{Identifier les composants :} Citer le miroir primaire (concave, parabolique), le miroir secondaire (plan, incliné à $45^\circ$) et l'oculaire (lentille convergente). Préciser leurs rôles respectifs.
    \item \textbf{Décrire le cheminement de la lumière :}
    \begin{itemize}
        \item Faisceau incident \emph{parallèle} (objet à l'infini) $\rightarrow$ réflexion sur le miroir primaire $\rightarrow$ convergence vers le foyer image $F'_1$ du primaire.
        \item Interception par le miroir secondaire \emph{avant} le foyer $F'_1$ $\rightarrow$ déviation de $90^\circ$ vers le côté du tube.
        \item Formation d'une \textbf{image réelle, inversée, réduite} au plan focal du primaire (au niveau du miroir secondaire dévié).
        \item Cette image réelle sert d'\textbf{objet réel} pour l'oculaire.
        \item L'oculaire fonctionne comme une \textbf{loupe} : il est placé de sorte que son foyer objet $F_{o2}$ coïncide avec l'image intermédiaire $\rightarrow$ émergence d'un faisceau \emph{parallèle} (image finale à l'infini, observation détendue).
    \end{itemize}
    \item \textbf{Préciser la nature des images :} Image intermédiaire : réelle, inversée. Image finale : virtuelle, inversée par rapport à l'objet initial (donc \emph{renversée} au total), à l'infini (réglage normal) ou à distance finie (réglage pour œil myope/hypermétrope).
    \item \textbf{Justifier le choix du miroir parabolique :} Éviter l'aberration sphérique pour les faisceaux parallèles incidents (stigmatisme rigoureux pour les rayons paraxiaux et marginaux).
\end{enumerate}
\cle{Réflexe Probatoire} : Toujours mentionner \og objet à l'infini \fg, \og image réelle au foyer du primaire \fg, \og oculaire utilisé comme loupe \fg, \og image finale à l'infini (réglage normal) \fg.
\end{methode}

\begin{exoresolu}[Fonctionnement et description -- Type Probatoire]
\textbf{Énoncé :} Un télescope de Newton est utilisé pour observer la Lune. Son miroir primaire a une distance focale $f'_1 = \SI{1200}{mm}$. L'oculaire utilisé a une distance focale $f'_2 = \SI{20}{mm}$. Le miroir secondaire est plan et incliné à $45^\circ$.
\begin{enumerate}
    \item Faire un schéma annoté du télescope en indiquant les trois éléments optiques principaux, le sens de la lumière, le foyer image $F'_1$ du primaire et le foyer objet $F_{o2}$ de l'oculaire. Préciser la position de l'image intermédiaire $A'B'$.
    \item Expliquer pourquoi le miroir primaire est de forme parabolique et non sphérique.
    \item Déterminer la nature (réelle/virtuelle, droite/inversée) et la position de l'image finale observée par l'astronome en réglage normal (œil détendu).
    \item Calculer le grossissement $G$ de cet instrument.
\end{enumerate}
\tcblower
\textbf{Solution détaillée :}
\begin{enumerate}
    \item \textbf{Schéma :} (Description textuelle pour le code \LaTeX, l'élève doit le tracer)
    \begin{itemize}
        \item Tube cylindrique. Miroir primaire (courbe concave) au fond à gauche. Miroir secondaire (segment droit incliné $45^\circ$) au centre du tube. Oculaire (lentille convexe) sur le côté droit du tube.
        \item Rayons incidents parallèles (venant de la droite sur le schéma horizontal) $\rightarrow$ réflexion sur primaire $\rightarrow$ convergence vers $F'_1$ (sur l'axe optique principal, à $\SI{1200}{mm}$ du sommet du primaire).
        \item Miroir secondaire placé \emph{avant} $F'_1$ (ex: à $\SI{1100}{mm}$ du primaire). Réflexion à $90^\circ$ vers l'oculaire.
        \item Image intermédiaire $A'B'$ formée \emph{au plan focal du primaire} (donc au niveau de $F'_1$), mais déviée latéralement par le secondaire. Elle est \textbf{réelle, inversée}.
        \item Oculaire positionné de sorte que son foyer objet $F_{o2}$ coïncide avec le plan de $A'B'$ (distance oculaire-$A'B'$ = $f'_2 = \SI{20}{mm}$).
    \end{itemize}
    \item \textbf{Miroir parabolique :} Un miroir sphérique souffre d'\textbf{aberration sphérique} : les rayons marginaux coupent l'axe plus près du miroir que les rayons paraxiaux. L'image d'un objet ponctuel à l'infini (étoile) est un disque flou (tache de diffusion). Un miroir \textbf{parabolique} est stigmatique rigoureux pour un faisceau incident parallèle : tous les rayons incidents parallèles à l'axe convergent en un unique point $F'$. Cela assure une netteté maximale (résolution angulaire limitée seulement par la diffraction).
    \item \textbf{Image finale (réglage normal) :}
    \begin{itemize}
        \item L'objet pour l'oculaire est l'image intermédiaire $A'B'$ (réelle, inversée), placée au foyer objet $F_{o2}$ de l'oculaire.
        \item L'oculaire (lentille convergente) transforme cet objet réel en un faisceau émergent \textbf{parallèle}.
        \item L'image finale est donc \textbf{virtuelle} (les rayons émergents ne se coupent pas réellement), \textbf{à l'infini}, et \textbf{inversée} par rapport à $A'B'$.
        \item Comme $A'B'$ est déjà inversée par rapport à l'objet initial (Lune), l'image finale est \textbf{renversée} (tête en bas, gauche-droite inversé) par rapport à la Lune.
    \end{itemize}
    \item \textbf{Grossissement $G$ :} Pour un télescope afocal (réglage normal, image à l'infini), le grossissement angulaire est donné par le rapport des distances focales :
    \[
    G = \frac{f'_{\text{objectif}}}{f'_{\text{oculaire}}} = \frac{f'_1}{f'_2} = \frac{\SI{1200}{mm}}{\SI{20}{mm}} = 60.
    \]
    \textbf{Conclusion :} Le télescope fournit un grossissement de \textbf{60 fois}. L'image finale est virtuelle, à l'infini, renversée, et 60 fois plus grande angulairement que l'objet vu à l'œil nu.
\end{enumerate}
\end{exoresolu}

\begin{methode}[Tracer la marche d'un faisceau lumineux dans un télescope de Newton]
\textbf{Objectif :} Réaliser un tracé géométrique rigoureux (construction de l'image intermédiaire par le miroir primaire, déviation par le secondaire, émergence par l'oculaire).
\begin{enumerate}
    \item \textbf{Repères et échelles :} Tracer l'axe optique principal (horizontal pour le primaire, vertical pour l'oculaire). Choisir une échelle pour les distances (ex: 1 cm pour \SI{100}{mm}) et une échelle pour les hauteurs (souvent plus grande pour voir les détails).
    \item \textbf{Miroir primaire (concave, focale $f'_1$) :} Objet à l'infini $\rightarrow$ rayons incidents \textbf{parallèles à l'axe}.
    \begin{itemize}
        \item Rayon 1 : Parallèle à l'axe $\rightarrow$ réflexion passant par le foyer image $F'_1$.
        \item Rayon 2 : Passant par le centre de courbure $C_1$ (si tracé) $\rightarrow$ réflexion sur lui-même (retour).
        \item Rayon 3 : Passant par le foyer objet $F_1$ (symétrique de $F'_1$) $\rightarrow$ réflexion parallèle à l'axe.
        \item \emph{Simplification Probatoire :} Pour un objet à l'infini, on trace au moins \textbf{deux rayons incidents parallèles} (un axial, un marginal ou à hauteur $h$). Leur intersection après réflexion donne $A'B'$.
    \end{itemize}
    \item \textbf{Miroir secondaire (plan, $45^\circ$) :} Appliquer la loi de Descartes (réflexion) : angle d'incidence = angle de réflexion = $45^\circ$.
    \begin{itemize}
        \item Le miroir secondaire dévie le faisceau convergent de $90^\circ$.
        \item L'image intermédiaire $A'B'$ est symétrique de sa position \og sans secondaire \fg par rapport au plan du miroir secondaire. Elle se forme dans le plan focal du primaire, décalée latéralement.
        \item \textbf{Astuce :} Traiter le miroir secondaire comme un miroir plan : l'image $A'B'$ est l'image par symétrie axiale (par rapport au plan du miroir) de l'image $A''B''$ qui se serait formée sans lui.
    \end{itemize}
    \item \textbf{Oculaire (lentille mince convergente, focale $f'_2$) :} L'image intermédiaire $A'B'$ est l'objet réel pour l'oculaire.
    \begin{itemize}
        \item \textbf{Réglage normal (image à l'infini) :} Placer $A'B'$ au foyer objet $F_{o2}$ de l'oculaire (distance $f'_2$).
        \item Tracer depuis un point $B'$ de l'image intermédiaire :
            \begin{itemize}
                \item Rayon passant par le centre optique $O_2$ de l'oculaire $\rightarrow$ non dévié.
                \item Rayon parallèle à l'axe de l'oculaire $\rightarrow$ réfracté passant par le foyer image $F'_{2}$.
            \end{itemize}
        \item Les rayons émergents sont \textbf{parallèles} entre eux $\rightarrow$ image finale à l'infini.
    \end{itemize}
    \item \textbf{Légender :} $F'_1$, $F_{o2}$, $F'_2$, $A'B'$ (image intermédiaire), sens des flèches, valeurs des distances focales.
\end{enumerate}
\cle{Attention} : Le miroir secondaire ne change \textbf{pas} la distance focale effective du primaire ni la taille de l'image intermédiaire $A'B'$ (seule sa position latérale change). Ne pas tracer de convergence vers un nouveau foyer après le secondaire.
\end{methode}

\begin{exoresolu}[Tracé géométrique complet -- Type Probatoire]
\textbf{Énoncé :} On considère un télescope de Newton de focale primaire $f'_1 = \SI{800}{mm}$ et d'oculaire $f'_2 = \SI{10}{mm}$. Le miroir secondaire est placé à $\SI{750}{mm}$ du sommet du miroir primaire. L'instrument est en réglage normal (image finale à l'infini). L'objet observé est une étoile (faisceau incident parallèle, incliné de $\alpha = \num{0,01}\,\text{rad}$ par rapport à l'axe).
\begin{enumerate}
    \item Sur une figure (échelle horizontale : 1 cm = \SI{100}{mm} ; échelle verticale : 1 cm = \SI{1}{mm}), tracer le cheminement d'un faisceau lumineux incident parallèle à l'axe (rayon axial) et d'un faisceau incident incliné de $\alpha$ (rayon principal marginal). Montrer la formation de l'image intermédiaire $A'B'$ et l'émergence par l'oculaire.
    \item Déterminer graphiquement et par le calcul la hauteur $h'$ de l'image intermédiaire $A'B'$.
    \item Vérifier que l'image finale émerge sous un angle $\alpha' = G \times \alpha$.
\end{enumerate}
\tcblower
\textbf{Solution détaillée :}
\begin{enumerate}
    \item \textbf{Construction du tracé (Description pour l'élève) :}
    \begin{itemize}
        \item \textbf{Partie primaire (horizontal) :} Axe horizontal. Miroir primaire en $x=0$. $F'_1$ à $x=\SI{80}{cm}$ ($\SI{800}{mm}$). Miroir secondaire (segment $45^\circ$) centré vers $x=\SI{75}{cm}$.
        \item \textbf{Rayon axial (bleu) :} Incident horizontal $\rightarrow$ réflexion sur primaire $\rightarrow$ passe par $F'_1$ $\rightarrow$ rencontre le secondaire $\rightarrow$ réfléchi verticalement vers le haut (vers l'oculaire).
        \item \textbf{Rayon principal incliné (rouge) :} Incident avec angle $\alpha = \num{0,01}\,\text{rad}$ (descendant si $\alpha>0$). Parallèle à l'axe incident $\rightarrow$ réflexion sur primaire $\rightarrow$ passe par $F'_1$ \emph{mais décalé}. Hauteur en $F'_1$ : $h' = f'_1 \times \tan\alpha \approx f'_1 \alpha = 800 \times \num{0,01} = \SI{8}{mm}$. Ce rayon coupe le plan focal en $B'$ ($h' = -\SI{8}{mm}$ si incident descendant).
        \item Ce rayon convergent vers $B'$ rencontre le secondaire $\rightarrow$ réfléchi (loi de Descartes, $45^\circ$) $\rightarrow$ part vers l'oculaire en visant le point $B'$ (image intermédiaire).
        \item \textbf{Oculaire (vertical) :} Placé au-dessus du secondaire. Son centre optique $O_2$ sur le rayon axial dévié. Son foyer objet $F_{o2}$ à $\SI{10}{mm}$ en dessous de $O_2$ (donc confondu avec le plan de $B'$).
        \item \textbf{Émergence :} Depuis $B'$ (hauteur $-\SI{8}{mm}$ par rapport à l'axe oculaire) :
            \begin{itemize}
                \item Rayon $B'O_2$ $\rightarrow$ sort non dévié.
                \item Rayon depuis $B'$ parallèle à l'axe oculaire $\rightarrow$ réfracté vers $F'_2$ (foyer image oculaire, $\SI{10}{mm}$ au-dessus de $O_2$).
            \end{itemize}
        \item Les deux rayons émergents sont \textbf{parallèles}, inclinés vers le haut. Angle émergent $\alpha'$.
    \end{itemize}
    \item \textbf{Hauteur $h'$ de l'image intermédiaire :}
    \begin{itemize}
        \item \textbf{Calcul :} Objet à l'infini, angle $\alpha$. Image au foyer du primaire.
        \[
        h' = f'_1 \tan\alpha \approx f'_1 \alpha = \SI{800e-3}{m} \times \num{0,01} = \SI{8e-3}{m} = \mathbf{\SI{8}{mm}}.
        \]
        (Signe négatif car image inversée).
        \item \textbf{Graphique :} Sur la figure, échelle verticale 1 cm = \SI{1}{mm}. $h' = \SI{8}{mm}$ $\rightarrow$ 8 cm sur le papier. Cohérent.
    \end{itemize}
    \item \textbf{Vérification de l'angle émergent $\alpha'$ :}
    \begin{itemize}
        \item \textbf{Par le grossissement :} $G = \frac{f'_1}{f'_2} = \frac{800}{10} = 80$.
        \[
        \alpha' = G \alpha = 80 \times \num{0,01} = \mathbf{\num{0,8}\,rad}.
        \]
        \item \textbf{Par la géométrie oculaire (loupe) :} L'objet pour l'oculaire a une hauteur $h' = -\SI{8}{mm}$ placé à $f'_2 = \SI{10}{mm}$ du centre $O_2$.
        L'angle sous-tendu par l'objet au centre de l'oculaire (angle émergent pour le rayon passant par $O_2$) est :
        \[
        \tan\alpha' = \frac{|h'|}{f'_2} = \frac{8}{10} = \num{0,8} \implies \alpha' \approx \mathbf{\num{0,8}\,rad}.
        \]
        \item \textbf{Cohérence :} Les deux méthodes donnent $\alpha' = \num{0,8}\,\text{rad}$. L'image finale est inversée par rapport à l'image intermédiaire (donc droite par rapport à l'objet initial ? Non : Primaire inverse une fois, Oculaire inverse une fois $\rightarrow$ Total : Renversée). Le rayon incident descendant ($\alpha>0$) donne un rayon émergent montant ($\alpha'>0$), confirmant l'inversion totale.
    \end{itemize}
    \textbf{Conclusion :} Le tracé géométrique confirme la formule $G = f'_1/f'_2$ et la formation d'une image intermédiaire de hauteur $h' = f'_1 \alpha$.
\end{enumerate}
\end{exoresolu}

\begin{methode}[Calculer le grossissement et la puissance d'un télescope de Newton]
\textbf{Objectif :} Appliquer les formules de grossissement $G$ et de puissance $P$ (vergences) en distinguant les cas : réglage normal (afocal) et réglage pour œil accommodé (image à distance finie).
\begin{enumerate}
    \item \textbf{Identifier les données :} $f'_{\text{obj}}$ (focale miroir primaire), $f'_{\text{oc}}$ (focale oculaire). Unités : \textbf{mètres} pour les vergences (dioptries), \textbf{mm ou m} pour $G$ (sans unité).
    \item \textbf{Réglage normal (image finale à l'infini, œil détendu) :}
    \begin{itemize}
        \item Grossissement angulaire : $\boxed{G = \dfrac{f'_{\text{obj}}}{f'_{\text{oc}}}}$.
        \item Puissance (vergence) de l'instrument : $\boxed{P = \dfrac{1}{f'_{\text{obj}}} + \dfrac{1}{f'_{\text{oc}}}}$ ? \textbf{FAUX pour télescope afocal.}
        \item \textbf{Correction :} Un système afocal (entrée parallèle, sortie parallèle) a une \textbf{vergence nulle} ($P=0\,\delta$). Il n'a pas de foyer image global.
        \item On parle de \textbf{puissance de grossissement} $G$, pas de vergence globale.
    \end{itemize}
    \item \textbf{Réglage pour œil accommodé (image finale à distance finie $D$, ex: $D = -\SI{250}{mm}$ pour point proche) :}
    \begin{itemize}
        \item L'oculaire doit former une image virtuelle à $D$. Relation de conjugaison pour l'oculaire (convention Descartes : $p<0$ objet réel, $p'>0$ image réelle, $p'<0$ image virtuelle, $f'>0$ convergente) : $\dfrac{1}{f'_{\text{oc}}} = \dfrac{1}{p'} - \dfrac{1}{p}$ avec $p' = D$ (négatif) et $p = -\overline{O_2A'B'}$ (objet réel pour l'oculaire, donc $p<0$).
        \item Distance oculaire-image intermédiaire : $\overline{O_2A'B'} = |p| = \dfrac{f'_{\text{oc}} |D|}{f'_{\text{oc}} + |D|}$ (valeur absolue $< f'_{\text{oc}}$).
        \item Grossissement angulaire (approché si angles petits) : $G \approx \dfrac{f'_{\text{obj}}}{|p|} = \dfrac{f'_{\text{obj}}}{f'_{\text{oc}}} \left(1 + \frac{f'_{\text{oc}}}{|D|}\right)$.
        \item Pour $D = -\SI{250}{mm}$ (vision nette de près), $G \approx G_{\text{normal}} \left(1 + \frac{f'_{\text{oc}}}{250}\right)$. Le grossissement augmente légèrement.
    \end{itemize}
    \item \textbf{Pupille de sortie :} Diamètre $d_s = \dfrac{D_{\text{primaire}}}{|G|}$. Condition confort : $d_s \approx \SIrange{2}{7}{mm}$ (pupille œil).
    \item \textbf{Pouvoir séparateur (Résolution) :} $\theta_{\min} \approx 1{,}22 \frac{\lambda}{D_{\text{primaire}}}$ (critère de Rayleigh). Indépendant de l'oculaire.
\end{enumerate}
\cle{Piège} : Ne jamais écrire $P = 1/f'_{\text{obj}} + 1/f'_{\text{oc}}$ pour un télescope. C'est la formule pour deux lentilles \emph{contact} ou microscope. Le télescope est \textbf{afocal} ($P=0$).
\end{methode}

\begin{exoresolu}[Grossissement, puissance, pupille de sortie -- Type Probatoire]
\textbf{Énoncé :} Un télescope de Newton a un miroir primaire de diamètre $D_1 = \SI{200}{mm}$ et de focale $f'_1 = \SI{1000}{mm}$. On dispose de deux oculaires : $O_1$ ($f'_2 = \SI{25}{mm}$) et $O_2$ ($f'_2 = \SI{10}{mm}$). L'observateur a une pupille d'entrée de $\SI{6}{mm}$ (nuit) et un point proche à $\SI{25}{cm}$.
\begin{enumerate}
    \item Calculer le grossissement $G_1$ et $G_2$ en réglage normal (image à l'infini).
    \item Calculer le diamètre de la pupille de sortie $d_{s1}$ et $d_{s2}$. Quel oculaire est le plus confortable pour l'observation nocturne ? Justifier.
    \item Calculer le grossissement $G_{2,\text{proche}}$ avec l'oculaire $O_2$ si l'observateur accommode au point proche ($\SI{25}{cm}$).
    \item Déterminer la puissance (vergence) globale de l'instrument en réglage normal. Expliquer le résultat.
    \item Calculer la résolution angulaire théorique (critère de Rayleigh, $\lambda = \SI{550}{nm}$) pour ce télescope. L'oculaire $O_2$ permet-il d'atteindre cette limite (œil humain $\approx 1'$ d'arc) ?
\end{enumerate}
\tcblower
\textbf{Solution détaillée :}
\begin{enumerate}
    \item \textbf{Grossissements (réglage normal) :}
    \[
    G_1 = \frac{f'_1}{f'_{O_1}} = \frac{1000}{25} = \mathbf{40}. \qquad
    G_2 = \frac{f'_1}{f'_{O_2}} = \frac{1000}{10} = \mathbf{100}.
    \]
    \item \textbf{Pupilles de sortie :} $d_s = \frac{D_1}{|G|}$.
    \[
    d_{s1} = \frac{200}{40} = \mathbf{\SI{5,0}{mm}}. \qquad
    d_{s2} = \frac{200}{100} = \mathbf{\SI{2,0}{mm}}.
    \]
    \textbf{Analyse :} La pupille de l'œil (nuit) est $\SI{6}{mm}$.
    \begin{itemize}
        \item $O_1$ : $d_{s1}=\SI{5}{mm} < \SI{6}{mm}$. Tout le faisceau entre dans l'œil. \textbf{Confortable}.
        \item $O_2$ : $d_{s2}=\SI{2}{mm} \ll \SI{6}{mm}$. L'œil ne reçoit qu'une petite partie de la lumière collectée. L'image paraît plus sombre, l'alignement de l'œil est critique (\og trou de serrure \fg). \textbf{Moin confortable}.
    \end{itemize}
    \textbf{Conclusion :} $O_1$ (\SI{25}{mm}) est plus confortable pour l'observation nocturne (nébuleuses, galaxies). $O_2$ réserve aux planètes/Lune (fort grossissement, objet lumineux).
    \item \textbf{Grossissement au point proche ($O_2$) :}
    Distance image finale $p' = D = -\SI{250}{mm}$ (virtuelle, même côté que l'objet pour l'oculaire).
    Relation de conjugaison (convention Descartes : $p<0$ objet réel, $p'<0$ image virtuelle) :
    \[
    \frac{1}{f'_2} = \frac{1}{p'} - \frac{1}{p} \implies \frac{1}{10} = \frac{1}{-250} - \frac{1}{p} \implies \frac{1}{p} = -\frac{1}{250} - \frac{1}{10} = -\frac{26}{250}
    \]
    \[
    p = -\frac{250}{26} \approx -\SI{9,62}{mm}.
    \]
    Distance oculaire-image intermédiaire $\overline{O_2A'B'} = |p| = \SI{9,62}{mm}$.
    Grossissement angulaire (formule exacte pour petits angles) :
    \[
    G_{2,\text{proche}} = \frac{f'_1}{|p|} = \frac{1000}{9,62} \approx \mathbf{104}.
    \]
    (Formule approchée : $G_2(1 + f'_2/250) = 100(1 + 10/250) = 104$. Cohérent).
    \item \textbf{Puissance globale (réglage normal) :}
    Entrée : faisceau parallèle (vergence $V_{\text{ent}} = 0$). Sortie : faisceau parallèle (vergence $V_{\text{sort}} = 0$).
    Vergence du système $P = V_{\text{sort}} - V_{\text{ent}} = \mathbf{0\,\delta}$ (dioptries).
    \textbf{Explication :} Le télescope est un système \textbf{afocal}. Il ne converge pas la lumière vers un point unique (pas de foyer image global), il modifie seulement l'angle des faisceaux (grossissement angulaire). Il ne possède pas de distance focale image globale.
    \item \textbf{Résolution angulaire (Rayleigh) :}
    \[
    \theta_{\min} = 1{,}22 \frac{\lambda}{D_1} = 1{,}22 \times \frac{550 \times 10^{-9}}{0,200} = 3,355 \times 10^{-6}\,\text{rad}.
    \]
    Conversion en secondes d'arc ($1\,\text{rad} \approx 206265''$) :
    \[
    \theta_{\min} \approx 3,355 \times 10^{-6} \times 206265 \approx \mathbf{0,69''}.
    \]
    \textbf{Résolution de l'œil :} $\approx 1' = 60''$.
    \textbf{Grossissement nécessaire pour rendre visible le détail :} $G_{\text{utile}} \approx \frac{60''}{0,69''} \approx 87$.
    \textbf{Test $O_2$ :} $G_2 = 100 > 87$. \textbf{Oui}, l'oculaire $O_2$ (100x) permet de résoudre la limite de diffraction du télescope (grossissement utile). $O_1$ (40x) ne le permet pas (sous-échantillonnage).
\end{enumerate}
\end{exoresolu}

\begin{methode}[Réaliser la mise au point d'un télescope de Newton]
\textbf{Objectif :} Expliquer la manipulation pratique et calculer le déplacement nécessaire de l'oculaire (ou du miroir primaire) pour passer de l'observation d'un objet à l'infini à un objet à distance finie, ou pour adapter à un œil myope/hypermétrope.
\begin{enumerate}
    \item \textbf{Principe :} La mise au point consiste à placer l'image intermédiaire $A'B'$ (formée par le primaire) \textbf{exactement dans le plan focal objet de l'oculaire} ($F_{o2}$).
    \item \textbf{Convention de signes (Vergence / Gauss) :} Pour le miroir primaire (réflexion), on utilise souvent la convention \og réel positif \fg : $f'_1>0$, $p_1<0$ (objet réel), $p'_1>0$ (image réelle). Relation : $\frac{1}{f'_1} = \frac{1}{p'_1} - \frac{1}{p_1}$. Pour l'oculaire (réfraction), convention Descartes standard : $f'_2>0$, $p_2<0$ (objet réel), $p'_2>0$ (image réelle), $p'_2<0$ (image virtuelle). Relation : $\frac{1}{f'_2} = \frac{1}{p'_2} - \frac{1}{p_2}$.
    \item \textbf{Cas 1 : Objet à l'infini $\rightarrow$ Objet à distance finie $d$ (ex: oiseau à 50 m).}
    \begin{itemize}
        \item Miroir primaire : $p_1 = d < 0$. Relation $\frac{1}{f'_1} = \frac{1}{p'_1} - \frac{1}{p_1} \implies \frac{1}{p'_1} = \frac{1}{f'_1} + \frac{1}{p_1}$.
        \item Comme $|d| \gg f'_1$, $p'_1 \approx f'_1 \left(1 + \frac{f'_1}{|d|}\right)$. L'image s'éloigne du miroir (recule vers l'intérieur du tube).
        \item Déplacement de l'image intermédiaire : $\Delta p'_1 = p'_1 - f'_1 \approx \frac{f'_1{}^2}{|d|}$ (valeur positive, l'image recule).
        \item \textbf{Action :} Il faut \textbf{éloigner l'oculaire du miroir primaire} (ou reculer le miroir primaire) de cette même valeur $\Delta p'_1$ pour que $A'B'$ reste sur $F_{o2}$.
    \end{itemize}
    \item \textbf{Cas 2 : Adaptation à un œil myope (point loin $P_L < \infty$) ou hypermétrope (point proche $P_P > \SI{25}{cm}$).}
    \begin{itemize}
        \item L'oculaire doit former l'image finale au point loin (myope) ou au point proche (hypermétrope).
        \item Calculer la distance objet $p_2$ (oculaire-image intermédiaire) via $\frac{1}{f'_2} = \frac{1}{p'_2} - \frac{1}{p_2}$.
        \item Comparer $|p_2|$ à $f'_2$ (position réglage normal).
        \item Myope ($p'_2 = P_L < 0$, $|P_L| < \infty$) : $|p_2| < f'_2$. Il faut \textbf{rapprocher l'oculaire de l'image intermédiaire} (avancer l'oculaire vers le miroir).
        \item Hypermétrope ($p'_2 = -P_P$, $|P_P| > \SI{25}{cm}$) : $|p_2| < f'_2$ aussi (mais moins). Il faut \textbf{avancer l'oculaire}.
        \item \emph{Note :} Pour un emmétrope réglé au point proche ($\SI{25}{cm}$), on avance aussi l'oculaire.
    \end{itemize}
    \item \textbf{Pratique :} On utilise le \textbf{porte-oculaire} (crémaillère) pour déplacer l'oculaire le long de l'axe optique. Ne pas toucher au miroir primaire (collimation perdue).
\end{enumerate}
\cle{Formule clé (objet lointain distance $d$) :} Déplacement mise au point $\Delta x \approx \dfrac{f'_1{}^2}{|d|}$. (Ex: $f'_1=\SI{1}{m}, d=\SI{100}{m} \rightarrow \Delta x = \SI{1}{cm}$).
\end{methode}

\begin{exoresolu}[Mise au point : objet proche et œil myope -- Type Probatoire]
\textbf{Énoncé :} Un télescope de Newton a $f'_1 = \SI{1200}{mm}$ et $f'_2 = \SI{15}{mm}$. Il est initialement au point sur une étoile (réglage normal : distance miroir primaire - oculaire = $L_0 = f'_1 + f'_2 = \SI{1215}{mm}$, en négligeant l'épaisseur du secondaire et le décalage latéral).
\begin{enumerate}
    \item L'observateur vise un oiseau perché à $d = \SI{30}{m}$. Calculer le déplacement $\Delta L$ du porte-oculaire nécessaire pour refaire la netteté. Sens du déplacement ?
    \item L'observateur est myope non corrigé : son point loin est à $P_L = -\SI{2,00}{m}$ (il voit net jusqu'à 2 m sans lunettes). L'instrument est réglé normalement sur une étoile (image à l'infini). L'observateur voit flou. De combien et dans quel sens doit-il déplacer l'oculaire pour voir net l'étoile sans lunettes ?
\end{enumerate}
\tcblower
\textbf{Solution détaillée :}
\begin{enumerate}
    \item \textbf{Objet à distance finie ($d = -\SI{30}{m} = -\SI{30000}{mm}$) :}
    \begin{itemize}
        \item Miroir primaire : $f'_1 = +\SI{1200}{mm}$. Objet $AB$ à $p_1 = d = -\SI{30000}{mm}$.
        \item Relation de conjugaison (convention réel positif pour miroir : $p<0$, $p'>0$) :
        \[
        \frac{1}{f'_1} = \frac{1}{p'_1} - \frac{1}{p_1} \implies \frac{1}{p'_1} = \frac{1}{f'_1} + \frac{1}{p_1} = \frac{1}{1200} + \frac{1}{-30000} = \frac{25 - 1}{30000} = \frac{24}{30000} = \frac{1}{1250}.
        \]
        \item Position image intermédiaire : $p'_1 = \mathbf{\SI{1250}{mm}}$.
        \item Position initiale (étoile) : $p'_{1,\infty} = f'_1 = \SI{1200}{mm}$.
        \item Déplacement de l'image : $\Delta p'_1 = 1250 - 1200 = \mathbf{+\SI{50}{mm}}$. L'image recule de $\SI{5}{cm}$ vers le fond du tube (loin du miroir).
        \item Pour garder l'image sur le foyer de l'oculaire ($F_{o2}$), il faut \textbf{reculer l'oculaire de $\SI{50}{mm}$} (l'éloigner du miroir primaire).
        \item \textbf{Vérification formule approchée :} $\Delta x \approx \frac{f'_1{}^2}{|d|} = \frac{1200^2}{30000} = \frac{1\,440\,000}{30\,000} = \SI{48}{mm} \approx \SI{50}{mm}$. Bon ordre de grandeur.
    \end{itemize}
    \item \textbf{Œil myope (Point loin $P_L = -\SI{2000}{mm}$) :}
    \begin{itemize}
        \item Objectif : Image finale virtuelle au point loin de l'œil : $p'_2 = P_L = -\SI{2000}{mm}$.
        \item Oculaire : $f'_2 = +\SI{15}{mm}$.
        \item Relation de conjugaison pour l'oculaire (convention Descartes : $p_2<0$ objet réel, $p'_2<0$ image virtuelle) :
        \[
        \frac{1}{f'_2} = \frac{1}{p'_2} - \frac{1}{p_2} \implies \frac{1}{15} = \frac{1}{-2000} - \frac{1}{p_2}
        \]
        \[
        \frac{1}{p_2} = -\frac{1}{2000} - \frac{1}{15} = -\frac{15 + 2000}{30000} = -\frac{2015}{30000}
        \]
        \[
        p_2 = -\frac{30000}{2015} \approx -\SI{14,89}{mm}.
        \]
        \item Distance oculaire-image intermédiaire requise : $|p_2| \approx \SI{14,89}{mm}$.
        \item Distance initiale (réglage normal, image à l'infini) : $|p_{2,\infty}| = f'_2 = \SI{15,00}{mm}$.
        \item Différence : $\Delta p_2 = 14,89 - 15,00 = -\SI{0,11}{mm}$.
        \item L'image intermédiaire doit être \textbf{rapprochée de l'oculaire} de $\SI{0,11}{mm}$.
        \item Comme l'image intermédiaire est fixe (étoile $\rightarrow$ au foyer du primaire), il faut \textbf{avancer l'oculaire vers le miroir primaire de $\SI{0,11}{mm}$}.
    \end{itemize}
    \textbf{Conclusion :} Pour l'oiseau à 30 m, reculer l'oculaire de \textbf{\SI{5}{cm}}. Pour l'œil myope (2 m), avancer l'oculaire de \textbf{\SI{0,11}{mm}} (quasi imperceptible, souvent dans la tolérance de mise au point).
\end{enumerate}
\end{exoresolu}

\begin{methode}[Comparer le télescope de Newton aux instruments à lentilles (Lunette astronomique / Télescope de Galilée)]
\textbf{Objectif :} Structurer une argumentation physique (avantages/inconvénients) pour une question de type dissertation ou QCM justifié.
\begin{enumerate}
    \item \textbf{Critère 1 : Aberration chromatique.}
    \begin{itemize}
        \item \textbf{Newton (Miroirs) :} \textbf{Zéro aberration chromatique.} La réflexion ne dépend pas de la longueur d'onde ($n=-1$ pour tous $\lambda$). Image parfaitement achromatique.
        \item \textbf{Lunette (Lentilles) :} Aberration chromatique longitudinale (foyer bleu plus court que foyer rouge) et transversale. Nécessite des doublets achromatiques (couronne + flint) ou apochromatiques (coûteux) pour corriger.
    \end{itemize}
    \item \textbf{Critère 2 : Ouverture / Diamètre / Prix.}
    \begin{itemize}
        \item \textbf{Newton :} Miroir unique (une face à polir). Support par l'arrière possible (cellule). \textbf{Coût $\propto D^2$ à $D^{2,5}$}. Facile à réaliser en grand diamètre (amateurs : 200-400 mm courants ; pro : 8-10 m).
        \item \textbf{Lunette :} Objectif = 2 à 4 lentilles (2 faces chacune $\rightarrow$ 4 à 8 surfaces à polir). Support par le bord seulement (risque de déformation). Verre optique cher. \textbf{Coût $\propto D^3$ à $D^4$}. Limité à $\approx \SI{1}{m}$ (Yerkes 1,02 m).
    \end{itemize}
    \item \textbf{Critère 3 : Obstruction centrale (Miroir secondaire).}
    \begin{itemize}
        \item \textbf{Newton :} Miroir secondaire + porte-oculaire $\rightarrow$ \textbf{obstruction centrale} (typiquement 15-25\% du diamètre). Conséquences : Perte de contraste (énergie dans les anneaux de diffraction), perte de lumière ($\approx (d/D)^2$), pics de diffraction (croix sur les étoiles brillantes).
        \item \textbf{Lunette :} \textbf{Aucune obstruction.} Meilleur contraste, images plus "propres" (planètes, doubles étoiles).
    \end{itemize}
    \item \textbf{Critère 4 : Encombrement / Monture.}
    \begin{itemize}
        \item \textbf{Newton :} Tube long ($L \approx f'_1$), lourd au bout (miroir + cellule). Nécessite monture équatoriale ou Dobson stable. Centre de gravité éloigné de la monture.
        \item \textbf{Lunette :} Tube long aussi ($L \approx f'_{\text{obj}}$), mais masse répartie. Plus facile à équilibrer.
    \end{itemize}
    \item \textbf{Critère 5 : Collimation (Alignement).}
    \begin{itemize}
        \item \textbf{Newton :} \textbf{Sensible.} Miroir primaire + secondaire + oculaire à aligner (collimation) régulièrement. Nécessite outils (laser, Cheshire).
        \item \textbf{Lunette :} \textbf{Stable.} Objectif fixé en usine. Quasi pas de maintenance.
    \end{itemize}
    \item \textbf{Synthèse Probatoire :} Newton = \textbf{Roi du rapport ouverture/prix}, idéal pour le ciel profond (nébuleuses, galaxies) grâce à grande ouverture sans chromatisme. Lunette = \textbf{Roi du contraste et robustesse}, idéal planétaire/lunaire, astrophotographie grand champ, débutants.
\end{enumerate}
\end{methode}

\begin{exoresolu}[Analyse comparative et choix d'instrument -- Type Probatoire]
\textbf{Énoncé :} Un club d'astronomie du Cameroun (région de Ngaoundéré, ciel noir, altitude \SI{1100}{m}) souhaite acquérir un instrument pour l'initiation des lycéens (observation visuelle planétaire et ciel profond) et l'astrophotographie de nébuleuses. Budget : \SI{1 500 000}{FCFA}. Deux options :
\begin{itemize}
    \item Option A : Lunette apochromatique (ED) 80 mm, $f/D = 7$ ($f' = \SI{560}{mm}$), sur monture équatoriale motorisée GoTo.
    \item Option B : Newton 200 mm, $f/D = 5$ ($f' = \SI{1000}{mm}$), sur monture Dobson manuelle (non motorisée).
\end{itemize}
Analyser les deux options au regard des critères : chromatisme, pouvoir séparateur, luminosité/étendue, encombrement/monture, astrophotographie. Conclure par un choix argumenté.
\tcblower
\textbf{Solution détaillée :}
\begin{center}
\begin{tabularx}{\linewidth}{|l|X|X|}
\hline
\rowcolor{popblueL} \textbf{Critère} & \textbf{Option A : Lunette APO 80/560} & \textbf{Option B : Newton 200/1000} \\ \hline
\textbf{Chromatisme} & \textbf{Excellent (APO).} Correction résiduelle quasi nulle. Images pures couleur. & \textbf{Nul (Miroirs).} Parfaitement achromatique. \\ \hline
\textbf{Pouvoir séparateur (Rayleigh)} & $\theta_A = 1{,}22 \frac{\SI{550}{nm}}{\SI{0,08}{m}} \approx \SI{1,73}{''}$ & $\theta_B = 1{,}22 \frac{\SI{550}{nm}}{\SI{0,20}{m}} \approx \mathbf{\SI{0,69}{''}}$ \\
& Limité par la diffraction (petit diamètre). & \textbf{2,5 fois meilleur.} Résout détails planétaires fins, doubles serrées. \\ \hline
\textbf{Luminosité / Étendue (Ciel profond)} & Surface $\propto 80^2 = 6400$. $f/D=7$ $\rightarrow$ champ large possible. & Surface $\propto 200^2 = \mathbf{40000}$. \textbf{6,25 fois plus de lumière.} $f/D=5$ $\rightarrow$ rapide. \\
& Faible collecte de photons pour nébuleuses. & \textbf{Idéal nébuleuses/galaxies.} \\ \hline
\textbf{Obstruction / Contraste} & \textbf{Aucune.} Contraste maximal. Planètes superbes. & Obstruction $\approx 20-25\%$ (secondaire 40-50mm). Perte contraste légère. Pics diffraction. \\ \hline
\textbf{Monture / Suivi} & \textbf{Équatoriale GoTo motorisée.} Suivi automatique, pointage aisé. Indispensable photo pose longue. & Dobson manuelle. \textbf{Pas de suivi.} Impossible photo pose longue ($> \SI{30}{s}$). Pointage manuel (apprentissage ciel). \\ \hline
\textbf{Encombrement / Transport} & Tube $\approx \SI{60}{cm}$, léger ($\approx \SI{4-5}{kg}$). Facile transport (moto-taxi, bus). & Tube $\approx \SI{100}{cm}$, lourd ($\approx \SI{15-20}{kg}$ + base Dobson). Difficile transport rural. \\ \hline
\textbf{Prix (Estimation FCFA)} & $\approx \SI{1 200 000}{FCFA} - \SI{1 500 000}{FCFA}$ (Tube + Monture GoTo). & $\approx \SI{400 000}{FCFA} - \SI{600 000}{FCFA}$ (Tube + Dobson). \textbf{Budget respecté largement.} \\ \hline
\end{tabularx}
\end{center}

\textbf{Analyse pour le contexte Cameroun (Ngaoundéré) :}
\begin{itemize}
    \item \textbf{Initiation lycéens :} Dobson (B) est pédagogique : apprend à repérer les objets (star-hopping), robuste, pas d'électronique fragile. Mais pas de suivi = fatigue pour groupes.
    \item \textbf{Astrophotographie nébuleuses :} \textbf{Critère éliminatoire pour B.} Sans monture équatoriale motorisée (équatoriale allemande $\approx \SI{800 000}{FCFA}$ de plus), photo pose longue impossible. Option A a le GoTo inclus.
    \item \textbf{Observation planétaire :} A (80 mm) limite à $\approx 160\times$ utile. B (200 mm) va à $\approx 400\times$ utile. B gagne sur le détail.
    \item \textbf{Observation ciel profond :} B gagne massivement (ouverture).
\end{itemize}

\textbf{Choix argumenté :}
\begin{itemize}
    \item \textbf{Si priorité ASTROPHOTO :} \textbf{Option A (Lunette APO 80 + Monture GoTo).} C'est la \textbf{seule} viable dans le budget pour faire de la photo de nébuleuses (suivi précis, champ large, pas de chromatisme). On sacrifie la résolution planétaire et la magnitude limite visuelle.
    \item \textbf{Si priorité OBSERVATION VISUELLE (Planétaire + Ciel profond) :} \textbf{Option B (Newton 200 Dobson).} Performance visuelle incomparable pour le prix. Mais \textbf{zéro astrophoto pose longue}. Il faut prévoir budget monture équatoriale plus tard.
    \item \textbf{Compromis réel pour le club :} Acheter le \textbf{Newton 200/1000 sur Dobson (Option B)} pour l'observation visuelle (impact "Wahou" garanti sur nébuleuses/planètes pour les élèves) ET une \textbf{petite lunette guide / photo (ex: 60-70 mm) sur une monture équatoriale d'occasion} pour l'initiation photo. Ou attendre budget pour monture EQ6-R Pro pour le Newton.
\end{itemize}
\textbf{Conclusion pédagogique :} Pour un club scolaire au Cameroun, le \textbf{Newton 200 mm Dobson} offre le meilleur retour sur investissement \emph{visuel} (émerveillement, apprentissage du ciel). L'astrophoto nécessite un budget monture distinct et une courbe d'apprentissage raide.
\end{exoresolu}

\begin{methode}[Calculer les caractéristiques du miroir secondaire (plan, 45°) et l'obstruction]
\textbf{Objectif :} Dimensionner le miroir secondaire (taille minimale, ellipse) et calculer le taux d'obstruction pour un newtonien donné.
\begin{enumerate}
    \item \textbf{Données :} $D_1$ (diamètre primaire), $f'_1$ (focale primaire), $d$ (distance sommet primaire - plan du secondaire), $D_{\text{oc}}$ (diamètre du champ oculaire / barreau de mise au point), $L_{\text{oc}}$ (distance secondaire - centre oculaire).
    \item \textbf{Géométrie du faisceau :} Le cône convergent du primaire a une ouverture angulaire définie par $D_1/f'_1$.
    \item \textbf{Diamètre minimal du secondaire (axe mineur $b$) :} Il doit intercepter \textbf{tout le cône axial} (pour l'image centrée) + le cône du \textbf{bord de champ} (pour l'image au bord du champ oculaire).
    \begin{itemize}
        \item Diamètre pour l'axe (cône axial à distance $d$) : $D_{\text{axe}} = D_1 \frac{f'_1 - d}{f'_1}$.
        \item Diamètre pour le bord de champ (éclairement plein champ) : Le rayon marginal du bord de champ passe par le bord du primaire et le bord du champ (hauteur $h_{\text{champ}} = D_{\text{oc}}/2$ au plan focal). Au niveau du secondaire, son décalage par rapport à l'axe est $h_{\text{champ}} \frac{d}{f'_1}$.
        \item \textbf{Formule exacte (axe + champ plein) :} Diamètre mineur $b = D_1 \frac{f'_1 - d}{f'_1} + D_{\text{oc}} \frac{d}{f'_1}$.
        \item Axe majeur $a = b / \cos 45^\circ = b\sqrt{2}$.
    \end{itemize}
    \item \textbf{Taux d'obstruction (surface) :} $\tau = \left(\frac{b}{D_1}\right)^2$ (si on approx le secondaire circulaire de diam $b$). Ou $\tau = \frac{a b}{D_1^2}$ pour l'ellipse exacte. Souvent donné en diamètre linéaire : $\epsilon = b/D_1$.
    \item \textbf{Position du foyer accessible :} Distance secondaire - plan focal = $f'_1 - d$. Doit être $\ge$ longueur du porte-oculaire + hauteur mise au point.
\end{enumerate}
\cle{Simplification Probatoire :} On demande souvent seulement le diamètre du secondaire pour l'axe (vignettage axial nul) : $D_{\text{sec}} = D_1 \frac{f'_1 - d}{f'_1}$. L'obstruction linéaire $\epsilon = \frac{f'_1 - d}{f'_1}$.
\end{methode}

\begin{exoresolu}[Dimensionnement miroir secondaire et obstruction -- Type Probatoire]
\textbf{Énoncé :} Un télescope de Newton a un miroir primaire $D_1 = \SI{250}{mm}$, $f'_1 = \SI{1250}{mm}$ ($f/D=5$). Le miroir secondaire plan est placé à $d = \SI{1100}{mm}$ du sommet du primaire. Le porte-oculaire a un diamètre intérieur $D_{\text{oc}} = \SI{31,75}{mm}$ (standard 1,25'') et son plan focal est à $L_{\text{oc}} = \SI{80}{mm}$ au-dessus du plan du secondaire.
\begin{enumerate}
    \item Calculer le diamètre minimal $b$ (axe mineur) du miroir secondaire elliptique pour éclairer pleinement l'axe optique (pas de vignettage axial).
    \item Calculer le diamètre minimal $b_{\text{champ}}$ pour éclairer pleinement le bord du champ d'un oculaire 31,75 mm (diamètre champ apparent $50^\circ$, focale $\SI{25}{mm}$ $\rightarrow$ champ réel $\approx \SI{22}{mm}$). On prendra $h_{\text{champ}} = \SI{11}{mm}$.
    \item Déduire l'obstruction linéaire $\epsilon$ et surfacique $\tau$ (approximée disque diamètre $b_{\text{champ}}$).
    \item Vérifier que la distance $f'_1 - d$ permet la mise au point avec cet oculaire (focale $\SI{25}{mm}$).
\end{enumerate}
\tcblower\textbf{Solution détaillée :}
\begin{enumerate}
    \item \textbf{Vignettage axial (cône central) :}
    Le cône convergent a un angle $\alpha$ tel que $\tan\alpha = \frac{D_1/2}{f'_1} = \frac{125}{1250} = 0,1$.
    À la distance $d = \SI{1100}{mm}$ du sommet, le rayon du cône est :
    \[
    r_{\text{axe}} = \frac{D_1}{2} \frac{f'_1 - d}{f'_1} = 125 \times \frac{1250 - 1100}{1250} = 125 \times \frac{150}{1250} = 125 \times 0,12 = \SI{15}{mm}.
    \]
    Diamètre minimal axe mineur pour l'axe : $b_{\text{axe}} = 2 r_{\text{axe}} = \mathbf{\SI{30,0}{mm}}$.
    \item \textbf{Vignettage bord de champ (illumination plein champ) :}
    Le rayon principal du bord de champ (hauteur $h_{\text{champ}} = \SI{11}{mm}$ au plan focal) passe par le centre du secondaire (pupille d'entrée).
    Au niveau du secondaire (distance $f'_1 - d = \SI{150}{mm}$ \emph{avant} le plan focal), la hauteur de ce rayon par rapport à l'axe est (triangles semblables) :
    \[
    y_{\text{champ}} = h_{\text{champ}} \frac{f'_1 - d}{f'_1} = 11 \times \frac{150}{1250} = 11 \times 0,12 = \SI{1,32}{mm}.
    \]
    Le rayon du cône axial au secondaire est $\SI{15}{mm}$. Le rayon total nécessaire pour capter le bord de champ est la somme (cône axial décentré) :
    \[
    r_{\text{total}} = r_{\text{axe}} + y_{\text{champ}} = 15 + 1,32 = \SI{16,32}{mm}.
    \]
    Diamètre mineur requis : $b_{\text{champ}} = 2 r_{\text{total}} = \mathbf{\SI{32,6}{mm}}$.
    (Axe majeur $a = b\sqrt{2} \approx \SI{46,1}{mm}$).
    \item \textbf{Obstruction :} On prend $b = \SI{32,6}{mm}$ (standard commercial souvent 33 ou 35 mm).
    \begin{itemize}
        \item Obstruction linéaire : $\epsilon = \frac{b}{D_1} = \frac{32,6}{250} = \mathbf{\num{0,13}}$ (13\%).
        \item Obstruction surfacique (approx disque) : $\tau = \epsilon^2 = \num{0,13}^2 = \mathbf{\num{0,017}}$ (1,7\%).
        \item Obstruction surfacique exacte (ellipse) : $\tau_{\text{ell}} = \frac{ab}{D_1^2} = \frac{b^2\sqrt{2}}{D_1^2} = \sqrt{2} \epsilon^2 \approx 1,414 \times \num{0,017} = \mathbf{\num{0,024}}$ (2,4\%).
    \end{itemize}
    \textbf{Très faible obstruction} (typique f/5 avec secondaire petit). Excellent contraste.
    \item \textbf{Mise au point :}
    Distance secondaire - plan focal primaire : $f'_1 - d = 1250 - 1100 = \SI{150}{mm}$.
    L'oculaire (focale $\SI{25}{mm}$) a son foyer objet à $\SI{25}{mm}$ de son centre optique.
    Le plan focal de l'oculaire doit coïncider avec le plan focal du primaire.
    Distance secondaire - centre oculaire = $150 - 25 = \SI{125}{mm}$.
    Donnée : $L_{\text{oc}} = \SI{80}{mm}$ (distance secondaire - plan focal oculaire ? Non, énoncé dit "plan focal à 80 mm au-dessus du plan du secondaire").
    Si le plan focal de l'oculaire est à $\SI{80}{mm}$ du secondaire, et que le foyer du primaire est à $\SI{150}{mm}$ du secondaire $\rightarrow$ \textbf{Incompatibilité}. Le foyer du primaire est \textbf{trop loin} (150 mm vs 80 mm).
    \textbf{Correction :} Soit l'énoncé implique que le porte-oculaire rentre de $\SI{70}{mm}$ dans le tube (rare), soit il faut déplacer le miroir primaire vers l'avant (collimation), soit choisir un oculaire à longue distance de mise au point (ex: 2 pouces, ou barlow).
    \textbf{Conclusion :} Avec ces cotes standards, la mise au point n'est \textbf{pas possible} sans modifier la position du miroir primaire (le reculer de $\approx \SI{70}{mm}$) ou utiliser un correcteur/barlow.
\end{enumerate}
\end{exoresolu}

\begin{attention}[Les 5 erreurs qui coûtent des points au Probatoire]
\begin{enumerate}
    \item \textbf{Confondre "Puissance" (Vergence en Dioptries) et "Grossissement" (Sans unité).}
    \emph{Erreur :} Écrire $P = 1/f'_1 + 1/f'_2$ pour un télescope.
    \emph{Correction :} Le télescope afocal a $P=0\,\delta$. Le grossissement est $G = f'_1/f'_2$. La "puissance" au sens optique (vergence) est nulle. La "puissance" au sens commercial = grossissement.
    \item \textbf{Oublier l'inversion de l'image (Renversée).}
    \emph{Erreur :} Dire "image finale droite" ou "image virtuelle, droite".
    \emph{Correction :} Miroir primaire $\rightarrow$ 1 inversion. Oculaire (loupe) $\rightarrow$ 1 inversion. Total = \textbf{Image renversée} (rotation $180^\circ$). Citer : "Image finale virtuelle, à l'infini, \textbf{renversée} par rapport à l'objet".
    \item \textbf{Tracer le miroir secondaire comme une lentille ou un miroir convergent.}
    \emph{Erreur :} Faire converger les rayons \emph{après} le miroir secondaire vers un nouveau foyer.
    \emph{Correction :} Miroir \textbf{PLAN}. Il ne change \textbf{rien} à la convergence du faisceau (ni focale, ni taille image). Il ne fait que \textbf{dévié de 90°}. L'image intermédiaire $A'B'$ se forme \emph{au même endroit} (plan focal du primaire) mais décalée latéralement.
    \item \textbf{Mauvaise gestion des signes (Conjugaison / Grossissement).}
    \emph{Erreur :} $G = -f'_1/f'_2$ (signe moins) ou calcul de $p'$ faux pour mise au point proche.
    \emph{Correction :} Grossissement \textbf{angulaire} $G = f'_1/f'_2$ (valeur positive, le signe gère l'inversion). Pour la mise au point : utiliser la relation de Descartes \textbf{rigoureuse avec signes} ($1/f' = 1/p' - 1/p$, $p<0$ objet réel, $p'>0$ image réelle, $p'<0$ image virtuelle). Vérifier l'ordre de grandeur ($\Delta x \approx f'^2/d$).
    \item \textbf{Négliger l'obstruction centrale dans la comparaison Newton vs Lunette.}
    \emph{Erreur :} Dire "Newton meilleur que lunette car pas de chromatisme et plus grand diamètre" sans mentionner l'obstruction.
    \emph{Correction :} Argumentation équilibrée : Newton = \textbf{+ Ouverture/Prix, + Pas chromatisme} MAIS \textbf{- Obstruction (contraste), - Collimation, - Encombrement}. Lunette = \textbf{+ Contraste, + Stable, - Chromatisme (sauf APO), - Prix/Diamètre}. Le choix dépend de l'usage (Planétaire/Contraste $\rightarrow$ Lunette ; Ciel profond/Prix $\rightarrow$ Newton).
\end{enumerate}
\end{attention}

\newpage

\coursec{Exercices}

\courssub{Je vérifie mes connaissances}

\begin{exercice}[QCM : Constitution et rôle des éléments]
Parmi les affirmations suivantes concernant le télescope de Newton, lesquelles sont exactes ?
\begin{enumerate}
    \item Le miroir primaire est un miroir plan.
    \item Le miroir secondaire est un miroir plan incliné à $45^\circ$ par rapport à l'axe optique.
    \item L'oculaire est une lentille convergente.
    \item L'image finale observée par l'œil est réelle et inversée par rapport à l'objet.
    \item Le grossissement $G$ s'exprime par $G = \dfrac{f'_{\text{oc}}}{F}$ où $F$ est la focale du miroir primaire.
\end{enumerate}
\end{exercice}

\begin{exercice}[Vrai ou Faux : Justifiez]
Déterminez si les affirmations suivantes sont vraies ou fausses. Justifiez votre réponse par une phrase ou une formule.
\begin{enumerate}
    \item Le télescope de Newton souffre d'aberration chromatique.
    \item La puissance $P$ d'un télescope s'exprime en dioptries (δ).
    \item Pour observer un objet à l'infini (étoiles), l'image intermédiaire formée par le miroir primaire se trouve dans le plan focal de l'oculaire.
    \item Augmenter le diamètre $D$ du miroir primaire augmente le grossissement de l'instrument.
    \item Le miroir secondaire sert à renvoyer l'image hors de l'axe du tube pour permettre l'observation sans masquer l'ouverture.
\end{enumerate}
\end{exercice}

\begin{exercice}[Calculs directs : Grossissement et puissance]
Un télescope de Newton possède un miroir primaire de distance focale $F = 1\,200\,\text{mm}$ et un oculaire de distance focale $f'_{\text{oc}} = 20\,\text{mm}$.
\begin{enumerate}
    \item Calculer le grossissement $G$ de cet instrument.
    \item Calculer sa puissance $P$ (en $\text{m}^{-1}$ ou $\delta$).
    \item Si l'on remplace l'oculaire par un modèle de $f'_{\text{oc}} = 10\,\text{mm}$, que deviennent $G$ et $P$ ?
\end{enumerate}
\end{exercice}

\begin{exercice}[Définitions et grandeurs caractéristiques]
Complétez le tableau ci-dessous en associant chaque grandeur à sa définition, sa formule (si applicable) et son unité SI.
\begin{center}
\begin{tabularx}{\linewidth}{|l|X|X|X|}\hline
\rowcolor{popblueL}
\textbf{Grandeur} & \textbf{Définition / Rôle} & \textbf{Formule} & \textbf{Unité SI} \\ \hline
Distance focale du miroir primaire $F$ & & & \\ \hline
Distance focale de l'oculaire $f'_{\text{oc}}$ & & & \\ \hline
Grossissement $G$ & & & \\ \hline
Puissance $P$ & & & \\ \hline
Diamètre du miroir $D$ & & & \\ \hline
Pouvoir séparateur (limite de diffraction) $\alpha_{\min}$ & & $\alpha_{\min} \approx 1,22\,\dfrac{\lambda}{D}$ & \\ \hline
\end{tabularx}
\end{center}
\end{exercice}

\courssub{Je m'entraîne}

\begin{exercice}[Tracé de la marche d'un faisceau — Réglage à l'infini]
On considère un télescope de Newton de focale $F = 600\,\text{mm}$ et d'oculaire $f'_{\text{oc}} = 15\,\text{mm}$, réglé pour l'observation d'un objet à l'infini (étoile). Un faisceau parallèle incident forme un angle $\alpha = 1^\circ$ avec l'axe optique.
\begin{enumerate}
    \item Sur une feuille, reproduire à l'échelle $1/5$ (1 cm pour 100 mm) le schéma de l'instrument : miroir primaire (représenté par sa lentille équivalente convergente de focale $F$), miroir secondaire (plan à $45^\circ$), oculaire. Positionner correctement les plans focaux.
    \item Tracer la marche de deux rayons incidents parallèles (bord supérieur et bord inférieur du faisceau) à travers tout l'instrument jusqu'à l'œil.
    \item Mesurer graphiquement l'angle $\alpha'$ sous lequel l'œil voit l'image finale. Vérifier que $G = \dfrac{\alpha'}{\alpha} \approx -\dfrac{F}{f'_{\text{oc}}}$.
\end{enumerate}
\end{exercice}

\begin{exercice}[Mise au point : Lune vs Objet terrestre]
Un astronome amateur observe la Lune (objet à l'infini pratique) avec son télescope de Newton ($F = 800\,\text{mm}$, $f'_{\text{oc}} = 10\,\text{mm}$). L'instrument est réglé à l'infini (image finale à l'infini). Il souhaite ensuite observer un arbre situé à $d = 500\,\text{m}$.
\begin{enumerate}
    \item Rappeler la position de l'image intermédiaire $A'B'$ formée par le miroir primaire dans les deux cas (Lune et arbre). Calculer l'écart $\Delta x$ entre ces deux positions d'image.
    \item Pour obtenir une image finale nette à l'infini pour l'arbre, dans quel sens (vers le miroir primaire ou vers l'extérieur) et de quelle valeur faut-il déplacer l'oculaire ? Justifier par la relation de conjugaison de l'oculaire.
    \item Faire un schéma annoté montrant les deux positions de l'oculaire.
\end{enumerate}
\end{exercice}

\begin{exercice}[Luminosité et diamètre : Comparaison avec une lunette]
On compare deux instruments de même focale $F = f'_1 = 900\,\text{mm}$ et de même oculaire $f'_{\text{oc}} = f'_2 = 18\,\text{mm}$ :
\begin{itemize}
    \item Lunette astronomique : objectif de diamètre $D_L = 50\,\text{mm}$.
    \item Télescope de Newton : miroir primaire de diamètre $D_T = 114\,\text{mm}$.
\end{itemize}
\begin{enumerate}
    \item Calculer le grossissement $G$ de chaque instrument. Sont-ils égaux ?
    \item Calculer le diamètre de la pupille de sortie $D'_S$ pour chaque instrument (pupille de l'œil supposée $D_{\text{œil}} = 6\,\text{mm}$). Lequel est le plus lumineux pour l'observation d'objets diffus (nébuleuses) ?
    \item Calculer le pouvoir séparateur théorique $\alpha_{\min}$ (en secondes d'arc) pour $\lambda = 550\,\text{nm}$ pour chaque instrument. Lequel montre le plus de détails sur les planètes ?
    \item Citer deux avantages pratiques du télescope de Newton par rapport à la lunette pour des diamètres supérieurs à $100\,\text{mm}$.
\end{enumerate}
\end{exercice}

\begin{exercice}[Dimensionnement d'un oculaire pour Saturne]
On souhaite observer les anneaux de Saturne avec un grossissement $G = 120$ à l'aide d'un télescope de Newton dont le miroir primaire a une focale $F = 1\,200\,\text{mm}$ et un diamètre $D = 150\,\text{mm}$.
\begin{enumerate}
    \item Calculer la distance focale $f'_{\text{oc}}$ de l'oculaire à utiliser.
    \item Le grossissement maximal utile (ou optimal) est souvent estimé par $G_{\max} \approx 2 \times D(\text{mm})$ ou limité par la diffraction à $G_{\text{diff}} \approx \dfrac{D}{0,7\,\text{mm}}$. Vérifier si le grossissement $G = 120$ est raisonnable pour ce diamètre.
    \item Si l'oculaire disponible le plus court a une focale de $4\,\text{mm}$, quel est le grossissement maximal atteignable ? Est-il utile ?
\end{enumerate}
\end{exercice}

\courssub{Je me prépare au Probatoire}

\begin{exercice}[Sujet type Probatoire : Étude complète d'un télescope de Newton]
\textbf{Contexte :} L'Observatoire du Mont Fébé (Yaoundé) dispose d'un télescope de Newton pour l'observation planétaire. Le miroir primaire a un diamètre $D = 200\,\text{mm}$ et une distance focale $F = 1\,000\,\text{mm}$. On dispose d'une série d'oculaires : $f'_{\text{oc}} \in \{25\,\text{mm}; 10\,\text{mm}; 6\,\text{mm}; 4\,\text{mm}\}$.
\begin{enumerate}
    \item \textbf{Description et schéma} : Faire un schéma annoté du télescope de Newton. Indiquer le miroir primaire (concave), le miroir secondaire (plan à $45^\circ$), l'oculaire, l'axe optique, les foyers $F_1$ (miroir) et $F_2$ (oculaire), l'image intermédiaire $A'B'$ et l'image finale $A''B''$ pour un objet à l'infini. Préciser la nature (réelle/virtuelle) et l'orientation de chaque image.
    \item \textbf{Grossissements disponibles} : Calculer les quatre grossissements $G$ obtenus avec les oculaires fournis. Les classer par ordre croissant.
    \item \textbf{Pupille de sortie} : Pour chaque oculaire, calculer le diamètre de la pupille de sortie $D'_S$. Sachant que la pupille de l'œil dilatée la nuit vaut environ $7\,\text{mm}$, quel(s) oculaire(s) donnent une image pleinement exploitée par l'œil ? Quel oculaire donne une image "sombre" (pupille de sortie $> 7\,\text{mm}$) ?
    \item \textbf{Pouvoir séparateur} : Calculer la limite de diffraction $\alpha_{\min}$ (en secondes d'arc) pour $\lambda = 550\,\text{nm}$. L'instrument peut-il séparer deux étoiles doubles dont la séparation angulaire est $0,6''$ ?
    \item \textbf{Mise au point sur Jupiter} : Jupiter est à l'infini. On utilise l'oculaire $f'_{\text{oc}} = 10\,\text{mm}$. L'image intermédiaire se forme dans le plan focal du miroir. Où doit-on placer l'oculaire par rapport à cette image intermédiaire pour que l'image finale soit à l'infini (œil au repos) ? Si l'observateur est myope et retire ses lunettes (myopie $-2\,\delta$), dans quel sens doit-il déplacer l'oculaire pour voir net sans lunettes ? Justifier.
\end{enumerate}
\end{exercice}

\begin{exercice}[Sujet type Probatoire : Comparaison Lunette vs Télescope pour l'astrophotographie]
\textbf{Contexte :} Un club d'astronomie du Lycée de Nkolbisson hésite entre l'achat d'une lunette apochromatique (objectif : $D_L = 80\,\text{mm}$, $f'_L = 480\,\text{mm}$) et un télescope de Newton (miroir : $D_T = 150\,\text{mm}$, $F_T = 750\,\text{mm}$) pour faire de l'astrophotographie du ciel profond (nébuleuses, galaxies) avec un capteur APS-C (diagonale $d = 28\,\text{mm}$).
\begin{enumerate}
    \item \textbf{Rapport d'ouverture (f/D)} : Calculer le rapport d'ouverture (nombre $f$) des deux instruments. Lequel est le plus "rapide" (temps de pose plus court pour même signal) ?
    \item \textbf{Champ réel} : Avec une caméra dont les pixels font $3,76\,\mu\text{m}$, calculer le champ angulaire couvert par le capteur (en degrés) pour chaque instrument. (Indication : $\text{champ} \approx \dfrac{\text{taille capteur}}{\text{focale}}$ en radians).
    \item \textbf{Échantillonnage} : Calculer la résolution angulaire par pixel (en $''/\text{pixel}$) pour chaque instrument. Lequel est le mieux échantillonné pour du "Lucky Imaging" planétaire (cible $\approx 0,2''/\text{pix}$) ? Lequel pour du ciel profond (cible $\approx 1''$ à $2''/\text{pix}$) ?
    \item \textbf{Obstruction centrale} : Le télescope de Newton a une obstruction centrale (miroir secondaire) de $40\,\text{mm}$ de diamètre. Calculer le rapport d'obstruction $\varepsilon = \dfrac{d_{\text{secondaire}}}{D_T}$. Quel est l'impact principal sur le contraste planétaire et sur la luminosité effective ?
    \item \textbf{Conclusion argumentée} : En synthétisant les résultats, quel instrument recommanderiez-vous pour : (a) la planétaire haute résolution, (b) le ciel profond large champ ? Justifier en 3 arguments physiques pour chaque cas.
\end{enumerate}
\end{exercice}

\begin{exercice}[Sujet type Probatoire : Problème de mise au point et aberrations]
\textbf{Contexte :} Un élève construit un télescope de Newton "maison" avec un miroir primaire récupéré ($F = 900\,\text{mm}$, $D = 114\,\text{mm}$, rapport $f/D = 7,9$) et un oculaire Plössl $f'_{\text{oc}} = 20\,\text{mm}$. Il observe la Lune (distance $3,84 \times 10^8\,\text{m}$, diamètre $3\,474\,\text{km}$).
\begin{enumerate}
    \item \textbf{Taille de l'image intermédiaire} : Calculer le diamètre $h'$ de l'image de la Lune formée par le miroir primaire dans son plan focal. Exprimer en mm.
    \item \textbf{Grossissement et champ apparent} : Calculer le grossissement $G$. Si le champ apparent de l'oculaire est $50^\circ$, quel est le champ réel (en degrés) ? La Lune tient-elle entièrement dans le champ ? (Diamètre apparent Lune $\approx 0,5^\circ$).
    \item \textbf{Erreur de mise au point} : L'élève n'arrive pas à faire la mise au point : l'image reste floue quel que soit le déplacement de l'oculaire. Il mesure la distance entre le sommet du miroir et le plan focal de l'oculaire (rentré au maximum) : $L = 880\,\text{mm}$.
        \begin{enumerate}
            \item Quelle devrait être la distance théorique $L_{\text{th}}$ entre le sommet du miroir et le plan focal de l'oculaire pour un réglage à l'infini (miroir secondaire plan supposé au niveau du foyer du primaire) ?
            \item En déduire l'erreur de construction (tube trop court ou trop long ?). Proposer une solution mécanique simple pour corriger sans recouper le tube.
        \end{enumerate}
    \item \textbf{Aberration de coma} : Le miroir primaire est un paraboloïde. Pour un objet hors axe à un angle $\theta = 1^\circ$, la tache de coma a une taille angulaire approximative $\delta_{\text{coma}} \approx \dfrac{3}{16} \dfrac{\theta}{(f/D)^2}$ (en radians). Calculer $\delta_{\text{coma}}$ en minutes d'arc. Comparer à la taille de la tache d'Airy (limite de diffraction). L'aberration de coma est-elle négligeable au centre du champ ? Au bord d'un capteur APS-C (demi-diagonale $14\,\text{mm}$, angle $\theta_{\max} \approx 0,9^\circ$) ?
    \item \textbf{Correcteur de coma} : On insère un correcteur de coma (lentille convergente de faible puissance) juste devant l'oculaire. Quel est l'effet sur le grossissement effectif et sur le rapport $f/D$ effectif ? Pourquoi est-il indispensable pour la photographie grand champ avec un Newton $f/5$ ou moins ?
\end{enumerate}
\end{exercice}

\newpage

\coursec{Corrigés des exercices}

\coursec{Exercices d'application et sujets type Probatoire — Télescope de Newton}

\begin{corrige}[Exercice 1 — QCM : Constitution et rôle des éléments]
\textbf{Affirmations exactes : 2 et 3.}
\begin{enumerate}
    \item \textbf{Fausse.} Le miroir primaire est un miroir \cle{concave} (parabolique en pratique, sphérique en première approximation) : c'est lui qui collecte la lumière et forme l'image intermédiaire réelle en son foyer. Un miroir plan ne formerait aucune image convergente.
    \item \textbf{Vraie.} Le miroir secondaire est un petit miroir \cle{plan} incliné à $45^\circ$ par rapport à l'axe optique : il dévie le faisceau de $90^\circ$ vers l'oculaire placé sur le côté du tube.
    \item \textbf{Vraie.} L'oculaire est une lentille \cle{convergente} de courte focale qui joue le rôle de loupe : il transforme l'image intermédiaire en image finale agrandie.
    \item \textbf{Fausse.} En réglage afocal, l'image intermédiaire est dans le plan focal objet de l'oculaire : l'image finale est donc \cle{virtuelle}, rejetée à l'infini (et inversée par rapport à l'objet). Une image réelle ne peut pas être observée directement par l'œil.
    \item \textbf{Fausse.} La formule correcte est $G = -\dfrac{F}{f'_{\text{oc}}}$ (le signe $-$ traduit l'inversion ; en valeur absolue $G = F/f'_{\text{oc}}$). La formule proposée est l'inverse.
\end{enumerate}
\textbf{Points de vigilance :} bien distinguer le rôle des trois éléments (primaire = objectif collecteur, secondaire = renvoi, oculaire = loupe) et la nature virtuelle de l'image finale.
\end{corrige}

\begin{corrige}[Exercice 2 — Vrai ou Faux : Justifiez]
\begin{enumerate}
    \item \textbf{Faux.} La réflexion sur un miroir ne dépend pas de la longueur d'onde (contrairement à la réfraction dans une lentille, dont l'indice varie avec $\lambda$). Le télescope de Newton est donc \cle{exempt d'aberration chromatique} (seul l'oculaire, de faible focale, peut en introduire une très légère).
    \item \textbf{Vrai.} La puissance $P = \dfrac{\alpha'}{\alpha}$ s'exprime en radians par mètre, c'est-à-dire en dioptries ($\delta$). Pour un instrument afocal, $P = \dfrac{1}{f'_{\text{oc}}}$ avec $f'_{\text{oc}}$ en mètres.
    \item \textbf{Vrai.} Pour un objet à l'infini, le miroir primaire forme l'image intermédiaire en son plan focal image. Pour que l'œil observe sans accommoder (image finale à l'infini), cette image intermédiaire doit coïncider avec le \cle{foyer objet de l'oculaire} : c'est le réglage \cle{afocal}.
    \item \textbf{Faux.} Le grossissement $G = F/f'_{\text{oc}}$ ne dépend pas de $D$. Augmenter $D$ augmente la \cle{luminosité} (flux collecté proportionnel à $D^2$) et améliore le \cle{pouvoir séparateur} ($\alpha_{\min} \approx 1,22\,\lambda/D$ diminue), mais pas le grossissement.
    \item \textbf{Vrai.} Sans miroir secondaire, l'observateur devrait placer sa tête dans le faisceau incident, devant l'ouverture du tube, et masquerait une partie de la lumière. Le miroir plan à $45^\circ$ renvoie l'image latéralement, hors du chemin incident.
\end{enumerate}
\end{corrige}

\begin{corrige}[Exercice 3 — Calculs directs : Grossissement et puissance]
\textbf{Données :} $F = \SI{1200}{mm}$, $f'_{\text{oc}} = \SI{20}{mm}$.
\begin{enumerate}
    \item Grossissement (valeur absolue) :
    \[ G = \frac{F}{f'_{\text{oc}}} = \frac{1200}{20} = 60. \]
    L'instrument grossit 60 fois (image inversée : $G = -60$ en valeur algébrique).
    \item Puissance de l'instrument afocal : $P = \dfrac{1}{f'_{\text{oc}}}$ avec $f'_{\text{oc}}$ en mètres :
    \[ P = \frac{1}{\SI{0,020}{m}} = \SI{50}{\delta}. \]
    \item Avec $f'_{\text{oc}} = \SI{10}{mm} = \SI{0,010}{m}$ :
    \[ G = \frac{1200}{10} = 120 \qquad ; \qquad P = \frac{1}{0,010} = \SI{100}{\delta}. \]
    Diviser la focale de l'oculaire par 2 double le grossissement et la puissance.
\end{enumerate}
\textbf{Point de vigilance :} convertir les focales en mètres pour la puissance en dioptries ; le grossissement, rapport de deux longueurs, est sans unité.
\end{corrige}

\begin{corrige}[Exercice 4 — Définitions et grandeurs caractéristiques]
\begin{center}
\begin{tabularx}{\linewidth}{|l|X|X|X|}\hline
\rowcolor{popblueL}
\textbf{Grandeur} & \textbf{Définition / Rôle} & \textbf{Formule} & \textbf{Unité SI} \\ \hline
Distance focale du miroir primaire $F$ & Distance du sommet du miroir concave à son foyer ; fixe la taille de l'image intermédiaire & $F = \dfrac{R}{2}$ (miroir sphérique) & mètre (m) \\ \hline
Distance focale de l'oculaire $f'_{\text{oc}}$ & Distance focale de la lentille convergente servant de loupe & $f'_{\text{oc}} = \dfrac{1}{C_{\text{oc}}}$ & mètre (m) \\ \hline
Grossissement $G$ & Rapport du diamètre apparent de l'image à celui de l'objet & $G = \dfrac{\alpha'}{\alpha} = -\dfrac{F}{f'_{\text{oc}}}$ & sans unité \\ \hline
Puissance $P$ & Aptitude à agrandir le diamètre apparent & $P = \dfrac{1}{f'_{\text{oc}}}$ & dioptrie ($\delta$) \\ \hline
Diamètre du miroir $D$ & Diamètre utile de l'ouverture collectrice ; fixe luminosité et résolution & --- & mètre (m) \\ \hline
Pouvoir séparateur $\alpha_{\min}$ & Plus petite séparation angulaire discernable (diffraction) & $\alpha_{\min} \approx 1,22\,\dfrac{\lambda}{D}$ & radian (rad) \\ \hline
\end{tabularx}
\end{center}
\textbf{Point de vigilance :} le grossissement est un rapport d'angles (sans dimension) ; la puissance s'exprime en dioptries ; le pouvoir séparateur en radians (convertir en secondes d'arc : $1\ \text{rad} = 206\,265''$).
\end{corrige}

\begin{corrige}[Exercice 5 — Tracé de la marche d'un faisceau — Réglage à l'infini]
\begin{enumerate}
    \item À l'échelle $1/5$ : $F = \SI{600}{mm} \rightarrow \SI{12}{cm}$ ; $f'_{\text{oc}} = \SI{15}{mm} \rightarrow \SI{3}{mm}$. On représente le miroir primaire par une lentille convergente équivalente de focale $F$ (le miroir « replie » le trajet, mais le tracé déplié est équivalent). Le foyer image $F_1$ du primaire est à $\SI{12}{cm}$ de celui-ci ; le foyer objet $F_2$ de l'oculaire est placé en $F_1$ (réglage afocal), donc l'oculaire est à $F + f'_{\text{oc}} = \SI{615}{mm}$ du primaire, soit $\SI{12,3}{cm}$ sur le schéma. Le miroir secondaire à $45^\circ$ est placé juste avant $F_1$ (il ne change pas la position de l'image, il la dévie).
    \item Tracé : les deux rayons incidents parallèles entre eux (angle $\alpha = 1^\circ$ avec l'axe) convergent, après réflexion sur le primaire, en un même point $B'$ du plan focal, situé à la hauteur $y = F\tan\alpha \approx 600 \times 0,0175 \approx \SI{10,5}{mm}$ (soit $\SI{2,1}{mm}$ sur le schéma) au-dessus de l'axe. Après le miroir secondaire, ces rayons arrivent sur l'oculaire ; comme $B'$ est au foyer objet de l'oculaire, ils ressortent \textbf{parallèles entre eux}, sous l'angle $\alpha'$ tel que $\tan\alpha' = \dfrac{y}{f'_{\text{oc}}}$.
    \item Vérification :
    \[ \tan\alpha' = \frac{F\tan\alpha}{f'_{\text{oc}}} = \frac{600 \times \tan 1^\circ}{15} = 40 \times 0,0175 = 0,698 \quad\Rightarrow\quad \alpha' \approx 34,9^\circ. \]
    Le rapport exact est $\dfrac{\tan\alpha'}{\tan\alpha} = \dfrac{F}{f'_{\text{oc}}} = 40$. La formule $G = \alpha'/\alpha$ (approximation des petits angles) donnerait $\alpha' \approx 40^\circ$ : elle devient imprécise ici car $\alpha'$ n'est plus petit. On retient :
    \[ G = \frac{\tan\alpha'}{\tan\alpha} \approx \frac{\alpha'}{\alpha} = \frac{F}{f'_{\text{oc}}} = 40 \quad (\text{signe } - \text{ : image inversée}). \]
\end{enumerate}
\textbf{Points de vigilance :} les rayons issus d'un point à l'infini (faisceau parallèle) convergent en un point du \emph{plan} focal, pas nécessairement au foyer ; à la sortie de l'oculaire, le faisceau ressort parallèle (image à l'infini) : c'est la signature du réglage afocal.
\end{corrige}

\begin{corrige}[Exercice 6 — Mise au point : Lune vs Objet terrestre]
\textbf{Données :} $F = \SI{800}{mm}$, $f'_{\text{oc}} = \SI{10}{mm}$, $d = \SI{500}{m}$.
\begin{enumerate}
    \item \textbf{Lune (objet à l'infini) :} l'image intermédiaire se forme dans le plan focal image du miroir : $x'_{\text{Lune}} = F = \SI{800}{mm}$ du sommet du miroir.

    \textbf{Arbre à $d = \SI{500}{m}$ :} relation de conjugaison (le miroir concave se comporte comme une lentille convergente de focale $F$) :
    \[ \frac{1}{F} = \frac{1}{d} + \frac{1}{x'_{\text{arbre}}} \;\Rightarrow\; x'_{\text{arbre}} = \frac{F\,d}{d - F} = \frac{0,800 \times 500}{500 - 0,800} = \frac{400}{499,2} \approx \SI{0,80128}{m} = \SI{801,28}{mm}. \]
    Écart : $\Delta x = x'_{\text{arbre}} - x'_{\text{Lune}} \approx 801,28 - 800 = \SI{1,28}{mm}$. L'image de l'objet rapproché se forme \emph{plus loin} du miroir que le plan focal.
    \item Pour que l'image finale reste à l'infini, l'image intermédiaire doit rester au foyer objet de l'oculaire. Comme elle s'est éloignée du miroir de $\Delta x \approx \SI{1,28}{mm}$, il faut \textbf{éloigner l'oculaire du miroir primaire} (tirer le porte-oculaire vers l'extérieur) de $\Delta x \approx \SI{1,3}{mm}$. Justification : pour l'oculaire, l'objet (image intermédiaire) doit vérifier $\overline{O_{\text{oc}}A'} = -f'_{\text{oc}}$ ; si $A'$ recule de $\Delta x$, l'oculaire doit reculer d'autant.
    \item Schéma annoté : on dessine le miroir primaire à gauche, deux positions de $A'$ ($F_1$ pour la Lune, $F_1 + \SI{1,3}{mm}$ pour l'arbre), et les deux positions correspondantes de l'oculaire, la seconde décalée de $\Delta x$ vers l'extérieur (côté observateur).
\end{enumerate}
\textbf{Point de vigilance :} le déplacement de mise au point est de l'ordre du millimètre : c'est pourquoi les porte-oculaires ont une course de quelques centimètres et une crémaillère de précision.
\end{corrige}

\begin{corrige}[Exercice 7 — Luminosité et diamètre : Comparaison avec une lunette]
\textbf{Données :} $F = \SI{900}{mm}$, $f'_{\text{oc}} = \SI{18}{mm}$ pour les deux ; $D_L = \SI{50}{mm}$, $D_T = \SI{114}{mm}$.
\begin{enumerate}
    \item Grossissements :
    \[ G_L = \frac{900}{18} = 50 \qquad ; \qquad G_T = \frac{900}{18} = 50. \]
    \textbf{Oui, ils sont égaux} : le grossissement ne dépend que des focales, pas du diamètre.
    \item Pupille de sortie : $D'_S = \dfrac{D}{G}$.
    \[ D'_{S,L} = \frac{50}{50} = \SI{1,0}{mm} \qquad ; \qquad D'_{S,T} = \frac{114}{50} = \SI{2,28}{mm}. \]
    Les deux pupilles sont inférieures à la pupille de l'œil ($\SI{6}{mm}$) : toute la lumière entre dans l'œil dans les deux cas. Mais le télescope collecte $\left(\dfrac{114}{50}\right)^2 \approx 5,2$ fois plus de lumière ; à grossissement égal, l'image qu'il donne est environ 5 fois plus lumineuse par unité de surface angulaire. \textbf{Le télescope est le plus lumineux} pour les nébuleuses.
    \item Pouvoir séparateur : $\alpha_{\min} = 1,22\,\dfrac{\lambda}{D}$, avec $\lambda = \SI{550}{nm} = \SI{5,5e-7}{m}$ ; conversion : $1\ \text{rad} = 206\,265''$.
    \[ \alpha_{\min,L} = \frac{1,22 \times 5,5\times10^{-7}}{0,050} = 1,34\times10^{-5}\ \text{rad} \approx 2,8''. \]
    \[ \alpha_{\min,T} = \frac{1,22 \times 5,5\times10^{-7}}{0,114} = 5,9\times10^{-6}\ \text{rad} \approx 1,2''. \]
    \textbf{Le télescope montre le plus de détails} (pouvoir séparateur plus de deux fois meilleur).
    \item Avantages pratiques du Newton pour $D > \SI{100}{mm}$ :
    \begin{itemize}
        \item \textbf{Absence d'aberration chromatique} (la réflexion ne disperse pas la lumière), alors qu'un grand objectif à lentilles exige des doublets/triplets coûteux ;
        \item \textbf{Coût et masse} : un miroir de grand diamètre est bien moins cher et plus léger qu'une lentille de même diamètre, et le tube est plus court (le faisceau est replié) ;
        \item un miroir ne se déforme pas sous son propre poids comme une lentille montée par les bords.
    \end{itemize}
\end{enumerate}
\textbf{Point de vigilance :} à grossissement égal, c'est le \emph{diamètre} qui tranche la luminosité et la résolution — d'où l'intérêt majeur du télescope.
\end{corrige}

\begin{corrige}[Exercice 8 — Dimensionnement d'un oculaire pour Saturne]
\textbf{Données :} $F = \SI{1200}{mm}$, $D = \SI{150}{mm}$, $G$ visé $= 120$.
\begin{enumerate}
    \item Focale de l'oculaire :
    \[ f'_{\text{oc}} = \frac{F}{G} = \frac{1200}{120} = \SI{10}{mm}. \]
    \item Grossissements limites :
    \[ G_{\max} \approx 2 \times D(\text{mm}) = 2 \times 150 = 300 \qquad ; \qquad G_{\text{diff}} \approx \frac{D}{0,7\ \text{mm}} = \frac{150}{0,7} \approx 214. \]
    Le grossissement $G = 120$ est \textbf{inférieur aux deux limites} : il est parfaitement raisonnable pour ce diamètre. L'image restera contrastée et exploitable si la turbulence atmosphérique le permet.
    \item Avec $f'_{\text{oc}} = \SI{4}{mm}$ :
    \[ G = \frac{1200}{4} = 300. \]
    Ce grossissement atteint $G_{\max} = 300$ et dépasse nettement $G_{\text{diff}} \approx 214$ : au-delà de $G_{\text{diff}}$, l'image n'apporte \textbf{aucun détail supplémentaire} (la diffraction fixe la résolution), elle devient simplement plus grande, plus sombre et plus floue (« grossissement vide »). Il n'est donc \textbf{pas utile}, sauf nuits de turbulence exceptionnellement faible.
\end{enumerate}
\textbf{Point de vigilance :} le grossissement utile est borné par le diamètre (diffraction) et par la turbulence ; augmenter $G$ sans augmenter $D$ ne révèle pas plus de détails.
\end{corrige}

\begin{corrige}[Exercice 9 — Sujet type Probatoire : Étude complète d'un télescope de Newton]
\textbf{Données :} $D = \SI{200}{mm}$, $F = \SI{1000}{mm}$ ; oculaires $\{\SI{25}{mm} ; \SI{10}{mm} ; \SI{6}{mm} ; \SI{4}{mm}\}$.
\begin{enumerate}
    \item \textbf{Schéma} (à reproduire) : tube horizontal, miroir primaire concave au fond, miroir secondaire plan à $45^\circ$ sur l'axe avant le foyer, oculaire latéral. Pour un objet à l'infini : les rayons parallèles convergent vers l'\textbf{image intermédiaire $A'B'$}, \emph{réelle} et \emph{inversée}, formée dans le plan focal $F_1$ du miroir ; le secondaire la renvoie latéralement ; placée au foyer objet $F_2$ de l'oculaire, elle donne une \textbf{image finale $A''B''$}, \emph{virtuelle}, \emph{inversée} par rapport à l'objet, rejetée à l'infini.
    \item Grossissements $G = F/f'_{\text{oc}}$ :
    \begin{center}
    \begin{tabular}{|c|cccc|}\hline
    $f'_{\text{oc}}$ (mm) & 25 & 10 & 6 & 4 \\ \hline
    $G$ & 40 & 100 & 167 & 250 \\ \hline
    \end{tabular}
    \end{center}
    Ordre croissant : $40 < 100 < 167 < 250$.
    \item Pupille de sortie $D'_S = D/G$ :
    \begin{center}
    \begin{tabular}{|c|cccc|}\hline
    $f'_{\text{oc}}$ (mm) & 25 & 10 & 6 & 4 \\ \hline
    $D'_S$ (mm) & 5,0 & 2,0 & 1,2 & 0,8 \\ \hline
    \end{tabular}
    \end{center}
    Toutes les pupilles de sortie sont $\leq \SI{7}{mm}$ : \textbf{tous les oculaires} donnent un faisceau entièrement reçu par l'œil (image pleinement exploitée). \textbf{Aucun} oculaire ne donne une image « sombre » au sens pupille $> \SI{7}{mm}$ (cela arriverait pour $G < D/7 \approx 29$, par exemple un oculaire de $\SI{40}{mm}$).
    \item Limite de diffraction :
    \[ \alpha_{\min} = 1,22\,\frac{\lambda}{D} = \frac{1,22 \times 5,5\times10^{-7}}{0,200} = 3,36\times10^{-6}\ \text{rad} \approx 0,69''. \]
    Deux étoiles séparées de $0,6'' < 0,69''$ : \textbf{non}, l'instrument ne peut pas les séparer (critère de Rayleigh non satisfait).
    \item \textbf{Mise au point :} pour une image finale à l'infini (œil emmétrope au repos), le foyer objet de l'oculaire doit coïncider avec l'image intermédiaire : l'oculaire est placé à $f'_{\text{oc}} = \SI{10}{mm}$ de $A'B'$. Pour l'observateur myope ($-2\ \delta$), l'image finale doit être virtuelle, située à son punctum remotum, soit $d = 1/2 = \SI{0,50}{m}$ devant l'œil. L'objet (image intermédiaire) doit alors être légèrement \emph{à l'intérieur} du foyer objet de l'oculaire. Estimation du déplacement (relation de Newton $x x' = -f'^2$) : $\Delta x \approx f'^2_{\text{oc}} \times |C| = (0,010)^2 \times 2 = \SI{2e-4}{m} = \SI{0,2}{mm}$. Il faut donc \textbf{rapprocher l'oculaire du miroir primaire} d'environ $\SI{0,2}{mm}$.
\end{enumerate}
\textbf{Barème indicatif (10 pts) :} schéma annoté complet 2,5 ; grossissements 1,5 ; pupilles de sortie + conclusion 2 ; pouvoir séparateur + conclusion 2 ; mise au point + cas du myope 2.
\textbf{Points de vigilance :} citer la condition afocale (foyers $F_1$ et $F_2$ confondus) ; justifier la nature virtuelle de l'image finale ; conversions nm $\rightarrow$ m et rad $\rightarrow$ secondes d'arc ; pour le myope, expliciter que l'image finale doit se former au punctum remotum (et non à l'infini).
\end{corrige}

\begin{corrige}[Exercice 10 — Sujet type Probatoire : Comparaison Lunette vs Télescope pour l'astrophotographie]
\textbf{Données :} lunette $D_L = \SI{80}{mm}$, $f'_L = \SI{480}{mm}$ ; Newton $D_T = \SI{150}{mm}$, $F_T = \SI{750}{mm}$ ; capteur $d = \SI{28}{mm}$, pixel $p = \SI{3,76}{\micro m}$.
\begin{enumerate}
    \item Rapports d'ouverture :
    \[ N_L = \frac{f'_L}{D_L} = \frac{480}{80} = 6 \qquad ; \qquad N_T = \frac{F_T}{D_T} = \frac{750}{150} = 5. \]
    Plus $f/D$ est petit, plus l'instrument est « rapide » (éclairement du capteur proportionnel à $1/N^2$) : \textbf{le télescope ($f/5$) est le plus rapide} ; à temps de pose égal, il donne un signal $(6/5)^2 = 1,44$ fois supérieur.
    \item Champ angulaire $= \dfrac{d}{f}$ (en radians) :
    \[ \theta_L = \frac{28}{480} = 5,83\times10^{-2}\ \text{rad} \approx 3,34^\circ \qquad ; \qquad \theta_T = \frac{28}{750} = 3,73\times10^{-2}\ \text{rad} \approx 2,14^\circ. \]
    La lunette couvre un champ plus large (focale plus courte).
    \item Échantillonnage : $e = \dfrac{p}{f}$ en radians, converti en secondes d'arc ($\times 206\,265$) :
    \[ e_L = \frac{3,76\times10^{-6}}{0,480} = 7,83\times10^{-6}\ \text{rad} \approx 1,62''/\text{pix} \qquad ; \qquad e_T = \frac{3,76\times10^{-6}}{0,750} = 5,01\times10^{-6}\ \text{rad} \approx 1,03''/\text{pix}. \]
    Aucun n'atteint $0,2''/\text{pix}$ : pour le « Lucky Imaging » planétaire, le \textbf{télescope} s'en approche le mieux (et une lentille de Barlow $\times 5$ donnerait $\approx 0,21''/\text{pix}$). Pour le ciel profond (cible $1''$ à $2''/\text{pix}$), \textbf{les deux conviennent}, la lunette étant idéale en grand champ.
    \item Obstruction : $\varepsilon = \dfrac{40}{150} \approx 0,27$ (27 \% en diamètre). Conséquences : perte de lumière $\varepsilon^2 \approx 7\ \%$ (négligeable) ; mais transfert d'énergie du pic central de la tache d'Airy vers les anneaux $\Rightarrow$ \textbf{baisse du contraste} sur les détails planétaires fins.
    \item \textbf{Conclusion :}
    \begin{itemize}
        \item \textbf{(a) Planétaire haute résolution : le télescope.} (1) Diamètre presque double $\Rightarrow$ pouvoir séparateur $\alpha_{\min,T} \approx 0,92''$ contre $1,72''$ pour la lunette ; (2) focale plus longue $\Rightarrow$ meilleur échantillonnage ($1,03''/\text{pix}$, compatible Barlow) ; (3) absence d'aberration chromatique résiduelle, cruciale en haute résolution.
        \item \textbf{(b) Ciel profond large champ : la lunette.} (1) Champ plus large ($3,3^\circ$ contre $2,1^\circ$) adapté aux grandes nébuleuses ; (2) échantillonnage $1,62''/\text{pix}$ pile dans la cible $1$--$2''/\text{pix}$ ; (3) pas d'obstruction centrale ni de coma (optique à lentilles corrigée), étoiles ponctuelles jusqu'au bord et contraste maximal.
    \end{itemize}
\end{enumerate}
\textbf{Barème indicatif (10 pts) :} $f/D$ 1,5 ; champs 2 ; échantillonnages 2 ; obstruction 1,5 ; conclusion argumentée 3.
\textbf{Points de vigilance :} conserver les unités cohérentes (mm partout) ; justifier « rapide » par l'éclairement en $1/N^2$ ; la conclusion doit mobiliser des arguments \emph{physiques} chiffrés, pas des goûts.
\end{corrige}

\begin{corrige}[Exercice 11 — Sujet type Probatoire : Problème de mise au point et aberrations]
\textbf{Données :} $F = \SI{900}{mm}$, $D = \SI{114}{mm}$, $f/D = 7,9$, $f'_{\text{oc}} = \SI{20}{mm}$ ; Lune : $d_L = \SI{3,84e8}{m}$, $\varnothing_L = \SI{3,474e6}{m}$.
\begin{enumerate}
    \item Diamètre angulaire de la Lune : $\alpha = \dfrac{\varnothing_L}{d_L} = \dfrac{3,474\times10^{6}}{3,84\times10^{8}} = 9,05\times10^{-3}\ \text{rad}$. Taille de l'image focale :
    \[ h' = F\,\alpha = 0,900 \times 9,05\times10^{-3} = 8,14\times10^{-3}\ \text{m} \approx \SI{8,1}{mm}. \]
    \item Grossissement : $G = \dfrac{F}{f'_{\text{oc}}} = \dfrac{900}{20} = 45$. Champ réel :
    \[ \theta_{\text{réel}} = \frac{\theta_{\text{apparent}}}{G} = \frac{50^\circ}{45} \approx 1,11^\circ. \]
    La Lune ($0,5^\circ$) occupe moins de la moitié du champ : \textbf{elle tient entièrement dans le champ}.
    \item \textbf{Erreur de mise au point :}
        \begin{enumerate}
            \item Pour un réglage à l'infini, le plan focal de l'oculaire doit coïncider avec le foyer du miroir : la distance sommet du miroir $\rightarrow$ plan focal de l'oculaire vaut
            \[ L_{\text{th}} = F = \SI{900}{mm}. \]
            \item Or $L = \SI{880}{mm} < \SI{900}{mm}$ : le \textbf{tube est trop court de 20 mm} ; l'image focale se forme 20 mm \emph{au-delà} de la position maximale accessible de l'oculaire, d'où l'impossibilité de mise au point. Solutions mécaniques sans recouper le tube : reculer le miroir primaire dans son barillet (vis de collimation / rallonge de support), ou intercaler une \textbf{rallonge (tube d'extension)} entre le porte-oculaire et l'oculaire, ou remplacer le porte-oculaire par un modèle à plus longue course.
        \end{enumerate}
    \item \textbf{Coma :} $\theta = 1^\circ = 1,745\times10^{-2}\ \text{rad}$ ; $(f/D)^2 = 7,9^2 = 62,4$.
    \[ \delta_{\text{coma}} = \frac{3}{16}\,\frac{\theta}{(f/D)^2} = \frac{0,1875 \times 1,745\times10^{-2}}{62,4} = 5,24\times10^{-5}\ \text{rad} \approx 10,8''. \]
    Tache d'Airy (rayon de Rayleigh) : $\alpha_{\min} = 1,22\,\lambda/D = \dfrac{1,22 \times 5,5\times10^{-7}}{0,114} = 5,9\times10^{-6}\ \text{rad} \approx 1,2''$.
    La coma à $1^\circ$ de l'axe est environ \textbf{9 fois} le rayon de la tache de diffraction : elle est significative et dégrade notablement la qualité des étoiles en bord de champ. \textbf{Au centre} ($\theta = 0$), la coma est nulle par construction (miroir parabolique parfaitement stigmatique sur axe) : négligeable. \textbf{Au bord du capteur} ($\theta_{\max} = 0,9^\circ$) : $\delta_{\text{coma}} \approx 5,24\times10^{-5} \times 0,9 = 4,7\times10^{-5}\ \text{rad} \approx 9,7''$ : les étoiles présentent une queue de comète visible — aberration non négligeable en photographie.
    \item \textbf{Correcteur de coma :} la lentille de faible puissance allonge légèrement la focale effective (facteur typique $\times 1,1$) : le \textbf{grossissement effectif augmente} d'autant et le rapport effectif passe de $f/7,9$ à environ $f/8,7$. En échange, elle corrige la coma sur tout le champ. Elle est \textbf{indispensable} pour un Newton rapide ($f/5$ ou moins) car la coma croît comme $1/(f/D)^2$ : à $f/5$, $(f/D)^2 = 25$ et la tache de coma au même angle serait $62,4/25 \approx 2,5$ fois plus grande qu'à $f/7,9$, rendant les bords du capteur inutilisables sans correction.
\end{enumerate}
\textbf{Barème indicatif (10 pts) :} taille image 1,5 ; grossissement + champ 1,5 ; diagnostic de mise au point 2,5 ; calcul de coma + comparaison 2,5 ; correcteur 2.
\textbf{Points de vigilance :} exprimer $\theta$ en radians dans la formule de la coma ; comparer des grandeurs dans la même unité ($'$ et $''$) ; pour le diagnostic, conclure sur le signe de l'écart ($L < L_{\text{th}} \Rightarrow$ tube trop court) et proposer une solution qui \emph{augmente} le tirage arrière.
\end{corrige}

\newpage

\coursec{Fiche bilan}

\courssub{Carte mentale de la leçon}
\begin{popfigure}
\begin{center}\tcbox[colback=popblue,colframe=popblue,coltext=white,arc=9pt,boxsep=4pt,left=6pt,right=6pt]{\bfseries\small Télescope de Newton}\end{center}
\vspace{-2pt}
\begin{tcbraster}[raster columns=3,raster equal height=rows,raster column skip=6pt,raster row skip=6pt,enhanced,blankest]
\begin{tcolorbox}[enhanced,colback=popblue,colframe=popblue,coltext=white,boxrule=0.9pt,arc=3pt,left=4pt,right=4pt,top=3pt,bottom=3pt,boxsep=1pt,halign=left]\bfseries\scriptsize Description Miroir concave + miroir plan + oculaire\end{tcolorbox}
\begin{tcolorbox}[enhanced,colback=poporange,colframe=poporange,coltext=white,boxrule=0.9pt,arc=3pt,left=4pt,right=4pt,top=3pt,bottom=3pt,boxsep=1pt,halign=left]\bfseries\scriptsize Fonctionnement Image réelle au foyer du primaire\end{tcolorbox}
\begin{tcolorbox}[enhanced,colback=popgreen,colframe=popgreen,coltext=white,boxrule=0.9pt,arc=3pt,left=4pt,right=4pt,top=3pt,bottom=3pt,boxsep=1pt,halign=left]\bfseries\scriptsize Tracé rayons 3 rayons caractéristiques (parallèle, sommet, foyer)\end{tcolorbox}
\begin{tcolorbox}[enhanced,colback=poppurple,colframe=poppurple,coltext=white,boxrule=0.9pt,arc=3pt,left=4pt,right=4pt,top=3pt,bottom=3pt,boxsep=1pt,halign=left]\bfseries\scriptsize Grossissement $G = -\dfrac{f'_o}{f'_e}$ Puissance $P = \dfrac{1}{f'_e}$\end{tcolorbox}
\begin{tcolorbox}[enhanced,colback=poppink,colframe=poppink,coltext=white,boxrule=0.9pt,arc=3pt,left=4pt,right=4pt,top=3pt,bottom=3pt,boxsep=1pt,halign=left]\bfseries\scriptsize Mise au point Déplacement oculaire (foyers confondus)\end{tcolorbox}
\begin{tcolorbox}[enhanced,colback=popgold,colframe=popgold,coltext=white,boxrule=0.9pt,arc=3pt,left=4pt,right=4pt,top=3pt,bottom=3pt,boxsep=1pt,halign=left]\bfseries\scriptsize Intérêt vs lentilles Pas d'aberration chromatique\end{tcolorbox}
\end{tcbraster}

\legende{Synthèse visuelle du télescope de Newton : composants, principes, formules et avantages.}
\end{popfigure}

\courssub{Définitions clés}
\begin{definition}[Définitions clés]
\begin{itemize}
  \item \cle{Télescope de Newton} : Instrument d'optique à réflexion utilisant un \cle{miroir concave primaire} (grand diamètre, grande focale image $f'_o > 0$, souvent parabolique), un \cle{miroir plan secondaire} (incliné à 45° par rapport à l'axe optique) et une \cle{oculaire} (lentille mince convergente de courte focale image $f'_e > 0$).
  \item \cle{Image intermédiaire} : Image réelle, inversée par rapport à l'objet, formée au foyer image $F'_o$ du miroir primaire (avant réflexion sur le miroir plan). Elle sert d'objet réel pour l'oculaire.
  \item \cle{Image finale} : Image virtuelle, inversée par rapport à l'objet initial, située à l'infini (réglage normal, œil au repos) ou à la distance de vision distincte $D = \SI{25}{cm}$ (réglage pour œil accommodé).
  \item \cle{Mise au point} : Action de déplacer l'oculaire le long de l'axe optique pour que son foyer objet $F_e$ coïncide exactement avec l'image intermédiaire (foyer image $F'_o$ du primaire dévié par le miroir plan).
  \item \cle{Pupille d'entrée} : Le miroir primaire (diaphragme limitant le faisceau entrant). Son diamètre $D$ détermine le pouvoir séparateur et la luminosité.
  \item \cle{Pupille de sortie} : Image du miroir primaire donnée par l'oculaire. L'œil de l'observateur doit s'y placer pour voir tout le champ.
\end{itemize}
\end{definition}

\courssub{Formules à connaître par cœur}
\begin{propriete}[Formules à connaître par cœur]
\begin{center}
\begin{tabularx}{\linewidth}{|l|X|l|}\hline
\rowcolor{popblueL}
\textbf{Grandeur} & \textbf{Formule} & \textbf{Unité SI} \\ \hline
Grossissement (réglage normal) & $G = -\dfrac{f'_o}{f'_e}$ & Sans unité \\ \hline
Puissance commerciale (oculaire) & $P = \dfrac{1}{f'_e}$ ($f'_e$ en m) & Dioptrie (D) ou $\text{m}^{-1}$ \\ \hline
Diamètre utile (pupille d'entrée) & $D = \text{diamètre miroir primaire}$ & Mètre (m) \\ \hline
Pouvoir séparateur (Rayleigh) & $\theta_{\min} \approx 1,22\,\dfrac{\lambda}{D}$ & Radian (rad) \\ \hline
Luminosité (éclairement image) & $\propto \left(\dfrac{D}{f'_o}\right)^2 = \dfrac{1}{N^2}$ & -- ($N = f'_o/D$ : nombre d'ouverture) \\ \hline
\end{tabularx}
\end{center}
\emph{Note :} Le signe $-$ de $G$ indique que l'image finale est inversée par rapport à l'objet (télescope astronomique). La puissance $P$ caractérise l'oculaire seul, le grossissement $G$ dépend du couple (primaire, oculaire).
\end{propriete}

\courssub{Méthode : Tracer la marche d'un faisceau incident parallèle à l'axe}
\begin{methode}[Tracer la marche d'un faisceau incident parallèle à l'axe (réglage normal)]
\begin{enumerate}
  \item Tracer le miroir concave primaire (représenté par son sommet $S_o$ et son foyer image $F'_o$) ; placer le miroir plan à 45° \textbf{avant} $F'_o$ sur le trajet de la lumière.
  \item Tracer l'oculaire (lentille mince convergente, centre optique $O_e$, foyers $F_e$ et $F'_e$) de sorte que son foyer objet $F_e$ coïncide avec l'image intermédiaire $F'_o$ (réglage normal : foyers confondus).
  \item \textbf{Rayon 1 (incident parallèle à l'axe)} : Réfléchi par le primaire en passant par $F'_o$, dévié à 90° par le miroir plan, traverse l'oculaire et en sort \textbf{parallèle à l'axe} (passe par $F'_e$).
  \item \textbf{Rayon 2 (incident sur le sommet $S_o$ du primaire)} : Réfléchi symétriquement par rapport à l'axe (loi de la réflexion), dévié par le miroir plan, passe par le centre optique $O_e$ de l'oculaire \textbf{sans déviation}.
  \item \textbf{Rayon 3 (incident vers le foyer objet $F_o$ du primaire)} : Pour un miroir concave, $F_o$ et $F'_o$ sont confondus en $F$. Ce rayon passe par $F$ avant réflexion, est réfléchi \textbf{parallèle à l'axe} par le primaire, dévié par le plan, et traverse l'oculaire en passant par son foyer image $F'_e$ (sort parallèle à l'axe).
\end{enumerate}
\end{methode}

\courssub{Méthode : Faire la mise au point}
\begin{methode}[Faire la mise au point (réglage normal)]
\begin{enumerate}
  \item Pointer le télescope vers un objet lointain (étoiles, antenne relais, sommet du Mont Cameroun, Lune) : les rayons incidents sont parallèles.
  \item Regarder dans l'oculaire : si l'image est floue, tourner la molette de mise au point.
  \item Cette molette \textbf{déplace l'oculaire le long de l'axe optique} (rapprochement ou éloignement du miroir plan) jusqu'à ce que l'image devienne nette. Cela correspond à la coïncidence du foyer objet de l'oculaire $F_e$ avec l'image intermédiaire (foyer $F'_o$ du primaire).
  \item Pour un œil myope ou presbyte : déplacer légèrement l'oculaire par rapport à la position de réglage normal pour former l'image finale à la distance de vision nette de l'observateur (entre $\SI{25}{cm}$ et l'infini).
\end{enumerate}
\end{methode}

\courssub{Intérêt du télescope de Newton par rapport aux lunettes astronomiques}
\begin{savaistu}[Intérêt du télescope de Newton par rapport aux lunettes astronomiques (lentilles)]
\begin{itemize}
  \item \textbf{Aucune aberration chromatique} : La réflexion sur un miroir suit la loi de Snell-Descartes ($i = r$) qui est \textbf{indépendante de la longueur d'onde} $\lambda$, contrairement à la réfraction dans le verre ($n(\lambda)$ variable). Pas de franges colorées sur les images.
  \item \textbf{Un seul miroir à polir avec précision} (le primaire) au lieu de 4 faces sphériques pour un objectif achromatique (doublet) + 2 faces pour l'oculaire $\rightarrow$ coût et temps de fabrication réduits pour les grands diamètres.
  \item \textbf{Support par l'arrière} : Le miroir primaire peut être soutenu sur toute sa surface arrière (cellule de support, systèmes actifs), évitant la déformation sous son propre poids qui limite les lentilles objectives à $\approx \SI{1}{m}$ de diamètre (ex: télescopes de \SI{8}{m} à \SI{10}{m} comme le VLT ou Keck).
  \item \textbf{Tube plus court} : Le trajet de la lumière est plié deux fois (longueur mécanique $\approx f'_o$ au lieu de $f'_o + f'_e$ pour une lunette), ce qui rend l'instrument plus compact, plus stable au vent et moins coûteux à abriter (dôme plus petit).
  \item \textbf{Applications au Cameroun} : Idéal pour l'observation des satellites de télécommunication (MTN, Orange, CAMTEL) en orbite géostationnaire, le suivi des débris spatiaux, ou l'astronomie amateur (observation de la Lune, planètes, nébuleuses) depuis les zones rurales peu polluées lumineusement.
\end{itemize}
\end{savaistu}

\courssub{Checklist avant le Probatoire}
\begin{aretenir}[Checklist avant le Probatoire]
\begin{itemize}
  \item $\square$ Je sais décrire les 3 composants principaux (miroir primaire concave parabolique, miroir plan secondaire incliné à 45°, oculaire convergente) et leur rôle respectif.
  \item $\square$ Je sais expliquer pourquoi l'image intermédiaire est réelle et inversée, et pourquoi l'image finale est virtuelle et inversée par rapport à l'objet.
  \item $\square$ Je sais tracer correctement les 3 rayons caractéristiques pour un faisceau incident parallèle à l'axe (réglage normal) : rayon parallèle $\to F'_o$, rayon au sommet $S_o$, rayon passant par $F_o$.
  \item $\square$ Je connais la formule du grossissement $G = -f'_o/f'_e$ et je sais l'appliquer avec les bonnes unités (m ou cm cohérents) ; j'interprète le signe $-$ (image inversée).
  \item $\square$ Je sais définir la puissance commerciale $P = 1/f'_e$ (en dioptries si $f'_e$ en mètres) et la distinguer du grossissement $G$.
  \item $\square$ Je sais expliquer la mise au point : déplacement de l'oculaire pour confondre $F_e$ et $F'_o$ (réglage normal) ou adapter à la vision de l'observateur.
  \item $\square$ Je sais citer au moins 3 avantages du Newton sur la lunette : pas d'aberration chromatique, grand diamètre possible (support arrière), tube court, coût réduit pour grande ouverture.
  \item $\square$ Je sais expliquer pourquoi il n'y a pas d'aberration chromatique (réflexion indépendante de $\lambda$, contrairement à la réfraction $n(\lambda)$).
  \item $\square$ Je sais calculer le pouvoir séparateur $\theta_{\min} = 1,22\lambda/D$ et en interpréter le sens (angle minimal résoluble, résolution angulaire).
  \item $\square$ Je sais identifier la pupille d'entrée (miroir primaire) et la pupille de sortie (image du primaire par l'oculaire) ; je sais que l'œil doit se placer à la pupille de sortie.
  \item $\square$ Je sais répondre à une question de type "Pourquoi le miroir primaire est-il parabolique et non sphérique ?" (élimination de l'aberration sphérique pour les grandes ouvertures / grands diamètres, mise au point unique pour les rayons parallèles).
  \item $\square$ Je sais schématiser le chemin optique complet (objet $\to$ primaire $\to$ plan $\to$ oculaire $\to$ œil) avec les sens des flèches lumineuses et la position des foyers.
\end{itemize}
\end{aretenir}

\findecours

\end{document}
